% ============================================================
%  Analyse Réelle II — Notes de Cours
%  Licence L1, Mathématiques
%  Partie 1 : Chapitres 7–8
%  « De l'intégrale aux fonctions de plusieurs variables »
% ============================================================
\documentclass[12pt,a4paper,openany]{book}

% --- Langues et encodage ---
\usepackage[utf8]{inputenc}
\usepackage[T1]{fontenc}
\usepackage[french]{babel}

% --- Géométrie ---
\usepackage[margin=2.5cm,top=3cm,bottom=3cm]{geometry}

% --- Mathématiques ---
\usepackage{amsmath,amssymb,amsthm,mathrsfs,mathtools}
\usepackage{caption}

% --- Graphiques ---
\usepackage{tikz}
\usetikzlibrary{calc,patterns,decorations.markings,arrows.meta,positioning,intersections}
\usepackage{pgfplots}
\usepgfplotslibrary{fillbetween}
\pgfplotsset{compat=1.18}

% --- Boîtes colorées ---
\usepackage[most]{tcolorbox}
\tcbuselibrary{theorems,skins,breakable}

% --- Liens hypertexte ---
\usepackage[colorlinks=true,linkcolor=blue!60!black,citecolor=green!50!black,urlcolor=blue!70!black]{hyperref}

% --- Index ---
\usepackage{makeidx}
\makeindex

% --- Listes ---
\usepackage{enumitem}

% --- En-têtes ---
\usepackage{fancyhdr}
\pagestyle{fancy}
\fancyhf{}
\fancyhead[LE]{\leftmark}
\fancyhead[RO]{\rightmark}
\fancyfoot[C]{\thepage}
\renewcommand{\headrulewidth}{0.4pt}

% --- Divers ---
\usepackage{xcolor}
\usepackage{microtype}

% ============================================================
%  Environnements de théorèmes
% ============================================================

% Environnements (amsthm + tcolorbox wrapping)
\theoremstyle{definition}
\newtheorem{definition}{Définition}[chapter]
\newtheorem{theoreme}{Théorème}[chapter]
\newtheorem{proposition}{Proposition}[chapter]
\newtheorem{lemme}{Lemme}[chapter]
\newtheorem{corollaire}{Corollaire}[chapter]
\newtheorem{remarque}{Remarque}[chapter]
\newtheorem{exemple}{Exemple}[chapter]
\newtheorem{exercice}{Exercice}[chapter]

\tcolorboxenvironment{definition}{enhanced,breakable,colback=blue!5!white,colframe=blue!60!black,left=8pt,right=8pt,top=6pt,bottom=6pt,before skip=12pt,after skip=12pt}
\tcolorboxenvironment{theoreme}{enhanced,breakable,colback=red!5,colframe=red!60!black,left=6pt,right=6pt,top=4pt,bottom=4pt,before skip=10pt,after skip=10pt}
\tcolorboxenvironment{proposition}{enhanced,breakable,colback=orange!5,colframe=orange!60!black,left=6pt,right=6pt,top=4pt,bottom=4pt,before skip=10pt,after skip=10pt}

% ============================================================
%  Commandes personnalisées
% ============================================================
\newcommand{\N}{\mathbb{N}}
\newcommand{\Z}{\mathbb{Z}}
\newcommand{\Q}{\mathbb{Q}}
\newcommand{\R}{\mathbb{R}}
\newcommand{\C}{\mathbb{C}}
\newcommand{\eps}{\varepsilon}
\newcommand{\abs}[1]{\left\lvert #1 \right\rvert}
\newcommand{\norm}[1]{\left\lVert #1 \right\rVert}
\newcommand{\dd}{\mathrm{d}}
\newcommand{\tomark}[1]{\texorpdfstring{#1}{}}

% Compteur de chapitre démarrant à 7
\setcounter{chapter}{6}

% ============================================================
\begin{document}

% ============================================================
%  Page de titre
% ============================================================
\begin{titlepage}
\begin{tikzpicture}[remember picture, overlay]
  \fill[left color=blue!15!white, right color=blue!5!white]
    (current page.south west) rectangle (current page.north east);
  \fill[color=orange!70!yellow]
    ([yshift=-2cm]current page.north west) rectangle ([yshift=-2.3cm]current page.north east);
  \fill[color=blue!80!black, rounded corners=8pt]
    ([xshift=3cm, yshift=-4cm]current page.north west) rectangle
    ([xshift=-3cm, yshift=-10cm]current page.north east);
  \node[anchor=center, text=white, font=\Huge\bfseries]
    at ([yshift=-6cm]current page.north) {Analyse R\'eelle II};
  \node[anchor=center, text=orange!80!yellow, font=\Large]
    at ([yshift=-7.5cm]current page.north) {Notes de Cours};
  \node[anchor=center, text=white!80, font=\large]
    at ([yshift=-8.8cm]current page.north) {Licence L3 --- 2025--2026};
  \node[anchor=center, text=blue!80!black, font=\Large\itshape]
    at ([yshift=-12cm]current page.north) {Ya\'e Ulrich Gaba};
  \draw[orange!70!yellow, line width=2pt]
    ([xshift=5cm, yshift=-13.5cm]current page.north west) --
    ([xshift=-5cm, yshift=-13.5cm]current page.north east);
  \node[anchor=center, text width=12cm, align=center, text=blue!60!black, font=\itshape]
    at ([yshift=-16cm]current page.north) {<<\,On peut toujours diviser les choses en deux : celles qui sont continues et celles qui ne le sont pas.\,>>\\[4pt]--- Henri Poincar\'e};
  \node[anchor=south, text=gray, font=\small]
    at ([yshift=1.5cm]current page.south) {\today};
\end{tikzpicture}
\end{titlepage}

% ============================================================
%  Table des matières
% ============================================================
\tableofcontents

% ============================================================
%  Préface
% ============================================================
\chapter*{Préface}
\addcontentsline{toc}{chapter}{Préface}

Ce cours constitue la suite naturelle d'\emph{Analyse Réelle~I}, dans lequel nous avons étudié les fondements de l'analyse : les nombres réels, les suites numériques, les fonctions continues et le calcul différentiel des fonctions d'une variable réelle. Le présent volume --- \emph{Analyse Réelle~II} --- prolonge et approfondit cette étude en abordant trois thèmes majeurs.

\medskip

\noindent\textbf{L'intégration.} Nous construisons rigoureusement l'intégrale de Riemann, en démontrons les propriétés fondamentales, et établissons le théorème fondamental de l'analyse qui relie intégration et dérivation. Les techniques classiques d'intégration (intégration par parties, changement de variable) sont présentées avec toute la rigueur requise.

\medskip

\noindent\textbf{Les suites et séries de fonctions.} L'étude de la convergence uniforme est un tournant conceptuel essentiel : elle permet de comprendre quand on peut intervertir limites et opérations d'analyse (continuité, intégration, dérivation). Les séries entières fournissent un cadre privilégié d'application de ces résultats.

\medskip

\noindent\textbf{Les fonctions de plusieurs variables.} Nous étendons les notions de continuité, de différentiabilité et d'intégration au cadre des fonctions de $\R^n$ dans $\R^p$, culminant avec le théorème d'inversion locale et le théorème des fonctions implicites.

\medskip

Tout au long de ce texte, nous adoptons les conventions de Bourbaki : les démonstrations sont complètes, les énoncés sont précis, et les contre-exemples sont mis en avant pour délimiter la portée des théorèmes. Le lecteur est encouragé à traiter les exercices proposés à la fin de chaque chapitre, classés par ordre de difficulté croissante.

\bigskip

\hfill\emph{Bonne lecture.}

% ############################################################
%  CHAPITRE 7 : INTÉGRALE DE RIEMANN
% ############################################################
\chapter{Intégrale de Riemann --- Définition, Propriétés, Théorème Fondamental}
\label{chap:integrale-riemann}
\index{intégrale de Riemann}

\section*{Introduction}
\addcontentsline{toc}{section}{Introduction}

Le problème du calcul de l'aire d'une figure plane est l'un des plus anciens de l'histoire des mathématiques. Dès l'Antiquité, Archimède sut calculer l'aire d'un segment parabolique par la \emph{méthode d'exhaustion}\index{méthode d'exhaustion} : il encadrait la figure par des polygones inscrits et circonscrits dont les aires convergeaient vers une valeur commune. C'est précisément cette idée --- approcher une quantité par des sommes finies de plus en plus fines --- qui est au c\oe ur de la construction de l'intégrale de Riemann.

Si $f \colon [a,b] \to \R$ est une fonction positive, l'aire du domaine compris entre le graphe de~$f$ et l'axe des abscisses est, intuitivement, la « somme » des aires de rectangles infiniment fins de largeur~$\dd x$ et de hauteur~$f(x)$. La formalisation rigoureuse de cette idée passe par la notion de \emph{subdivision} d'un intervalle et par les \emph{sommes de Darboux}\index{sommes de Darboux}, que nous introduisons dans la section suivante.

% ============================================================
\section{Subdivisions et sommes de Darboux}
\label{sec:subdivisions-darboux}
\index{subdivision}

Soit $[a,b]$ un intervalle compact de~$\R$, avec $a < b$.

\begin{definition}[Subdivision]\label{def:subdivision}
  On appelle \emph{subdivision}\index{subdivision} (ou \emph{partition}) de l'intervalle $[a,b]$ toute famille finie de points
  \[
    \sigma = (x_0, x_1, \ldots, x_n)
  \]
  telle que $a = x_0 < x_1 < \cdots < x_n = b$. L'entier $n$ est le \emph{nombre de sous-intervalles} de la subdivision, et les intervalles $[x_{i-1}, x_i]$, pour $1 \leqslant i \leqslant n$, sont les \emph{sous-intervalles} de~$\sigma$. On note $\abs{\sigma} = \max_{1 \leqslant i \leqslant n}(x_i - x_{i-1})$ le \emph{pas} de la subdivision.
\end{definition}

\begin{definition}[Sommes de Darboux]\label{def:sommes-darboux}
  Soit $f \colon [a,b] \to \R$ une fonction bornée et $\sigma = (x_0, \ldots, x_n)$ une subdivision de~$[a,b]$. Pour chaque $i \in \{1, \ldots, n\}$, on pose
  \[
    m_i = \inf_{x \in [x_{i-1}, x_i]} f(x), \qquad
    M_i = \sup_{x \in [x_{i-1}, x_i]} f(x).
  \]
  La \emph{somme de Darboux inférieure}\index{somme de Darboux!inférieure} de $f$ relativement à~$\sigma$ est
  \[
    s(f, \sigma) = \sum_{i=1}^{n} m_i \, (x_i - x_{i-1}),
  \]
  et la \emph{somme de Darboux supérieure}\index{somme de Darboux!supérieure} est
  \[
    S(f, \sigma) = \sum_{i=1}^{n} M_i \, (x_i - x_{i-1}).
  \]
\end{definition}

\begin{center}
\begin{tikzpicture}[scale=0.9]
  \begin{axis}[
    axis lines=middle,
    xlabel={$x$},
    ylabel={$y$},
    xmin=-0.3, xmax=5.5,
    ymin=-0.3, ymax=3.5,
    xtick={0.5,1.5,2.5,3.5,4.5},
    xticklabels={$x_0$,$x_1$,$x_2$,$x_3$,$x_4$},
    ytick=\empty,
    width=12cm, height=7cm,
    title={Sommes de Darboux inférieure (bleu) et supérieure (rouge)},
    every axis title/.style={at={(0.5,1.05)}},
  ]
    % Fonction
    \addplot[domain=0.5:4.5, samples=100, thick, black] {0.3*(x-1)^2 + 0.5*sin(deg(x)) + 1};

    % Somme inférieure (bleu)
    \addplot[ybar interval, fill=blue!20, draw=blue!60, opacity=0.5]
      coordinates {(0.5,1.01) (1.5,0.85) (2.5,1.18) (3.5,2.08) (4.5,2.08)};

    % Somme supérieure (rouge)
    \addplot[ybar interval, fill=red!15, draw=red!60, opacity=0.4]
      coordinates {(0.5,1.37) (1.5,1.37) (2.5,2.08) (3.5,3.15) (4.5,3.15)};

    % Fonction (redessiner par-dessus)
    \addplot[domain=0.5:4.5, samples=100, very thick, black] {0.3*(x-1)^2 + 0.5*sin(deg(x)) + 1};
  \end{axis}
\end{tikzpicture}
\end{center}

\begin{definition}[Raffinement]\label{def:raffinement}
  Soient $\sigma$ et $\sigma'$ deux subdivisions de $[a,b]$. On dit que $\sigma'$ est un \emph{raffinement}\index{raffinement} de~$\sigma$ si $\sigma \subset \sigma'$ (au sens ensembliste). La subdivision $\sigma \cup \sigma'$ est le raffinement commun le plus grossier de $\sigma$ et~$\sigma'$.
\end{definition}

\begin{proposition}\label{prop:raffinement-darboux}
  Soit $f \colon [a,b] \to \R$ bornée. Si $\sigma'$ est un raffinement de~$\sigma$, alors
  \[
    s(f, \sigma) \leqslant s(f, \sigma') \leqslant S(f, \sigma') \leqslant S(f, \sigma).
  \]
\end{proposition}

\begin{proof}
  Il suffit de traiter le cas où $\sigma'$ est obtenue en ajoutant un seul point~$c$ à~$\sigma$. Le cas général s'en déduit par récurrence.

  Soit $\sigma = (x_0, \ldots, x_n)$ et supposons $c \in \,]x_{k-1}, x_k[$ pour un certain $k \in \{1, \ldots, n\}$. Posons
  \[
    m_k = \inf_{[x_{k-1}, x_k]} f, \quad
    m_k' = \inf_{[x_{k-1}, c]} f, \quad
    m_k'' = \inf_{[c, x_k]} f.
  \]
  Puisque $[x_{k-1}, c] \subset [x_{k-1}, x_k]$ et $[c, x_k] \subset [x_{k-1}, x_k]$, on a $m_k' \geqslant m_k$ et $m_k'' \geqslant m_k$. Par conséquent,
  \begin{align*}
    s(f, \sigma') - s(f, \sigma)
    &= m_k'(c - x_{k-1}) + m_k''(x_k - c) - m_k(x_k - x_{k-1}) \\
    &= (m_k' - m_k)(c - x_{k-1}) + (m_k'' - m_k)(x_k - c) \\
    &\geqslant 0.
  \end{align*}
  Ainsi $s(f, \sigma) \leqslant s(f, \sigma')$. Le raisonnement pour les sommes supérieures est analogue (les inégalités sont renversées). L'inégalité $s(f, \sigma') \leqslant S(f, \sigma')$ est immédiate car $m_i \leqslant M_i$ pour tout~$i$.
\end{proof}

\begin{corollaire}\label{cor:inf-leq-sup}
  Pour toutes subdivisions $\sigma$ et $\sigma'$ de $[a,b]$ (pas nécessairement l'une raffinement de l'autre),
  \[
    s(f, \sigma) \leqslant S(f, \sigma').
  \]
\end{corollaire}

\begin{proof}
  La subdivision $\sigma \cup \sigma'$ est un raffinement commun. Par la proposition~\ref{prop:raffinement-darboux},
  \[
    s(f, \sigma) \leqslant s(f, \sigma \cup \sigma') \leqslant S(f, \sigma \cup \sigma') \leqslant S(f, \sigma'). \qedhere
  \]
\end{proof}

% ============================================================
\section{Intégrale inférieure et supérieure, fonctions intégrables}
\label{sec:integrale-inf-sup}
\index{intégrale!inférieure}\index{intégrale!supérieure}

Le corollaire~\ref{cor:inf-leq-sup} montre que l'ensemble des sommes de Darboux inférieures est majoré (par toute somme supérieure) et l'ensemble des sommes supérieures est minoré (par toute somme inférieure). La borne supérieure de l'un et la borne inférieure de l'autre existent donc dans~$\R$.

\begin{definition}[Intégrales inférieure et supérieure]\label{def:int-inf-sup}
  Soit $f \colon [a,b] \to \R$ bornée. On définit l'\emph{intégrale inférieure} et l'\emph{intégrale supérieure} de $f$ sur $[a,b]$ par
  \[
    \underline{\int_a^b} f = \sup_{\sigma} s(f, \sigma), \qquad
    \overline{\int_a^b} f = \inf_{\sigma} S(f, \sigma),
  \]
  où le supremum et l'infimum sont pris sur toutes les subdivisions $\sigma$ de~$[a,b]$.
\end{definition}

Par le corollaire~\ref{cor:inf-leq-sup}, on a toujours
\[
  \underline{\int_a^b} f \leqslant \overline{\int_a^b} f.
\]

\begin{definition}[Fonction intégrable au sens de Riemann]\label{def:integrabilite-riemann}
  Une fonction bornée $f \colon [a,b] \to \R$ est dite \emph{intégrable au sens de Riemann}\index{intégrable au sens de Riemann} sur $[a,b]$ si
  \[
    \underline{\int_a^b} f = \overline{\int_a^b} f.
  \]
  La valeur commune est alors notée $\displaystyle\int_a^b f(x)\,\dd x$, ou plus brièvement $\displaystyle\int_a^b f$, et s'appelle l'\emph{intégrale de Riemann}\index{intégrale de Riemann} de $f$ sur~$[a,b]$. L'ensemble des fonctions intégrables au sens de Riemann sur $[a,b]$ est noté $\mathcal{R}([a,b])$.
\end{definition}

\begin{theoreme}[Critère d'intégrabilité de Riemann]\label{thm:critere-integrabilite}
  \index{critère d'intégrabilité}
  Soit $f \colon [a,b] \to \R$ bornée. Alors $f$ est intégrable au sens de Riemann si et seulement si
  \[
    \forall\, \eps > 0,\; \exists\, \sigma \text{ subdivision de } [a,b], \quad S(f,\sigma) - s(f,\sigma) < \eps.
  \]
\end{theoreme}

\begin{proof}
  \emph{Condition suffisante.} Supposons que pour tout $\eps > 0$, il existe une subdivision $\sigma$ telle que $S(f,\sigma) - s(f,\sigma) < \eps$. Alors
  \[
    0 \leqslant \overline{\int_a^b} f - \underline{\int_a^b} f
    \leqslant S(f,\sigma) - s(f,\sigma) < \eps.
  \]
  Comme ceci vaut pour tout $\eps > 0$, on conclut que $\displaystyle\underline{\int_a^b} f = \overline{\int_a^b} f$, donc $f$ est intégrable.

  \emph{Condition nécessaire.} Supposons $f$ intégrable et soit $\eps > 0$. Par définition du supremum et de l'infimum, il existe des subdivisions $\sigma_1$ et $\sigma_2$ telles que
  \[
    \underline{\int_a^b} f - \frac{\eps}{2} < s(f, \sigma_1) \leqslant \underline{\int_a^b} f
    \quad \text{et} \quad
    \overline{\int_a^b} f \leqslant S(f, \sigma_2) < \overline{\int_a^b} f + \frac{\eps}{2}.
  \]
  Posons $\sigma = \sigma_1 \cup \sigma_2$. Comme $\sigma$ raffine à la fois $\sigma_1$ et $\sigma_2$, on a
  \[
    S(f, \sigma) - s(f, \sigma) \leqslant S(f, \sigma_2) - s(f, \sigma_1)
    < \overline{\int_a^b} f + \frac{\eps}{2} - \underline{\int_a^b} f + \frac{\eps}{2}
    = \eps,
  \]
  puisque l'intégrabilité donne $\displaystyle\overline{\int_a^b} f = \underline{\int_a^b} f$.
\end{proof}

\begin{center}
\begin{tikzpicture}[scale=0.85]
  \begin{axis}[
    axis lines=middle,
    xlabel={$x$},
    ylabel={},
    xmin=-0.2, xmax=6.5,
    ymin=-0.3, ymax=3.8,
    xtick=\empty,
    ytick=\empty,
    width=13cm, height=6.5cm,
    title={Convergence des sommes de Darboux lors du raffinement},
    every axis title/.style={at={(0.5,1.05)}},
  ]
    % Fonction
    \addplot[domain=0.5:6, samples=200, very thick, black] {0.15*(x-2)^2 + 0.8*sin(deg(0.8*x)) + 1.5};

    % Somme inférieure — subdivision plus fine
    \addplot[ybar interval, fill=blue!25, draw=blue!50, opacity=0.6]
      coordinates {
        (0.5,1.3) (1.25,1.1) (2.0,0.9) (2.75,1.0) (3.5,1.5)
        (4.25,2.0) (5.0,2.2) (5.5,2.5) (6.0,2.5)
      };

    % Somme supérieure — subdivision plus fine
    \addplot[ybar interval, fill=red!15, draw=red!50, opacity=0.4]
      coordinates {
        (0.5,1.8) (1.25,1.3) (2.0,1.1) (2.75,1.5) (3.5,2.0)
        (4.25,2.5) (5.0,2.7) (5.5,3.0) (6.0,3.0)
      };

    \addplot[domain=0.5:6, samples=200, very thick, black] {0.15*(x-2)^2 + 0.8*sin(deg(0.8*x)) + 1.5};
  \end{axis}
\end{tikzpicture}
\end{center}

% ============================================================
\section{Fonctions intégrables}
\label{sec:fonctions-integrables}
\index{fonctions intégrables}

\begin{theoreme}\label{thm:continue-integrable}
  Toute fonction continue sur $[a,b]$ est intégrable au sens de Riemann.
  \index{fonction continue!intégrabilité}
\end{theoreme}

\begin{proof}
  Soit $f \colon [a,b] \to \R$ continue. Puisque $[a,b]$ est compact, $f$ est uniformément continue sur $[a,b]$ (théorème de Heine). Soit $\eps > 0$. Il existe $\delta > 0$ tel que
  \[
    \forall\, x, y \in [a,b], \quad \abs{x - y} < \delta \implies \abs{f(x) - f(y)} < \frac{\eps}{b - a}.
  \]
  Choisissons une subdivision $\sigma = (x_0, \ldots, x_n)$ de pas $\abs{\sigma} < \delta$. Sur chaque sous-intervalle $[x_{i-1}, x_i]$, la fonction $f$ est continue et atteint ses bornes (théorème des valeurs extrêmes). Notons $M_i = \max_{[x_{i-1},x_i]} f$ et $m_i = \min_{[x_{i-1},x_i]} f$. Puisque $x_i - x_{i-1} < \delta$, les points où le maximum et le minimum sont atteints sont distants de moins de~$\delta$, d'où
  \[
    M_i - m_i < \frac{\eps}{b - a}.
  \]
  Par conséquent,
  \[
    S(f, \sigma) - s(f, \sigma) = \sum_{i=1}^{n}(M_i - m_i)(x_i - x_{i-1})
    < \frac{\eps}{b-a} \sum_{i=1}^{n}(x_i - x_{i-1}) = \frac{\eps}{b-a} \cdot (b - a) = \eps.
  \]
  Le critère d'intégrabilité (théorème~\ref{thm:critere-integrabilite}) permet de conclure.
\end{proof}

\begin{theoreme}\label{thm:monotone-integrable}
  Toute fonction monotone sur $[a,b]$ est intégrable au sens de Riemann.
  \index{fonction monotone!intégrabilité}
\end{theoreme}

\begin{proof}
  Traitons le cas où $f$ est croissante ; le cas décroissant est analogue. La fonction $f$ est bornée sur $[a,b]$ (par $f(a)$ et $f(b)$). Soit $\eps > 0$. Considérons la subdivision régulière $\sigma_n = (x_0, \ldots, x_n)$ avec $x_i = a + i\,\frac{b-a}{n}$ pour $0 \leqslant i \leqslant n$.

  Puisque $f$ est croissante, sur chaque sous-intervalle $[x_{i-1}, x_i]$, on a $m_i = f(x_{i-1})$ et $M_i = f(x_i)$. Ainsi,
  \begin{align*}
    S(f, \sigma_n) - s(f, \sigma_n)
    &= \sum_{i=1}^{n} \bigl(f(x_i) - f(x_{i-1})\bigr) \cdot \frac{b-a}{n} \\
    &= \frac{b-a}{n} \sum_{i=1}^{n} \bigl(f(x_i) - f(x_{i-1})\bigr) \\
    &= \frac{b-a}{n} \bigl(f(b) - f(a)\bigr),
  \end{align*}
  par télescopage. En choisissant $n$ assez grand pour que $\frac{(b-a)(f(b)-f(a))}{n} < \eps$, on obtient $S(f,\sigma_n) - s(f,\sigma_n) < \eps$, et le critère d'intégrabilité s'applique.
\end{proof}

\begin{definition}[Fonction continue par morceaux]\label{def:continue-par-morceaux}
  \index{fonction continue par morceaux}
  Une fonction $f \colon [a,b] \to \R$ est dite \emph{continue par morceaux} s'il existe une subdivision $a = c_0 < c_1 < \cdots < c_p = b$ telle que, pour tout $k \in \{1, \ldots, p\}$, la restriction de $f$ à $]c_{k-1}, c_k[$ se prolonge en une fonction continue sur $[c_{k-1}, c_k]$.
\end{definition}

\begin{theoreme}\label{thm:cpm-integrable}
  Toute fonction continue par morceaux sur $[a,b]$ est intégrable au sens de Riemann.
\end{theoreme}

\begin{proof}
  Soit $f$ continue par morceaux, de subdivision associée $(c_0, \ldots, c_p)$. Soit $\eps > 0$. Sur chaque intervalle $[c_{k-1}, c_k]$, le prolongement continu $\tilde{f}_k$ est intégrable d'après le théorème~\ref{thm:continue-integrable}. Il existe donc une subdivision $\sigma_k$ de $[c_{k-1}, c_k]$ telle que $S(\tilde{f}_k, \sigma_k) - s(\tilde{f}_k, \sigma_k) < \eps / p$. Comme $f$ et $\tilde{f}_k$ ne diffèrent qu'en au plus deux points (les extrémités), et que la modification d'une fonction bornée en un nombre fini de points ne change pas ses sommes de Darboux, la subdivision $\sigma = \sigma_1 \cup \cdots \cup \sigma_p$ satisfait $S(f,\sigma) - s(f,\sigma) < \eps$.
\end{proof}

% ============================================================
\section{Propriétés de l'intégrale}
\label{sec:proprietes-integrale}
\index{intégrale!propriétés}

Dans cette section, $f$ et $g$ désignent des fonctions intégrables au sens de Riemann sur $[a,b]$.

\begin{theoreme}[Linéarité]\label{thm:linearite}
  \index{linéarité de l'intégrale}
  Si $f, g \in \mathcal{R}([a,b])$ et $\alpha, \beta \in \R$, alors $\alpha f + \beta g \in \mathcal{R}([a,b])$ et
  \[
    \int_a^b (\alpha f + \beta g) = \alpha \int_a^b f + \beta \int_a^b g.
  \]
\end{theoreme}

\begin{proof}
  Traitons d'abord l'homogénéité. Si $\alpha \geqslant 0$, pour toute subdivision $\sigma$, $s(\alpha f, \sigma) = \alpha\, s(f, \sigma)$ et $S(\alpha f, \sigma) = \alpha\, S(f, \sigma)$, d'où $S(\alpha f,\sigma) - s(\alpha f, \sigma) = \alpha(S(f,\sigma) - s(f,\sigma))$, et l'intégrabilité de $\alpha f$ suit de celle de~$f$. Si $\alpha < 0$, on a $s(\alpha f, \sigma) = \alpha\, S(f, \sigma)$ et $S(\alpha f, \sigma) = \alpha\, s(f, \sigma)$, et le même raisonnement s'applique.

  Pour l'additivité, on utilise l'inégalité
  \[
    \inf_{[x_{i-1},x_i]} f + \inf_{[x_{i-1},x_i]} g \leqslant \inf_{[x_{i-1},x_i]} (f+g)
  \]
  et l'inégalité analogue pour le supremum :
  \[
    \sup_{[x_{i-1},x_i]}(f+g) \leqslant \sup_{[x_{i-1},x_i]} f + \sup_{[x_{i-1},x_i]} g.
  \]
  Par conséquent, pour toute subdivision $\sigma$,
  \[
    s(f,\sigma) + s(g,\sigma) \leqslant s(f+g,\sigma) \leqslant S(f+g,\sigma) \leqslant S(f,\sigma) + S(g,\sigma).
  \]
  Soit $\eps > 0$. Choisissons $\sigma_1$ et $\sigma_2$ telles que $S(f,\sigma_1) - s(f,\sigma_1) < \eps/2$ et $S(g,\sigma_2) - s(g,\sigma_2) < \eps/2$. Posons $\sigma = \sigma_1 \cup \sigma_2$. Alors
  \[
    S(f+g,\sigma) - s(f+g,\sigma) \leqslant \bigl(S(f,\sigma) - s(f,\sigma)\bigr) + \bigl(S(g,\sigma) - s(g,\sigma)\bigr) < \eps.
  \]
  Donc $f + g$ est intégrable. Pour la valeur de l'intégrale, les encadrements ci-dessus donnent, en passant au supremum et à l'infimum :
  \[
    \int_a^b f + \int_a^b g \leqslant \int_a^b(f+g) \leqslant \int_a^b f + \int_a^b g. \qedhere
  \]
\end{proof}

\begin{theoreme}[Positivité]\label{thm:positivite}
  \index{positivité de l'intégrale}
  Si $f \in \mathcal{R}([a,b])$ et $f \geqslant 0$ sur $[a,b]$, alors
  \[
    \int_a^b f \geqslant 0.
  \]
\end{theoreme}

\begin{proof}
  Si $f \geqslant 0$, alors pour toute subdivision~$\sigma$, les bornes inférieures $m_i \geqslant 0$, d'où $s(f,\sigma) \geqslant 0$. Par passage au supremum, $\displaystyle\int_a^b f = \sup_\sigma s(f,\sigma) \geqslant 0$.
\end{proof}

\begin{corollaire}[Croissance]\label{cor:croissance}
  \index{croissance de l'intégrale}
  Si $f, g \in \mathcal{R}([a,b])$ et $f \leqslant g$ sur $[a,b]$, alors
  \[
    \int_a^b f \leqslant \int_a^b g.
  \]
\end{corollaire}

\begin{proof}
  La fonction $g - f \geqslant 0$ est intégrable par linéarité, et $\displaystyle\int_a^b (g - f) \geqslant 0$ par positivité. La linéarité conclut.
\end{proof}

\begin{theoreme}[Relation de Chasles]\label{thm:chasles}
  \index{relation de Chasles}
  Soit $f \in \mathcal{R}([a,b])$ et $c \in \,]a,b[\,$. Alors $f$ est intégrable sur $[a,c]$ et sur $[c,b]$, et
  \[
    \int_a^b f = \int_a^c f + \int_c^b f.
  \]
\end{theoreme}

\begin{proof}
  Soit $\eps > 0$. Il existe une subdivision $\sigma$ de $[a,b]$ telle que $S(f,\sigma) - s(f,\sigma) < \eps$. Posons $\sigma' = \sigma \cup \{c\}$, qui est un raffinement de~$\sigma$, d'où $S(f,\sigma') - s(f,\sigma') \leqslant S(f,\sigma) - s(f,\sigma) < \eps$.

  La subdivision $\sigma'$ se décompose naturellement en $\sigma'_1$ (points dans $[a,c]$) et $\sigma'_2$ (points dans $[c,b]$), et l'on a
  \[
    S(f,\sigma') - s(f,\sigma') = \bigl(S(f|_{[a,c]}, \sigma'_1) - s(f|_{[a,c]}, \sigma'_1)\bigr) + \bigl(S(f|_{[c,b]}, \sigma'_2) - s(f|_{[c,b]}, \sigma'_2)\bigr).
  \]
  Chaque terme étant positif et leur somme étant inférieure à $\eps$, chaque terme est inférieur à~$\eps$. Donc $f$ est intégrable sur $[a,c]$ et sur $[c,b]$. L'égalité des intégrales suit de $s(f,\sigma') = s(f|_{[a,c]},\sigma'_1) + s(f|_{[c,b]},\sigma'_2)$ et du passage au supremum.
\end{proof}

\begin{theoreme}[Inégalité de la valeur moyenne]\label{thm:ineg-valeur-moyenne}
  \index{inégalité de la valeur moyenne}
  Si $f \in \mathcal{R}([a,b])$, alors $\abs{f} \in \mathcal{R}([a,b])$ et
  \[
    \abs*{\int_a^b f} \leqslant \int_a^b \abs{f}.
  \]
\end{theoreme}

\begin{proof}
  Montrons d'abord que $\abs{f}$ est intégrable. Pour tout sous-intervalle $[x_{i-1}, x_i]$,
  \[
    \sup_{[x_{i-1},x_i]} \abs{f} - \inf_{[x_{i-1},x_i]} \abs{f}
    \leqslant \sup_{[x_{i-1},x_i]} f - \inf_{[x_{i-1},x_i]} f,
  \]
  car $\bigl\lvert\abs{f(x)} - \abs{f(y)}\bigr\rvert \leqslant \abs{f(x) - f(y)}$ pour tous $x, y$. Donc $S(\abs{f},\sigma) - s(\abs{f},\sigma) \leqslant S(f,\sigma) - s(f,\sigma)$, et l'intégrabilité de $\abs{f}$ découle de celle de~$f$.

  Pour l'inégalité, on utilise $-\abs{f} \leqslant f \leqslant \abs{f}$. Par croissance de l'intégrale,
  \[
    -\int_a^b \abs{f} \leqslant \int_a^b f \leqslant \int_a^b \abs{f},
  \]
  ce qui revient à $\displaystyle\abs*{\int_a^b f} \leqslant \int_a^b \abs{f}$.
\end{proof}

\begin{theoreme}[Inégalité de Cauchy--Schwarz intégrale]\label{thm:cauchy-schwarz-integrale}
  \index{inégalité de Cauchy--Schwarz}
  Si $f, g \in \mathcal{R}([a,b])$, alors
  \[
    \left(\int_a^b f\, g\right)^2 \leqslant \left(\int_a^b f^2\right)\left(\int_a^b g^2\right).
  \]
\end{theoreme}

% ============================================================
\section{Théorèmes fondamentaux}
\label{sec:theoremes-fondamentaux}
\index{théorème fondamental de l'analyse}

\begin{theoreme}[Premier théorème fondamental de l'analyse]\label{thm:premier-fondamental}
  \index{théorème fondamental!premier}
  Soit $f \colon [a,b] \to \R$ continue. La fonction
  \[
    F \colon [a,b] \to \R, \quad F(x) = \int_a^x f(t)\,\dd t,
  \]
  est dérivable sur $[a,b]$ et $F'(x) = f(x)$ pour tout $x \in [a,b]$.
\end{theoreme}

\begin{proof}
  Soit $x \in [a,b]$ et $h \neq 0$ tel que $x + h \in [a,b]$. Par la relation de Chasles,
  \[
    F(x+h) - F(x) = \int_a^{x+h} f(t)\,\dd t - \int_a^x f(t)\,\dd t = \int_x^{x+h} f(t)\,\dd t.
  \]
  On peut écrire
  \[
    \frac{F(x+h) - F(x)}{h} - f(x) = \frac{1}{h}\int_x^{x+h} \bigl(f(t) - f(x)\bigr)\,\dd t.
  \]
  Soit $\eps > 0$. Puisque $f$ est continue en~$x$, il existe $\delta > 0$ tel que $\abs{t - x} < \delta$ implique $\abs{f(t) - f(x)} < \eps$. Si $0 < \abs{h} < \delta$, alors pour tout $t$ entre $x$ et $x+h$, on a $\abs{t - x} \leqslant \abs{h} < \delta$, d'où $\abs{f(t) - f(x)} < \eps$. Par conséquent (en supposant $h > 0$, le cas $h < 0$ étant analogue),
  \[
    \abs*{\frac{F(x+h)-F(x)}{h} - f(x)}
    = \abs*{\frac{1}{h}\int_x^{x+h}(f(t)-f(x))\,\dd t}
    \leqslant \frac{1}{h}\int_x^{x+h}\abs{f(t)-f(x)}\,\dd t
    < \frac{1}{h} \cdot \eps \cdot h = \eps.
  \]
  Ceci montre que $\displaystyle\lim_{h \to 0}\frac{F(x+h)-F(x)}{h} = f(x)$, c'est-à-dire $F'(x) = f(x)$.
\end{proof}

\begin{center}
\begin{tikzpicture}[scale=0.9]
  \begin{axis}[
    axis lines=middle,
    xlabel={$x$},
    ylabel={},
    xmin=-0.5, xmax=5.5,
    ymin=-0.5, ymax=4.5,
    xtick={0.5, 3.5},
    xticklabels={$a$, $x$},
    ytick=\empty,
    width=11cm, height=7cm,
    title={$F(x) = \int_a^x f(t)\,\dd t$ : aire sous la courbe},
    every axis title/.style={at={(0.5,1.05)}},
  ]
    \addplot[domain=0.5:5, samples=100, very thick, black, name path=f] {0.2*(x-1)^2 + 0.5*sin(deg(x)) + 1.2};
    \addplot[domain=0.5:3.5, samples=100, draw=none, name path=zero] {0};
    \addplot[fill=blue!20, opacity=0.6] fill between[of=f and zero, soft clip={domain=0.5:3.5}];
    \node at (axis cs:2,0.7) {$F(x)$};
    \node at (axis cs:4.3,3.2) {$y = f(t)$};
  \end{axis}
\end{tikzpicture}
\end{center}

\begin{corollaire}\label{cor:primitive}
  Toute fonction continue sur un intervalle admet une primitive.
\end{corollaire}

\begin{theoreme}[Second théorème fondamental --- formule de Newton--Leibniz]\label{thm:newton-leibniz}
  \index{formule de Newton--Leibniz}\index{théorème fondamental!second}
  Soit $f \colon [a,b] \to \R$ continue et $G$ une primitive de $f$ sur $[a,b]$ (c'est-à-dire $G$ est dérivable et $G' = f$). Alors
  \[
    \int_a^b f(t)\,\dd t = G(b) - G(a).
  \]
  On note $G(b) - G(a) = \bigl[G(t)\bigr]_a^b$.
\end{theoreme}

\begin{proof}
  Soit $F(x) = \int_a^x f(t)\,\dd t$. Par le premier théorème fondamental, $F' = f = G'$ sur $[a,b]$. Donc $(F - G)' = 0$ sur $[a,b]$, ce qui implique que $F - G$ est constante (conséquence du théorème des accroissements finis). Notons $c = F(a) - G(a) = 0 - G(a) = -G(a)$. Alors $F(x) = G(x) - G(a)$ pour tout $x \in [a,b]$, et en particulier
  \[
    \int_a^b f(t)\,\dd t = F(b) = G(b) - G(a). \qedhere
  \]
\end{proof}

% ============================================================
\section{Techniques d'intégration}
\label{sec:techniques-integration}
\index{techniques d'intégration}

\begin{theoreme}[Intégration par parties]\label{thm:ipp}
  \index{intégration par parties}
  Soient $u, v \colon [a,b] \to \R$ de classe $\mathcal{C}^1$. Alors
  \[
    \int_a^b u(t)\, v'(t)\,\dd t = \bigl[u(t)\,v(t)\bigr]_a^b - \int_a^b u'(t)\, v(t)\,\dd t.
  \]
\end{theoreme}

\begin{proof}
  La fonction $uv$ est de classe $\mathcal{C}^1$ et $(uv)' = u'v + uv'$. En intégrant les deux membres entre $a$ et $b$ et en appliquant le second théorème fondamental,
  \[
    u(b)v(b) - u(a)v(a) = \int_a^b u'(t)\,v(t)\,\dd t + \int_a^b u(t)\,v'(t)\,\dd t.
  \]
  Le résultat s'obtient par réarrangement.
\end{proof}

\begin{exemple}\label{ex:ipp1}
  Calculons $\displaystyle I = \int_0^1 t\, e^t\,\dd t$.
  On pose $u(t) = t$ et $v'(t) = e^t$, d'où $u'(t) = 1$ et $v(t) = e^t$. Ainsi
  \[
    I = \bigl[t\, e^t\bigr]_0^1 - \int_0^1 e^t\,\dd t = e - \bigl[e^t\bigr]_0^1 = e - (e - 1) = 1.
  \]
\end{exemple}

\begin{exemple}\label{ex:ipp2}
  Calculons $\displaystyle J = \int_0^{\pi} t\,\sin t\,\dd t$.
  On pose $u(t) = t$ et $v'(t) = \sin t$, d'où $u'(t) = 1$ et $v(t) = -\cos t$. Ainsi
  \[
    J = \bigl[-t\cos t\bigr]_0^{\pi} + \int_0^{\pi}\cos t\,\dd t = \pi + \bigl[\sin t\bigr]_0^{\pi} = \pi + 0 = \pi.
  \]
\end{exemple}

\begin{exemple}\label{ex:ipp3}
  Calculons $\displaystyle K = \int_1^e (\ln t)^2\,\dd t$.
  On pose $u(t) = (\ln t)^2$ et $v'(t) = 1$, d'où $u'(t) = \frac{2\ln t}{t}$ et $v(t) = t$. Ainsi
  \[
    K = \bigl[t(\ln t)^2\bigr]_1^e - 2\int_1^e \ln t\,\dd t.
  \]
  Or $\displaystyle\int_1^e \ln t\,\dd t = \bigl[t\ln t - t\bigr]_1^e = (e - e) - (0 - 1) = 1$, d'où $K = e - 2$.
\end{exemple}

\begin{theoreme}[Changement de variable]\label{thm:changement-variable}
  \index{changement de variable}
  Soit $f \colon [a,b] \to \R$ continue et $\varphi \colon [\alpha, \beta] \to [a,b]$ de classe $\mathcal{C}^1$ avec $\varphi(\alpha) = a$ et $\varphi(\beta) = b$. Alors
  \[
    \int_a^b f(x)\,\dd x = \int_\alpha^\beta f\bigl(\varphi(t)\bigr)\,\varphi'(t)\,\dd t.
  \]
\end{theoreme}

\begin{proof}
  Soit $F$ une primitive de $f$ (qui existe d'après le corollaire~\ref{cor:primitive}). La fonction composée $F \circ \varphi$ est de classe $\mathcal{C}^1$ sur $[\alpha, \beta]$ et
  \[
    (F \circ \varphi)'(t) = F'(\varphi(t))\,\varphi'(t) = f(\varphi(t))\,\varphi'(t).
  \]
  Par le second théorème fondamental,
  \[
    \int_\alpha^\beta f(\varphi(t))\,\varphi'(t)\,\dd t = F(\varphi(\beta)) - F(\varphi(\alpha)) = F(b) - F(a) = \int_a^b f(x)\,\dd x. \qedhere
  \]
\end{proof}

\begin{exemple}\label{ex:cv1}
  Calculons $\displaystyle\int_0^1 \frac{2t}{1 + t^2}\,\dd t$.
  On pose $x = \varphi(t) = t^2$, d'où $\varphi'(t) = 2t$, $\varphi(0) = 0$, $\varphi(1) = 1$. Alors
  \[
    \int_0^1 \frac{2t}{1+t^2}\,\dd t = \int_0^1 \frac{\dd x}{1 + x} = \bigl[\ln(1+x)\bigr]_0^1 = \ln 2.
  \]
\end{exemple}

\begin{exemple}\label{ex:cv2}
  Calculons $\displaystyle\int_0^1 \sqrt{1 - t^2}\,\dd t$.
  On pose $t = \sin\theta$, $\dd t = \cos\theta\,\dd\theta$, $\theta \in [0, \pi/2]$. Alors
  \[
    \int_0^1 \sqrt{1-t^2}\,\dd t = \int_0^{\pi/2} \cos^2\theta\,\dd\theta = \int_0^{\pi/2}\frac{1 + \cos 2\theta}{2}\,\dd\theta = \frac{\pi}{4}.
  \]
  On retrouve l'aire d'un quart de disque de rayon~$1$.
\end{exemple}

\begin{exemple}\label{ex:cv3}
  Calculons $\displaystyle\int_0^{\pi/2} \frac{\dd t}{1 + \sin t}$.
  On effectue le changement de variable $u = \tan(t/2)$, d'où $\sin t = \frac{2u}{1+u^2}$ et $\dd t = \frac{2\,\dd u}{1+u^2}$. Pour $t = 0$, $u = 0$ ; pour $t = \pi/2$, $u = 1$. Ainsi
  \[
    \int_0^{\pi/2}\frac{\dd t}{1+\sin t}
    = \int_0^1 \frac{1}{1 + \frac{2u}{1+u^2}}\cdot\frac{2\,\dd u}{1+u^2}
    = \int_0^1 \frac{2\,\dd u}{(1+u)^2}
    = \left[-\frac{2}{1+u}\right]_0^1 = -1 + 2 = 1.
  \]
\end{exemple}

\medskip

\noindent\textbf{Exemples classiques.}

\begin{exemple}\label{ex:sin2}
  \emph{Intégrale de $\sin^2 x$.} En utilisant la linéarisation $\sin^2 x = \frac{1 - \cos 2x}{2}$,
  \[
    \int_0^{\pi} \sin^2 x\,\dd x = \frac{1}{2}\int_0^{\pi}(1 - \cos 2x)\,\dd x = \frac{1}{2}\left[\pi - \frac{\sin 2x}{2}\right]_0^{\pi} = \frac{\pi}{2}.
  \]
\end{exemple}

\begin{exemple}\label{ex:arctan}
  \emph{Intégrale de $1/(1+x^2)$.} On a $\displaystyle\int_0^1 \frac{\dd x}{1+x^2} = \bigl[\arctan x\bigr]_0^1 = \frac{\pi}{4}$.
\end{exemple}

\begin{exemple}\label{ex:fractions-rationnelles}
  \emph{Fractions rationnelles.} Calculons $\displaystyle\int_2^3\frac{\dd x}{x^2 - 1}$.
  La décomposition en éléments simples donne $\frac{1}{x^2-1} = \frac{1}{2}\left(\frac{1}{x-1} - \frac{1}{x+1}\right)$, d'où
  \[
    \int_2^3 \frac{\dd x}{x^2-1} = \frac{1}{2}\left[\ln\abs*{\frac{x-1}{x+1}}\right]_2^3 = \frac{1}{2}\left(\ln\frac{1}{2} - \ln\frac{1}{3}\right) = \frac{1}{2}\ln\frac{3}{2}.
  \]
\end{exemple}

% ============================================================
\section{Intégrales impropres (intégrales généralisées)}
\label{sec:integrales-impropres}
\index{intégrale impropre}\index{intégrale généralisée}

Lorsque l'intervalle d'intégration est non borné ou que l'intégrande n'est pas bornée, l'intégrale de Riemann au sens strict ne s'applique plus. On étend la notion par passage à la limite.

\begin{definition}[Intégrale impropre en $+\infty$]\label{def:integrale-impropre-infty}
  Soit $f \colon [a, +\infty[\, \to \R$ continue (ou continue par morceaux). On dit que l'\emph{intégrale impropre}\index{intégrale impropre} $\displaystyle\int_a^{+\infty} f(t)\,\dd t$ \emph{converge} si la limite
  \[
    \lim_{A \to +\infty} \int_a^A f(t)\,\dd t
  \]
  existe et est finie. Cette limite est alors notée $\displaystyle\int_a^{+\infty} f(t)\,\dd t$.
\end{definition}

\begin{definition}[Intégrale impropre en une singularité]\label{def:integrale-impropre-sing}
  Si $f \colon \,]a, b] \to \R$ est continue et non bornée au voisinage de~$a$, on pose
  \[
    \int_a^b f(t)\,\dd t = \lim_{\eps \to 0^+} \int_{a+\eps}^b f(t)\,\dd t,
  \]
  lorsque cette limite existe.
\end{definition}

\begin{theoreme}[Critère de comparaison]\label{thm:comparaison-impropre}
  \index{critère de comparaison}
  Soient $f, g \colon [a, +\infty[\, \to \R$ continues avec $0 \leqslant f(t) \leqslant g(t)$ pour tout $t \geqslant a$.
  \begin{enumerate}[label=(\roman*)]
    \item Si $\displaystyle\int_a^{+\infty} g$ converge, alors $\displaystyle\int_a^{+\infty} f$ converge et $\displaystyle\int_a^{+\infty} f \leqslant \int_a^{+\infty} g$.
    \item Si $\displaystyle\int_a^{+\infty} f$ diverge, alors $\displaystyle\int_a^{+\infty} g$ diverge.
  \end{enumerate}
\end{theoreme}

\begin{proof}
  Posons $F(A) = \int_a^A f(t)\,\dd t$ et $G(A) = \int_a^A g(t)\,\dd t$. Puisque $f \geqslant 0$, la fonction $F$ est croissante. De même pour $G$. L'hypothèse $f \leqslant g$ implique $F(A) \leqslant G(A)$ pour tout $A \geqslant a$.

  \emph{(i)} Si $G(A)$ converge vers $\ell \in \R$, alors $F(A) \leqslant G(A) \leqslant \ell$ pour tout $A$. Donc $F$ est croissante et majorée, d'où convergente, et sa limite satisfait $\int_a^{+\infty} f \leqslant \ell = \int_a^{+\infty} g$.

  \emph{(ii)} Contraposée de~(i).
\end{proof}

\begin{exemple}\label{ex:riemann-integrale}
  \emph{Intégrales de Riemann.}\index{intégrale de Riemann!impropre} Pour $\alpha \in \R$, l'intégrale $\displaystyle\int_1^{+\infty}\frac{\dd t}{t^\alpha}$ converge si et seulement si $\alpha > 1$. En effet, si $\alpha \neq 1$,
  \[
    \int_1^A \frac{\dd t}{t^\alpha} = \frac{A^{1-\alpha} - 1}{1-\alpha},
  \]
  qui converge lorsque $A \to +\infty$ si et seulement si $1 - \alpha < 0$, c'est-à-dire $\alpha > 1$, et la valeur est $\frac{1}{\alpha - 1}$. Si $\alpha = 1$, $\int_1^A \frac{\dd t}{t} = \ln A \to +\infty$.
\end{exemple}

\begin{remarque}\label{rem:gamma-beta}
  La \emph{fonction Gamma d'Euler}\index{fonction Gamma}, définie pour $s > 0$ par
  \[
    \Gamma(s) = \int_0^{+\infty} t^{s-1}\,e^{-t}\,\dd t,
  \]
  est une intégrale impropre convergente. Elle vérifie $\Gamma(n+1) = n!$ pour tout $n \in \N$ et la relation fonctionnelle $\Gamma(s+1) = s\,\Gamma(s)$. La \emph{fonction Bêta}\index{fonction Bêta}, définie pour $p, q > 0$ par
  \[
    B(p,q) = \int_0^1 t^{p-1}(1-t)^{q-1}\,\dd t,
  \]
  est liée à la fonction Gamma par $B(p,q) = \frac{\Gamma(p)\,\Gamma(q)}{\Gamma(p+q)}$.
\end{remarque}

% ============================================================
\section{Formules de Wallis et de Stirling}
\label{sec:wallis-stirling}

\begin{definition}[Intégrales de Wallis]\label{def:wallis}
  \index{intégrales de Wallis}
  Pour $n \in \N$, on pose
  \[
    W_n = \int_0^{\pi/2} \sin^n t\,\dd t = \int_0^{\pi/2} \cos^n t\,\dd t.
  \]
\end{definition}

\begin{proposition}[Formule de récurrence]\label{prop:wallis-recurrence}
  Pour tout $n \geqslant 2$,
  \[
    W_n = \frac{n-1}{n}\,W_{n-2},
  \]
  avec $W_0 = \dfrac{\pi}{2}$ et $W_1 = 1$.
\end{proposition}

\begin{proof}
  Par intégration par parties, en posant $u(t) = \sin^{n-1}t$ et $v'(t) = \sin t$, on a $u'(t) = (n-1)\sin^{n-2}t\cos t$ et $v(t) = -\cos t$. Ainsi
  \begin{align*}
    W_n &= \bigl[-\sin^{n-1}t\,\cos t\bigr]_0^{\pi/2} + (n-1)\int_0^{\pi/2}\sin^{n-2}t\,\cos^2 t\,\dd t \\
    &= 0 + (n-1)\int_0^{\pi/2}\sin^{n-2}t\,(1 - \sin^2 t)\,\dd t \\
    &= (n-1)(W_{n-2} - W_n).
  \end{align*}
  D'où $W_n + (n-1)W_n = (n-1)W_{n-2}$, soit $nW_n = (n-1)W_{n-2}$.
\end{proof}

\begin{theoreme}[Formule de Wallis]\label{thm:wallis}
  \index{formule de Wallis}
  \[
    \lim_{n \to +\infty} \frac{(2n)!!^2}{(2n-1)!!^2 \cdot (2n+1)} = \frac{\pi}{2},
  \]
  ou de manière équivalente, $W_n \sim \sqrt{\dfrac{\pi}{2n}}$ lorsque $n \to +\infty$.
\end{theoreme}

\begin{theoreme}[Formule de Stirling]\label{thm:stirling}
  \index{formule de Stirling}
  Lorsque $n \to +\infty$,
  \[
    n! \sim \sqrt{2\pi n}\left(\frac{n}{e}\right)^n.
  \]
\end{theoreme}

% ============================================================
\section{Exercices}
\label{sec:exercices-chap7}

\begin{exercice}
  Calculer les intégrales suivantes :
  \begin{enumerate}[label=(\alph*)]
    \item $\displaystyle\int_0^1 (3t^2 - 2t + 1)\,\dd t$ ;
    \item $\displaystyle\int_0^{\pi/2} \cos^3 t\,\dd t$ ;
    \item $\displaystyle\int_1^2 \frac{\ln t}{t}\,\dd t$.
  \end{enumerate}
\end{exercice}

\begin{exercice}
  Soit $f \colon [0,1] \to \R$ définie par $f(x) = x^2$. Calculer explicitement $s(f, \sigma_n)$ et $S(f, \sigma_n)$ pour la subdivision régulière $\sigma_n$ d'ordre~$n$, et vérifier que les deux suites convergent vers $\frac{1}{3}$.
\end{exercice}

\begin{exercice}
  Montrer que si $f \colon [a,b] \to \R$ est continue et $f \geqslant 0$, avec $\displaystyle\int_a^b f = 0$, alors $f = 0$.
\end{exercice}

\begin{exercice}[$\star\star$]
  Soit $f \colon [0,1] \to \R$ continue. Montrer que
  \[
    \lim_{n \to +\infty} \frac{1}{n}\sum_{k=1}^{n} f\!\left(\frac{k}{n}\right) = \int_0^1 f(x)\,\dd x.
  \]
  \emph{(Sommes de Riemann.)}\index{sommes de Riemann}
\end{exercice}

\begin{exercice}[$\star$]
  Calculer par intégration par parties : $\displaystyle\int_0^1 \arctan t\,\dd t$.
\end{exercice}

\begin{exercice}[$\star$]
  Calculer par changement de variable : $\displaystyle\int_0^{\pi/4}\frac{\sin t}{\cos^2 t}\,\dd t$.
\end{exercice}

\begin{exercice}[$\star\star$]
  Calculer $\displaystyle\int_0^{+\infty} \frac{\dd t}{(1+t^2)^2}$ par le changement de variable $t = \tan\theta$.
\end{exercice}

\begin{exercice}[$\star\star$]
  Montrer que $\displaystyle\int_0^{+\infty} e^{-t^2}\,\dd t$ converge. \emph{(On ne demande pas la valeur exacte.)}
\end{exercice}

\begin{exercice}[$\star$]
  Étudier la convergence de $\displaystyle\int_0^1\frac{\dd t}{\sqrt{t}}$ et calculer sa valeur.
\end{exercice}

\begin{exercice}[$\star\star$]
  Soit $f \colon [a,b] \to \R$ de classe $\mathcal{C}^1$ avec $f(a) = 0$. Montrer que
  \[
    \int_a^b f(t)^2\,\dd t \leqslant (b-a)^2 \int_a^b f'(t)^2\,\dd t.
  \]
  \emph{(Inégalité de Poincaré.)}\index{inégalité de Poincaré}
\end{exercice}

\begin{exercice}[$\star\star\star$]
  Soit $f \colon [0,1] \to \R$ continue et positive. Montrer que
  \[
    \lim_{n \to +\infty}\left(\int_0^1 f(x)^n\,\dd x\right)^{1/n} = \sup_{x \in [0,1]} f(x).
  \]
\end{exercice}

\begin{exercice}[$\star\star\star$]
  Montrer la formule de Stirling : $n! \sim \sqrt{2\pi n}\left(\frac{n}{e}\right)^n$, en utilisant les intégrales de Wallis et la méthode de Laplace.
\end{exercice}

\bigskip
\begin{center}
\fbox{\parbox{0.85\textwidth}{%
\textbf{Résumé du chapitre~\thechapter.}
\begin{itemize}[nosep]
  \item L'intégrale de Riemann est définie via les sommes de Darboux : $f$ est intégrable si $\sup s(f,\sigma) = \inf S(f,\sigma)$.
  \item Toute fonction continue ou monotone sur $[a,b]$ est intégrable.
  \item L'intégrale est linéaire, positive, croissante et vérifie la relation de Chasles.
  \item Le premier théorème fondamental : $F(x) = \int_a^x f$ est une primitive de~$f$.
  \item Le second théorème fondamental (Newton--Leibniz) : $\int_a^b f = G(b) - G(a)$.
  \item Techniques : intégration par parties et changement de variable.
  \item Les intégrales impropres étendent l'intégrale aux intervalles non bornés ou aux fonctions non bornées.
\end{itemize}}}
\end{center}

% ############################################################
%  CHAPITRE 8 : SUITES ET SÉRIES DE FONCTIONS
% ############################################################
\chapter{Suites et Séries de Fonctions --- Convergence Uniforme}
\label{chap:suites-series-fonctions}
\index{convergence uniforme}

\section*{Introduction}
\addcontentsline{toc}{section}{Introduction}

En Analyse~I, nous avons étudié les suites numériques et les séries numériques. Une question naturelle se pose : que se passe-t-il lorsqu'on considère des suites ou des séries dont les termes sont eux-mêmes des \emph{fonctions} ? Si $(f_n)_{n \in \N}$ est une suite de fonctions continues convergeant vers une fonction~$f$, la limite $f$ est-elle encore continue ? Peut-on intervertir limite et intégrale ? Limite et dérivée ?

La réponse est en général \textbf{négative}, comme le montrent des contre-exemples classiques. C'est la notion de \emph{convergence uniforme}\index{convergence uniforme}, distincte de la convergence simple, qui fournit les conditions suffisantes pour ces interversions. Ce chapitre est consacré à l'étude de cette notion fondamentale.

% ============================================================
\section{Convergence simple et convergence uniforme}
\label{sec:cvs-cvu}

Soit $A$ un sous-ensemble de~$\R$ et $(f_n)_{n \in \N}$ une suite de fonctions de $A$ dans~$\R$.

\begin{definition}[Convergence simple]\label{def:convergence-simple}
  \index{convergence simple}
  La suite $(f_n)$ \emph{converge simplement} (ou \emph{ponctuellement}) vers $f \colon A \to \R$ si
  \[
    \forall\, x \in A,\quad \lim_{n \to +\infty} f_n(x) = f(x),
  \]
  c'est-à-dire :
  \[
    \forall\, x \in A, \quad \forall\, \eps > 0, \quad \exists\, N \in \N, \quad \forall\, n \geqslant N, \quad \abs{f_n(x) - f(x)} < \eps.
  \]
  L'entier $N$ dépend en général de $x$ \textbf{et} de~$\eps$.
\end{definition}

\begin{definition}[Convergence uniforme]\label{def:convergence-uniforme}
  \index{convergence uniforme!définition}
  La suite $(f_n)$ \emph{converge uniformément} vers $f \colon A \to \R$ si
  \[
    \forall\, \eps > 0, \quad \exists\, N \in \N, \quad \forall\, n \geqslant N, \quad \forall\, x \in A, \quad \abs{f_n(x) - f(x)} < \eps,
  \]
  autrement dit, si $\displaystyle\sup_{x \in A}\abs{f_n(x) - f(x)} \xrightarrow[n \to +\infty]{} 0$.
\end{definition}

\begin{remarque}\label{rem:cvu-cvs}
  La convergence uniforme implique la convergence simple, mais la réciproque est fausse. La différence cruciale est que dans la convergence uniforme, l'entier $N$ ne dépend que de $\eps$, et non de~$x$. On peut reformuler : la convergence uniforme signifie que $\norm{f_n - f}_\infty \to 0$, où $\norm{\cdot}_\infty$ est la norme de la convergence uniforme (ou norme du supremum).
\end{remarque}

\begin{center}
\begin{tikzpicture}[scale=1]
  \begin{axis}[
    axis lines=middle,
    xlabel={$x$},
    ylabel={$y$},
    xmin=-0.2, xmax=1.3,
    ymin=-0.3, ymax=1.5,
    xtick={0,1},
    ytick={0,1},
    width=11cm, height=7cm,
    title={Convergence uniforme : le graphe de $f_n$ reste dans le tube de largeur $2\eps$},
    every axis title/.style={at={(0.5,1.08)}},
    legend style={at={(0.97,0.5)}, anchor=east, font=\small},
  ]
    % Fonction limite
    \addplot[domain=0:1, samples=100, very thick, black] {x};
    \addlegendentry{$f(x)=x$}

    % Tube epsilon
    \addplot[domain=0:1, samples=50, dashed, blue!60, thick] {x + 0.15};
    \addplot[domain=0:1, samples=50, dashed, blue!60, thick] {x - 0.15};

    % Fonctions approchantes
    \addplot[domain=0:1, samples=100, red!70, thick] {x + 0.1*sin(deg(6*x))};
    \addlegendentry{$f_n$ (dans le tube)}

    % Annotations
    \draw[<->, blue!60, thick] (axis cs:0.5,0.35) -- (axis cs:0.5,0.65) node[midway, right, blue!80] {$2\eps$};
  \end{axis}
\end{tikzpicture}
\end{center}

\begin{theoreme}[Critère de Cauchy uniforme]\label{thm:cauchy-uniforme}
  \index{critère de Cauchy uniforme}
  La suite $(f_n)$ converge uniformément sur $A$ si et seulement si
  \[
    \forall\, \eps > 0, \quad \exists\, N \in \N, \quad \forall\, m, n \geqslant N, \quad \forall\, x \in A, \quad \abs{f_m(x) - f_n(x)} < \eps.
  \]
\end{theoreme}

\begin{proof}
  \emph{Condition nécessaire.} Si $(f_n)$ converge uniformément vers~$f$, soit $\eps > 0$. Il existe $N$ tel que pour tout $n \geqslant N$ et tout $x \in A$, $\abs{f_n(x) - f(x)} < \eps/2$. Pour $m, n \geqslant N$,
  \[
    \abs{f_m(x) - f_n(x)} \leqslant \abs{f_m(x) - f(x)} + \abs{f(x) - f_n(x)} < \frac{\eps}{2} + \frac{\eps}{2} = \eps.
  \]

  \emph{Condition suffisante.} Supposons la condition de Cauchy uniforme satisfaite. Pour chaque $x \in A$, la suite numérique $(f_n(x))$ est de Cauchy dans~$\R$ (complet), donc converge. Posons $f(x) = \lim_{n} f_n(x)$.

  Soit $\eps > 0$ et $N$ donné par la condition de Cauchy. Pour $m, n \geqslant N$ et $x \in A$, $\abs{f_m(x) - f_n(x)} < \eps$. En faisant tendre $m \to +\infty$ à $n$ et $x$ fixés, on obtient $\abs{f(x) - f_n(x)} \leqslant \eps$ pour tout $n \geqslant N$ et tout $x \in A$. Ceci montre la convergence uniforme.
\end{proof}

\begin{exemple}[Contre-exemple fondamental]\label{ex:xn-pas-cvu}
  \index{convergence simple!sans convergence uniforme}
  Considérons $f_n \colon [0,1] \to \R$ définie par $f_n(x) = x^n$. Pour tout $x \in [0,1[$, on a $f_n(x) = x^n \to 0$, et $f_n(1) = 1 \to 1$. La suite converge simplement vers
  \[
    f(x) = \begin{cases} 0 & \text{si } x \in [0,1[, \\ 1 & \text{si } x = 1. \end{cases}
  \]
  La convergence n'est \textbf{pas} uniforme sur $[0,1]$. En effet,
  \[
    \sup_{x \in [0,1]}\abs{f_n(x) - f(x)} \geqslant \abs{f_n(1 - 1/n) - f(1-1/n)} = \left(1 - \frac{1}{n}\right)^n \xrightarrow[n \to +\infty]{} \frac{1}{e} \neq 0.
  \]
  De plus, chaque $f_n$ est continue, mais la limite $f$ est discontinue en $x = 1$ : la convergence simple de fonctions continues ne préserve pas la continuité.
\end{exemple}

% ============================================================
\section{Propriétés de la convergence uniforme}
\label{sec:proprietes-cvu}

Le théorème suivant est l'un des résultats centraux de ce chapitre.

\begin{theoreme}[Limite uniforme de fonctions continues]\label{thm:limite-uniforme-continue}
  \index{convergence uniforme!et continuité}
  Soit $(f_n)$ une suite de fonctions continues sur $A \subset \R$ convergeant uniformément vers $f \colon A \to \R$. Alors $f$ est continue sur~$A$.
\end{theoreme}

\begin{proof}
  Soit $x_0 \in A$ et $\eps > 0$. Puisque la convergence est uniforme, il existe $N \in \N$ tel que
  \[
    \forall\, n \geqslant N, \quad \forall\, x \in A, \quad \abs{f_n(x) - f(x)} < \frac{\eps}{3}.
  \]
  Fixons un tel $n \geqslant N$. La fonction $f_n$ est continue en $x_0$, donc il existe $\delta > 0$ tel que
  \[
    \forall\, x \in A, \quad \abs{x - x_0} < \delta \implies \abs{f_n(x) - f_n(x_0)} < \frac{\eps}{3}.
  \]
  Pour $x \in A$ avec $\abs{x - x_0} < \delta$, on a alors
  \begin{align*}
    \abs{f(x) - f(x_0)}
    &\leqslant \abs{f(x) - f_n(x)} + \abs{f_n(x) - f_n(x_0)} + \abs{f_n(x_0) - f(x_0)} \\
    &< \frac{\eps}{3} + \frac{\eps}{3} + \frac{\eps}{3} = \eps.
  \end{align*}
  Ceci montre que $f$ est continue en~$x_0$. Comme $x_0$ est arbitraire, $f$ est continue sur~$A$.
\end{proof}

\begin{remarque}\label{rem:contreexemple-cvs-continuite}
  Le contre-exemple~\ref{ex:xn-pas-cvu} ($f_n(x) = x^n$ sur $[0,1]$) montre que l'hypothèse de convergence uniforme est essentielle : la convergence simple de fonctions continues peut produire une limite discontinue.
\end{remarque}

\begin{theoreme}[Interversion limite-intégrale]\label{thm:interversion-limite-integrale}
  \index{interversion!limite-intégrale}
  Soit $(f_n)$ une suite de fonctions intégrables au sens de Riemann sur $[a,b]$ convergeant uniformément vers $f \colon [a,b] \to \R$. Alors $f$ est intégrable sur $[a,b]$ et
  \[
    \lim_{n \to +\infty} \int_a^b f_n(t)\,\dd t = \int_a^b f(t)\,\dd t = \int_a^b \lim_{n \to +\infty} f_n(t)\,\dd t.
  \]
\end{theoreme}

\begin{proof}
  Montrons d'abord que $f$ est intégrable. Soit $\eps > 0$. Par convergence uniforme, il existe $N$ tel que pour tout $n \geqslant N$ et tout $x \in [a,b]$, $\abs{f_n(x) - f(x)} < \frac{\eps}{3(b-a)}$, c'est-à-dire
  \[
    f_n(x) - \frac{\eps}{3(b-a)} < f(x) < f_n(x) + \frac{\eps}{3(b-a)}.
  \]
  Fixons un tel $n \geqslant N$. Puisque $f_n$ est intégrable, il existe une subdivision $\sigma$ telle que $S(f_n,\sigma) - s(f_n,\sigma) < \eps/3$. Sur chaque sous-intervalle $[x_{i-1},x_i]$,
  \[
    \sup_{[x_{i-1},x_i]} f - \inf_{[x_{i-1},x_i]} f
    \leqslant \left(\sup_{[x_{i-1},x_i]} f_n + \frac{\eps}{3(b-a)}\right) - \left(\inf_{[x_{i-1},x_i]} f_n - \frac{\eps}{3(b-a)}\right).
  \]
  D'où
  \begin{align*}
    S(f,\sigma) - s(f,\sigma)
    &\leqslant S(f_n,\sigma) - s(f_n,\sigma) + \frac{2\eps}{3(b-a)}\sum_i (x_i - x_{i-1}) \\
    &< \frac{\eps}{3} + \frac{2\eps}{3} = \eps.
  \end{align*}
  Donc $f$ est intégrable.

  Pour la convergence des intégrales, on a
  \[
    \abs*{\int_a^b f_n - \int_a^b f}
    = \abs*{\int_a^b (f_n - f)}
    \leqslant \int_a^b \abs{f_n - f}
    \leqslant (b-a)\,\sup_{[a,b]}\abs{f_n - f}.
  \]
  Comme $\sup_{[a,b]}\abs{f_n - f} \to 0$ par convergence uniforme, on conclut $\int_a^b f_n \to \int_a^b f$.
\end{proof}

\begin{theoreme}[Interversion limite-dérivée]\label{thm:interversion-limite-derivee}
  \index{interversion!limite-dérivée}
  Soit $(f_n)$ une suite de fonctions de classe $\mathcal{C}^1$ sur $[a,b]$ vérifiant :
  \begin{enumerate}[label=(\roman*)]
    \item il existe $x_0 \in [a,b]$ tel que la suite $(f_n(x_0))$ converge ;
    \item la suite $(f_n')$ converge uniformément sur $[a,b]$ vers une fonction $g$.
  \end{enumerate}
  Alors $(f_n)$ converge uniformément sur $[a,b]$ vers une fonction $f$ de classe $\mathcal{C}^1$, et $f' = g$, c'est-à-dire
  \[
    \left(\lim_{n \to +\infty} f_n\right)' = \lim_{n \to +\infty} f_n'.
  \]
\end{theoreme}

\begin{proof}
  Posons $\ell = \lim_{n} f_n(x_0)$. Par le théorème fondamental, pour tout $x \in [a,b]$,
  \[
    f_n(x) = f_n(x_0) + \int_{x_0}^x f_n'(t)\,\dd t.
  \]
  Puisque $(f_n')$ converge uniformément vers $g$ sur $[a,b]$, le théorème~\ref{thm:interversion-limite-integrale} donne
  \[
    \int_{x_0}^x f_n'(t)\,\dd t \xrightarrow[n \to +\infty]{} \int_{x_0}^x g(t)\,\dd t
  \]
  uniformément en $x$ (car $\abs{\int_{x_0}^x (f_n' - g)} \leqslant (b-a)\norm{f_n'-g}_\infty$). Donc $(f_n)$ converge uniformément vers
  \[
    f(x) = \ell + \int_{x_0}^x g(t)\,\dd t.
  \]
  La fonction $g$ est continue sur $[a,b]$ (comme limite uniforme de fonctions continues $f_n'$, par le théorème~\ref{thm:limite-uniforme-continue}). Par le premier théorème fondamental, $f$ est de classe $\mathcal{C}^1$ et $f' = g$.
\end{proof}

% ============================================================
\section{Séries de fonctions}
\label{sec:series-fonctions}
\index{série de fonctions}

\begin{definition}[Convergence d'une série de fonctions]\label{def:convergence-serie-fonctions}
  Soit $(u_n)_{n \in \N}$ une suite de fonctions de $A$ dans~$\R$. La \emph{série de fonctions} $\sum u_n$ est la suite des sommes partielles $S_N = \sum_{n=0}^{N} u_n$. On dit que la série
  \begin{itemize}[nosep]
    \item \emph{converge simplement} sur $A$ si la suite $(S_N)$ converge simplement ;
    \item \emph{converge uniformément} sur $A$ si la suite $(S_N)$ converge uniformément ;
    \item \emph{converge normalement}\index{convergence normale} sur $A$ si la série numérique $\sum_{n \geqslant 0} \norm{u_n}_\infty$ converge, où $\norm{u_n}_\infty = \sup_{x \in A}\abs{u_n(x)}$.
  \end{itemize}
\end{definition}

\begin{proposition}\label{prop:normale-uniforme}
  La convergence normale entraîne la convergence uniforme, qui entraîne la convergence simple. Les réciproques sont fausses en général.
\end{proposition}

\begin{proof}
  Supposons la convergence normale. Pour $N < M$,
  \[
    \sup_{x \in A}\abs{S_M(x) - S_N(x)} = \sup_{x \in A}\abs*{\sum_{n=N+1}^{M} u_n(x)}
    \leqslant \sum_{n=N+1}^{M} \norm{u_n}_\infty.
  \]
  Comme la série $\sum \norm{u_n}_\infty$ converge, son reste tend vers~$0$, donc $(S_N)$ vérifie le critère de Cauchy uniforme. L'implication convergence uniforme $\Rightarrow$ convergence simple est immédiate.
\end{proof}

\begin{theoreme}[Critère de Weierstrass --- test $M$]\label{thm:weierstrass-M}
  \index{critère de Weierstrass}\index{test $M$}
  Soit $(u_n)$ une suite de fonctions de $A$ dans~$\R$. S'il existe une suite $(M_n)$ de réels positifs telle que
  \[
    \forall\, n \in \N, \quad \forall\, x \in A, \quad \abs{u_n(x)} \leqslant M_n
    \qquad \text{et} \qquad \sum_{n=0}^{+\infty} M_n < +\infty,
  \]
  alors la série $\sum u_n$ converge normalement (donc uniformément) sur~$A$.
\end{theoreme}

\begin{proof}
  On a $\norm{u_n}_\infty = \sup_{x \in A}\abs{u_n(x)} \leqslant M_n$ pour tout~$n$. La convergence de $\sum M_n$ entraîne celle de $\sum \norm{u_n}_\infty$ par comparaison de séries à termes positifs. Donc $\sum u_n$ converge normalement.
\end{proof}

Les théorèmes d'interversion s'appliquent naturellement aux séries de fonctions via leurs sommes partielles :

\begin{corollaire}[Continuité d'une somme de série]\label{cor:continuite-serie}
  Si chaque $u_n$ est continue sur $A$ et si $\sum u_n$ converge uniformément sur~$A$, alors $\sum_{n=0}^{+\infty} u_n$ est continue sur~$A$.
\end{corollaire}

\begin{corollaire}[Intégration terme à terme]\label{cor:integration-terme-a-terme}
  \index{intégration terme à terme}
  Si chaque $u_n$ est intégrable sur $[a,b]$ et si $\sum u_n$ converge uniformément sur $[a,b]$, alors
  \[
    \int_a^b \sum_{n=0}^{+\infty} u_n(t)\,\dd t = \sum_{n=0}^{+\infty} \int_a^b u_n(t)\,\dd t.
  \]
\end{corollaire}

\begin{corollaire}[Dérivation terme à terme]\label{cor:derivation-terme-a-terme}
  \index{dérivation terme à terme}
  Si chaque $u_n$ est de classe $\mathcal{C}^1$ sur $[a,b]$, si $\sum u_n(x_0)$ converge en un point $x_0$, et si $\sum u_n'$ converge uniformément sur $[a,b]$, alors $\sum u_n$ converge uniformément vers une fonction de classe $\mathcal{C}^1$ et
  \[
    \left(\sum_{n=0}^{+\infty} u_n\right)' = \sum_{n=0}^{+\infty} u_n'.
  \]
\end{corollaire}

% ============================================================
\section{Séries entières}
\label{sec:series-entieres}
\index{série entière}

\begin{definition}[Série entière]\label{def:serie-entiere}
  On appelle \emph{série entière} toute série de fonctions de la forme $\sum_{n \geqslant 0} a_n x^n$, où $(a_n)_{n \in \N}$ est une suite de réels (ou de complexes) appelés les \emph{coefficients} de la série.
\end{definition}

\begin{definition}[Rayon de convergence]\label{def:rayon-convergence}
  \index{rayon de convergence}
  Le \emph{rayon de convergence} de la série entière $\sum a_n x^n$ est
  \[
    R = \frac{1}{\limsup_{n \to +\infty} \abs{a_n}^{1/n}} \in [0, +\infty].
  \]
  La série converge absolument pour $\abs{x} < R$ et diverge pour $\abs{x} > R$.
\end{definition}

\begin{theoreme}\label{thm:cvu-compact-se}
  \index{série entière!convergence uniforme}
  Si $R > 0$ est le rayon de convergence de $\sum a_n x^n$, alors pour tout $0 < r < R$, la série converge normalement (donc uniformément) sur $[-r, r]$.
\end{theoreme}

\begin{proof}
  Soit $0 < r < R$. Choisissons $\rho$ tel que $r < \rho < R$. Puisque $\sum a_n \rho^n$ converge absolument, la suite $(\abs{a_n}\rho^n)$ est bornée : il existe $M > 0$ tel que $\abs{a_n}\rho^n \leqslant M$ pour tout~$n$. Pour $\abs{x} \leqslant r$,
  \[
    \abs{a_n x^n} \leqslant \abs{a_n}\, r^n = \abs{a_n}\,\rho^n \cdot \left(\frac{r}{\rho}\right)^n \leqslant M \left(\frac{r}{\rho}\right)^n.
  \]
  Puisque $r/\rho < 1$, la série géométrique $\sum M(r/\rho)^n$ converge. Le critère de Weierstrass conclut.
\end{proof}

\begin{theoreme}[Dérivation des séries entières]\label{thm:derivation-se}
  \index{série entière!dérivation}
  Soit $\sum a_n x^n$ une série entière de rayon de convergence $R > 0$. Alors sa somme $f(x) = \sum_{n=0}^{+\infty} a_n x^n$ est de classe $\mathcal{C}^\infty$ sur $]-R, R[$ et
  \[
    f'(x) = \sum_{n=1}^{+\infty} n\, a_n\, x^{n-1}, \qquad \forall\, x \in \,]-R, R[.
  \]
  De plus, la série dérivée a le même rayon de convergence~$R$.
\end{theoreme}

\begin{proof}
  Montrons que la série dérivée $\sum n a_n x^{n-1}$ a le même rayon de convergence. On utilise la formule de Hadamard. Posons $b_n = (n+1)a_{n+1}$. Alors
  \[
    \abs{b_n}^{1/n} = (n+1)^{1/n}\,\abs{a_{n+1}}^{1/n}.
  \]
  Puisque $(n+1)^{1/n} \to 1$ et $\abs{a_{n+1}}^{1/n} = \bigl(\abs{a_{n+1}}^{1/(n+1)}\bigr)^{(n+1)/n}$, on obtient $\limsup \abs{b_n}^{1/n} = \limsup \abs{a_n}^{1/n}$, donc le rayon de convergence est le même.

  Soit $0 < r < R$. Sur $[-r,r]$, la série $\sum n a_n x^{n-1}$ converge uniformément (par le théorème~\ref{thm:cvu-compact-se} appliqué à la série dérivée). Le corollaire~\ref{cor:derivation-terme-a-terme} donne alors $f'(x) = \sum_{n=1}^{+\infty} n a_n x^{n-1}$ sur $]-r,r[$. Comme ceci vaut pour tout $r < R$, le résultat est vrai sur $]-R,R[$. Par récurrence, $f$ est $\mathcal{C}^\infty$.
\end{proof}

\begin{corollaire}[Intégration des séries entières]\label{cor:integration-se}
  \index{série entière!intégration}
  Sous les mêmes hypothèses, pour tout $x \in \,]-R,R[$,
  \[
    \int_0^x f(t)\,\dd t = \sum_{n=0}^{+\infty} \frac{a_n}{n+1}\,x^{n+1}.
  \]
\end{corollaire}

\begin{theoreme}[Théorème d'Abel radial]\label{thm:abel-radial}
  \index{théorème d'Abel radial}
  Soit $\sum a_n x^n$ une série entière de rayon de convergence $R > 0$. Si la série numérique $\sum a_n R^n$ converge (de somme $S$), alors
  \[
    \lim_{x \to R^-} \sum_{n=0}^{+\infty} a_n x^n = S = \sum_{n=0}^{+\infty} a_n R^n.
  \]
\end{theoreme}

\begin{exemple}\label{ex:series-entieres-classiques}
  \emph{Développements en série entière classiques.}\index{développement en série entière}
  \begin{align*}
    e^x &= \sum_{n=0}^{+\infty}\frac{x^n}{n!}, & R &= +\infty, \\
    \sin x &= \sum_{n=0}^{+\infty}\frac{(-1)^n\,x^{2n+1}}{(2n+1)!}, & R &= +\infty, \\
    \cos x &= \sum_{n=0}^{+\infty}\frac{(-1)^n\,x^{2n}}{(2n)!}, & R &= +\infty, \\
    \frac{1}{1-x} &= \sum_{n=0}^{+\infty} x^n, & R &= 1, \\
    \ln(1+x) &= \sum_{n=1}^{+\infty}\frac{(-1)^{n-1}\,x^n}{n}, & R &= 1, \\
    \arctan x &= \sum_{n=0}^{+\infty}\frac{(-1)^n\,x^{2n+1}}{2n+1}, & R &= 1.
  \end{align*}
\end{exemple}

\begin{theoreme}[Unicité du développement en série entière]\label{thm:unicite-dse}
  \index{série entière!unicité}
  Si $\sum_{n=0}^{+\infty} a_n x^n = \sum_{n=0}^{+\infty} b_n x^n$ pour tout $x$ dans un voisinage de~$0$, alors $a_n = b_n$ pour tout $n \in \N$.
\end{theoreme}

\begin{proof}
  Posons $f(x) = \sum a_n x^n = \sum b_n x^n$. Par le théorème~\ref{thm:derivation-se}, $f$ est $\mathcal{C}^\infty$ et $f^{(n)}(0) = n!\,a_n = n!\,b_n$ pour tout~$n$, d'où $a_n = b_n$.
\end{proof}

% ============================================================
\section{Approximation}
\label{sec:approximation}

\begin{definition}[Polynômes de Bernstein]\label{def:bernstein}
  \index{polynômes de Bernstein}
  Soit $f \colon [0,1] \to \R$ continue. Le $n$-ième \emph{polynôme de Bernstein} de $f$ est
  \[
    B_n(f)(x) = \sum_{k=0}^{n} f\!\left(\frac{k}{n}\right) \binom{n}{k} x^k (1-x)^{n-k}.
  \]
\end{definition}

\begin{theoreme}[Théorème d'approximation de Weierstrass]\label{thm:weierstrass-approximation}
  \index{théorème de Weierstrass!approximation}
  Soit $f \colon [0,1] \to \R$ continue. Alors la suite $(B_n(f))_{n \geqslant 1}$ converge uniformément vers $f$ sur $[0,1]$.

  En d'autres termes, toute fonction continue sur un segment peut être approchée uniformément par des polynômes.
\end{theoreme}

\noindent\textbf{Idée de la preuve (via les polynômes de Bernstein).}
On utilise l'interprétation probabiliste : si $X_1, \ldots, X_n$ sont des variables aléatoires de Bernoulli indépendantes de paramètre~$x$, alors $B_n(f)(x) = \mathbb{E}\bigl[f(\bar{X}_n)\bigr]$ où $\bar{X}_n = \frac{1}{n}\sum_{k=1}^n X_k$. Par la loi des grands nombres, $\bar{X}_n \to x$ en probabilité, et la continuité uniforme de $f$ sur le compact $[0,1]$ permet de conclure que $\mathbb{E}[f(\bar{X}_n)] \to f(x)$ uniformément.

Plus précisément, soit $\eps > 0$. Par continuité uniforme de~$f$, il existe $\delta > 0$ tel que $\abs{s - t} < \delta$ implique $\abs{f(s) - f(t)} < \eps/2$. Soit $M = \norm{f}_\infty$. On décompose
\[
  \abs{B_n(f)(x) - f(x)} \leqslant \sum_{k=0}^{n}\abs*{f\!\left(\frac{k}{n}\right) - f(x)}\binom{n}{k}x^k(1-x)^{n-k}
\]
en séparant les indices $k$ tels que $\abs{k/n - x} < \delta$ (contribution $< \eps/2$) et ceux tels que $\abs{k/n - x} \geqslant \delta$ (contribution majorée par $\frac{2M}{\delta^2}\cdot\frac{x(1-x)}{n} \leqslant \frac{M}{2n\delta^2}$ grâce à l'identité $\sum_{k=0}^n (k/n - x)^2\binom{n}{k}x^k(1-x)^{n-k} = \frac{x(1-x)}{n}$). En choisissant $n$ assez grand, la somme totale est inférieure à~$\eps$.

\begin{center}
\begin{tikzpicture}[scale=1]
  \begin{axis}[
    axis lines=middle,
    xlabel={$x$},
    ylabel={$y$},
    xmin=-0.1, xmax=1.15,
    ymin=-0.2, ymax=1.5,
    xtick={0,0.5,1},
    ytick=\empty,
    width=11cm, height=7cm,
    title={Approximation de $f$ par les polynômes de Bernstein $B_n(f)$},
    every axis title/.style={at={(0.5,1.08)}},
    legend style={at={(0.03,0.97)}, anchor=north west, font=\small},
  ]
    % Fonction cible
    \addplot[domain=0:1, samples=200, very thick, black]
      {abs(x-0.5) + 0.2*sin(deg(5*x))};
    \addlegendentry{$f$}

    % B_3
    \addplot[domain=0:1, samples=100, thick, red!70, dashed]
      {0.5*(1-x)^3 + 0.55*3*x*(1-x)^2 + 0.15*3*x^2*(1-x) + 0.7*x^3};
    \addlegendentry{$B_3(f)$}

    % B_10 (approximation plus fine)
    \addplot[domain=0:1, samples=100, thick, blue!70, densely dotted]
      {0.2*(1-x)^5 + 0.45*x*(1-x)^4*5 + 0.38*x^2*(1-x)^3*10 + 0.15*x^3*(1-x)^2*10 + 0.42*x^4*(1-x)*5 + 0.7*x^5};
    \addlegendentry{$B_{10}(f)$ (approx.)}

  \end{axis}
\end{tikzpicture}
\end{center}

% ============================================================
\section{Exercices}
\label{sec:exercices-chap8}

\begin{exercice}
  Étudier la convergence simple et la convergence uniforme de la suite $(f_n)$ sur $[0,1]$ définie par $f_n(x) = nx\,e^{-nx}$.
\end{exercice}

\begin{exercice}
  Soit $f_n(x) = \frac{x}{1 + nx^2}$ pour $x \in \R$. Montrer que $(f_n)$ converge simplement vers~$0$ mais pas uniformément sur~$\R$.
\end{exercice}

\begin{exercice}[$\star$]
  Soit $f_n(x) = \frac{\sin(nx)}{n}$ pour $x \in \R$. Montrer que $(f_n)$ converge uniformément vers~$0$ sur~$\R$.
\end{exercice}

\begin{exercice}[$\star$]
  Soit $f_n \colon [0,1] \to \R$ définie par $f_n(x) = nx^n(1-x)$. Montrer que $(f_n)$ converge simplement vers~$0$ sur $[0,1]$ mais que $\int_0^1 f_n(x)\,\dd x \not\to 0$. Pourquoi cela ne contredit-il pas le théorème~\ref{thm:interversion-limite-integrale} ?
\end{exercice}

\begin{exercice}[$\star\star$]
  Montrer que la série $\displaystyle\sum_{n=1}^{+\infty}\frac{\sin(nx)}{n^2}$ converge uniformément sur~$\R$ et que sa somme est continue.
\end{exercice}

\begin{exercice}[$\star\star$]
  Montrer que la série $\displaystyle\sum_{n=0}^{+\infty}\frac{x^n}{n!}$ converge uniformément sur tout segment $[-R, R]$ mais pas uniformément sur~$\R$.
\end{exercice}

\begin{exercice}[$\star\star$]
  Soit $f(x) = \displaystyle\sum_{n=1}^{+\infty}\frac{(-1)^{n-1}}{n}\,x^n$ pour $x \in \,]-1, 1]$. Montrer que $f(x) = \ln(1+x)$. En déduire, à l'aide du théorème d'Abel radial, que $\ln 2 = \sum_{n=1}^{+\infty}\frac{(-1)^{n-1}}{n}$.
\end{exercice}

\begin{exercice}[$\star$]
  Montrer que la série $\displaystyle\sum_{n=1}^{+\infty}\frac{x^n}{n}$ converge normalement sur $[-r, r]$ pour tout $0 < r < 1$, mais pas normalement sur $]-1, 1[$.
\end{exercice}

\begin{exercice}[$\star\star$]
  Calculer le rayon de convergence et la somme de $\displaystyle\sum_{n=0}^{+\infty} (n+1)\,x^n$.
\end{exercice}

\begin{exercice}[$\star\star$]
  Montrer que pour tout $x \in \,]-1, 1[$,
  \[
    \sum_{n=0}^{+\infty}\frac{x^{2n+1}}{2n+1} = \frac{1}{2}\ln\frac{1+x}{1-x}.
  \]
\end{exercice}

\begin{exercice}[$\star\star\star$]
  Soit $(f_n)$ une suite de fonctions de classe $\mathcal{C}^1$ sur $[0,1]$ telle que $(f_n')$ converge uniformément vers $g$ et $(f_n(0))$ converge. Montrer que $(f_n)$ converge uniformément et que la limite est de classe $\mathcal{C}^1$ de dérivée~$g$.
\end{exercice}

\begin{exercice}[$\star\star\star$]
  Montrer que la fonction $f(x) = \displaystyle\sum_{n=0}^{+\infty}\frac{\cos(3^n x)}{2^n}$ est continue sur $\R$ mais n'est dérivable en aucun point. \emph{(Fonction de Weierstrass.)}\index{fonction de Weierstrass}
\end{exercice}

\bigskip
\begin{center}
\fbox{\parbox{0.85\textwidth}{%
\textbf{Résumé du chapitre~\thechapter.}
\begin{itemize}[nosep]
  \item Convergence simple : $f_n(x) \to f(x)$ pour chaque $x$ ($N$ dépend de $x$ et $\eps$).
  \item Convergence uniforme : $\norm{f_n - f}_\infty \to 0$ ($N$ ne dépend que de $\eps$).
  \item La convergence uniforme préserve la continuité (théorème clé).
  \item Sous convergence uniforme, on peut intervertir limite et intégrale.
  \item Pour intervertir limite et dérivée, il faut la convergence uniforme des dérivées.
  \item Convergence normale $\Rightarrow$ convergence uniforme (critère de Weierstrass).
  \item Les séries entières convergent uniformément sur tout compact du disque ouvert de convergence.
  \item On peut dériver et intégrer une série entière terme à terme dans le disque ouvert.
  \item Théorème de Weierstrass : toute fonction continue sur un segment est limite uniforme de polynômes.
\end{itemize}}}
\end{center}

% ============================================================
%  Chapitre 9 : Introduction au Calcul Différentiel en Plusieurs Variables
% ============================================================

\chapter{Introduction au Calcul Différentiel en Plusieurs Variables}
\label{chap:calcul-diff-plusieurs-variables}
\index{calcul différentiel!plusieurs variables}

\section*{Introduction}
\addcontentsline{toc}{section}{Introduction}

Les phénomènes naturels dépendent le plus souvent de plusieurs grandeurs simultanément. La température en un point de l'atmosphère est fonction de trois coordonnées spatiales et du temps ; l'énergie cinétique d'un système mécanique dépend des vitesses de chacune de ses particules ; le profit d'une entreprise dépend des prix, des quantités produites et des coûts des facteurs de production. Le passage de l'analyse en une variable à l'analyse en plusieurs variables n'est donc pas un luxe théorique : il est imposé par la nature même des objets que l'on souhaite modéliser.

Ce chapitre étend au cadre des fonctions de~$\R^n$ dans~$\R^p$ les notions fondamentales étudiées en Analyse Réelle~I : limites, continuité, dérivabilité. Nous verrons que certains résultats se généralisent sans difficulté, tandis que d'autres exigent des précautions nouvelles. En particulier, l'existence des \emph{dérivées partielles}\index{dérivée partielle} ne garantit ni la continuité, ni la différentiabilité d'une fonction --- un phénomène sans analogue en dimension un. La notion correcte est celle de \emph{différentiabilité}\index{différentiabilité}, qui fait intervenir une \emph{application linéaire} approchant la fonction au voisinage d'un point. C'est autour de cette idée centrale que s'organise tout le calcul différentiel en plusieurs variables.

% ============================================================
\section{Topologie de \texorpdfstring{$\R^n$}{Rn}}
\label{sec:topologie-Rn}
\index{topologie!de $\R^n$}

\subsection{Normes sur \texorpdfstring{$\R^n$}{Rn}}
\label{subsec:normes-Rn}

Nous munissons~$\R^n$ de sa structure naturelle d'espace vectoriel réel. Un point de~$\R^n$ est un $n$-uplet $x = (x_1, \ldots, x_n)$ avec $x_i \in \R$ pour tout~$i$.

\begin{definition}[Norme sur $\R^n$]\label{def:norme-Rn}
  On appelle \emph{norme}\index{norme} sur $\R^n$ toute application $\norm{\cdot} \colon \R^n \to \R^+$ vérifiant, pour tous $x, y \in \R^n$ et tout $\lambda \in \R$ :
  \begin{enumerate}[label=(\roman*)]
    \item \emph{Séparation} : $\norm{x} = 0 \iff x = 0$ ;
    \item \emph{Homogénéité} : $\norm{\lambda x} = \abs{\lambda} \, \norm{x}$ ;
    \item \emph{Inégalité triangulaire}\index{inégalité triangulaire} : $\norm{x + y} \leqslant \norm{x} + \norm{y}$.
  \end{enumerate}
\end{definition}

\begin{exemple}
  \label{ex:normes-classiques}
  Les trois normes les plus courantes sur~$\R^n$ sont :
  \begin{itemize}
    \item la \emph{norme euclidienne}\index{norme!euclidienne} :
    $\displaystyle \norm{x}_2 = \sqrt{\sum_{i=1}^n x_i^2}$ ;
    \item la \emph{norme~$1$}\index{norme!$\ell^1$} (ou norme de Manhattan) :
    $\displaystyle \norm{x}_1 = \sum_{i=1}^n \abs{x_i}$ ;
    \item la \emph{norme infinie}\index{norme!infinie} (ou norme du sup) :
    $\displaystyle \norm{x}_\infty = \max_{1 \leqslant i \leqslant n} \abs{x_i}$.
  \end{itemize}
\end{exemple}

\begin{remarque}
  \label{rem:inegalites-normes}
  On vérifie aisément que pour tout $x \in \R^n$ :
  \[
    \norm{x}_\infty \leqslant \norm{x}_2 \leqslant \norm{x}_1 \leqslant n \, \norm{x}_\infty.
  \]
  Ces inégalités montrent que les trois normes sont \emph{comparables} au sens que la convergence d'une suite dans l'une entraîne la convergence dans les deux autres.
\end{remarque}

Le résultat fondamental suivant affirme que ce phénomène est général en dimension finie.

\begin{theoreme}[Équivalence des normes en dimension finie]
  \label{thm:equivalence-normes}
  \index{équivalence des normes}
  Sur un espace vectoriel réel de dimension finie, toutes les normes sont équivalentes. Plus précisément, si $\norm{\cdot}_a$ et $\norm{\cdot}_b$ sont deux normes sur~$\R^n$, il existe des constantes $\alpha, \beta > 0$ telles que
  \[
    \alpha \, \norm{x}_a \leqslant \norm{x}_b \leqslant \beta \, \norm{x}_a
    \qquad \text{pour tout } x \in \R^n.
  \]
\end{theoreme}

En conséquence, les notions topologiques (ouverts, fermés, convergence, compacité) sont les mêmes pour toutes les normes sur~$\R^n$. On pourra donc choisir, selon les besoins de la démonstration, la norme la plus commode.

\begin{center}
\begin{tikzpicture}[scale=2.2]
  % Boule norme 1
  \draw[thick, blue!70!black, fill=blue!10, fill opacity=0.3]
    (1,0) -- (0,1) -- (-1,0) -- (0,-1) -- cycle;
  % Boule norme 2
  \draw[thick, red!70!black, fill=red!10, fill opacity=0.3]
    (0,0) circle[radius=1];
  % Boule norme infinie
  \draw[thick, green!50!black, fill=green!10, fill opacity=0.3]
    (-1,-1) rectangle (1,1);
  % Axes
  \draw[->] (-1.4,0) -- (1.4,0) node[right] {$x_1$};
  \draw[->] (0,-1.4) -- (0,1.4) node[above] {$x_2$};
  % Légende
  \node[blue!70!black] at (0.55,0.55) {\small $\norm{\cdot}_1$};
  \node[red!70!black] at (0.85,0.55) {\small $\norm{\cdot}_2$};
  \node[green!50!black] at (1.12,0.85) {\small $\norm{\cdot}_\infty$};
  \node[below] at (0,-1.6) {\small Boules unités fermées dans $\R^2$ pour les trois normes classiques};
\end{tikzpicture}
\end{center}

\subsection{Boules, ouverts, fermés}
\label{subsec:boules-ouverts-fermes}

Dans toute la suite, sauf mention contraire, $\R^n$ est muni d'une norme quelconque, notée~$\norm{\cdot}$.

\begin{definition}[Boule ouverte et boule fermée]\label{def:boules}
  Soient $a \in \R^n$ et $r > 0$. La \emph{boule ouverte}\index{boule!ouverte} de centre~$a$ et de rayon~$r$ est
  \[
    B(a, r) = \{ x \in \R^n : \norm{x - a} < r \}.
  \]
  La \emph{boule fermée}\index{boule!fermée} de centre~$a$ et de rayon~$r$ est
  \[
    \overline{B}(a, r) = \{ x \in \R^n : \norm{x - a} \leqslant r \}.
  \]
\end{definition}

\begin{definition}[Ouvert, fermé]\label{def:ouvert-ferme}
  Soit $\Omega \subset \R^n$.
  \begin{enumerate}[label=(\roman*)]
    \item $\Omega$ est \emph{ouvert}\index{ouvert} si pour tout $a \in \Omega$, il existe $r > 0$ tel que $B(a, r) \subset \Omega$.
    \item $\Omega$ est \emph{fermé}\index{fermé} si son complémentaire $\R^n \setminus \Omega$ est ouvert, ou de manière équivalente, si $\Omega$ contient toutes ses valeurs d'adhérence.
  \end{enumerate}
\end{definition}

\begin{definition}[Partie compacte]\label{def:compact-Rn}
  Une partie $K \subset \R^n$ est dite \emph{compacte}\index{compact} si elle est fermée et bornée (c'est-à-dire contenue dans une boule de rayon fini). De manière équivalente, $K$ est compacte si et seulement si de toute suite d'éléments de~$K$ on peut extraire une sous-suite convergente dont la limite appartient à~$K$.
\end{definition}

\subsection{Suites dans \texorpdfstring{$\R^n$}{Rn}}
\label{subsec:suites-Rn}

\begin{proposition}[Convergence composante par composante]
  \label{prop:convergence-composante}
  \index{convergence!composante par composante}
  Soit $(x^{(k)})_{k \in \N}$ une suite de~$\R^n$, où $x^{(k)} = (x_1^{(k)}, \ldots, x_n^{(k)})$, et soit $\ell = (\ell_1, \ldots, \ell_n) \in \R^n$. Alors
  \[
    x^{(k)} \xrightarrow[k \to +\infty]{} \ell
    \quad \iff \quad
    \forall\, i \in \{1, \ldots, n\},\;
    x_i^{(k)} \xrightarrow[k \to +\infty]{} \ell_i.
  \]
\end{proposition}

\begin{proof}
  $(\Rightarrow)$ Pour tout~$i$, on a $\abs{x_i^{(k)} - \ell_i} \leqslant \norm{x^{(k)} - \ell}_\infty \leqslant \norm{x^{(k)} - \ell}$ (en utilisant l'équivalence des normes). Donc $\norm{x^{(k)} - \ell} \to 0$ entraîne $x_i^{(k)} \to \ell_i$.

  $(\Leftarrow)$ Réciproquement, $\norm{x^{(k)} - \ell}_\infty = \max_i \abs{x_i^{(k)} - \ell_i} \to 0$ si chaque composante converge.
\end{proof}

% ============================================================
\section{Limites et continuité dans \texorpdfstring{$\R^n$}{Rn}}
\label{sec:limites-continuite-Rn}
\index{continuité!dans $\R^n$}

\subsection{Définitions}

\begin{definition}[Limite d'une fonction de plusieurs variables]\label{def:limite-plusieurs-variables}
  Soient $\Omega \subset \R^n$ un ouvert, $a \in \Omega$ (ou $a$ point d'accumulation de~$\Omega$), $f \colon \Omega \setminus \{a\} \to \R^p$ et $\ell \in \R^p$. On dit que $f$ admet la \emph{limite}\index{limite!fonction de plusieurs variables}~$\ell$ en~$a$, et on écrit $\lim_{x \to a} f(x) = \ell$, si
  \[
    \forall\, \eps > 0,\; \exists\, \delta > 0,\; \forall\, x \in \Omega,\;
    0 < \norm{x - a} < \delta \implies \norm{f(x) - \ell} < \eps.
  \]
\end{definition}

\begin{definition}[Continuité]\label{def:continuite-Rn}
  Soient $\Omega \subset \R^n$ un ouvert et $f \colon \Omega \to \R^p$. On dit que $f$ est \emph{continue}\index{continuité} en $a \in \Omega$ si $\lim_{x \to a} f(x) = f(a)$. On dit que $f$ est continue sur~$\Omega$ si elle est continue en tout point de~$\Omega$.
\end{definition}

\subsection{Limites itérées et limite double}

Un piège redoutable en dimension supérieure est la confusion entre \emph{limite double} et \emph{limites itérées}\index{limites itérées}.

\begin{exemple}[Limites itérées $\neq$ limite double]
  \label{ex:limites-iterees}
  Considérons la fonction $f \colon \R^2 \setminus \{(0,0)\} \to \R$ définie par
  \[
    f(x,y) = \frac{xy}{x^2 + y^2}.
  \]
  Les limites itérées existent et sont nulles :
  \[
    \lim_{x \to 0} \Bigl( \lim_{y \to 0} f(x,y) \Bigr) = \lim_{x \to 0} 0 = 0,
    \qquad
    \lim_{y \to 0} \Bigl( \lim_{x \to 0} f(x,y) \Bigr) = \lim_{y \to 0} 0 = 0.
  \]
  Cependant, le long de la droite $y = x$, on a $f(x,x) = x^2/(2x^2) = 1/2$, et le long de la droite $y = -x$, on a $f(x,-x) = -1/2$. La limite double $\lim_{(x,y) \to (0,0)} f(x,y)$ \emph{n'existe pas}\index{limite double!non-existence}, bien que les deux limites itérées existent et coïncident.
\end{exemple}

\begin{remarque}
  \label{rem:limites-iterees}
  Cet exemple illustre un principe fondamental : pour montrer qu'une limite double n'existe pas, il suffit de trouver deux chemins d'approche donnant des limites distinctes. En revanche, pour montrer qu'une limite existe, il faut la démontrer pour \emph{tous} les chemins simultanément, ce qui nécessite en général une majoration directe en~$\eps$-$\delta$.
\end{remarque}

\subsection{Continuité et compacité}

Le théorème suivant étend au cadre multidimensionnel le résultat classique des bornes atteintes.

\begin{theoreme}[Bornes atteintes sur un compact]
  \label{thm:bornes-atteintes-compact}
  \index{théorème des bornes atteintes!sur un compact}
  Soient $K \subset \R^n$ un compact non vide et $f \colon K \to \R$ une fonction continue. Alors $f$ est bornée et atteint ses bornes : il existe $a, b \in K$ tels que
  \[
    f(a) = \inf_{x \in K} f(x), \qquad f(b) = \sup_{x \in K} f(x).
  \]
\end{theoreme}

\begin{proof}
  Montrons que $f$ atteint son maximum (le raisonnement est analogue pour le minimum). Posons $M = \sup_{x \in K} f(x)$. Si $M = -\infty$, alors $K = \varnothing$, ce qui est exclu. Sinon, il existe une suite $(x^{(k)})_{k \in \N}$ d'éléments de~$K$ telle que $f(x^{(k)}) \to M$. Comme $K$ est compact, on peut extraire une sous-suite $(x^{(\varphi(k))})$ convergeant vers un point $b \in K$. Par continuité de~$f$, on a $f(b) = \lim_k f(x^{(\varphi(k))}) = M$. En particulier, $M \in \R$ et $f$ est majorée.
\end{proof}

% ============================================================
\section{Dérivées partielles}
\label{sec:derivees-partielles}
\index{dérivée partielle}

\subsection{Définition et notation}

\begin{definition}[Dérivée partielle]\label{def:derivee-partielle}
  Soient $\Omega \subset \R^n$ un ouvert, $f \colon \Omega \to \R$ et $a = (a_1, \ldots, a_n) \in \Omega$. La \emph{dérivée partielle}\index{dérivée partielle} de $f$ par rapport à la $i$-ème variable en~$a$ est, lorsque la limite existe :
  \[
    \frac{\partial f}{\partial x_i}(a)
    = \lim_{h \to 0} \frac{f(a_1, \ldots, a_{i-1}, a_i + h, a_{i+1}, \ldots, a_n) - f(a)}{h}.
  \]
  On note aussi $\partial_i f(a)$, $f_{x_i}(a)$ ou $D_i f(a)$.
\end{definition}

\begin{remarque}
  \label{rem:partielle-gel-variables}
  Calculer $\dfrac{\partial f}{\partial x_i}(a)$ revient à dériver la fonction d'une variable $t \mapsto f(a_1, \ldots, a_{i-1}, t, a_{i+1}, \ldots, a_n)$ en $t = a_i$. Autrement dit, on \emph{gèle} toutes les variables sauf la $i$-ème et on dérive par rapport à celle-ci.
\end{remarque}

\begin{center}
\begin{tikzpicture}[scale=1]
  \begin{axis}[
    axis lines=middle,
    xlabel={$x$}, ylabel={$z$},
    xmin=-0.3, xmax=3.5,
    ymin=-0.5, ymax=3.5,
    xtick={1.5}, xticklabels={$a_1$},
    ytick=\empty,
    width=10cm, height=7cm,
    samples=100,
    clip=false
  ]
    % Courbe f(x, a_2)
    \addplot[thick, blue!70!black, domain=0.2:3.2] {0.3*(x-1)^3 - 0.5*(x-1) + 2};
    % Point
    \addplot[only marks, mark=*, mark size=2.5pt, red!70!black] coordinates {(1.5, 1.9625)};
    % Tangente
    \addplot[thick, red!70!black, dashed, domain=0.5:2.5] {1.9625 + 0.175*(x-1.5)};
    % Labels
    \node[blue!70!black, right] at (axis cs:3.2, 2.8) {$z = f(x, a_2)$};
    \node[red!70!black, above right] at (axis cs:1.5, 1.9625) {$(a_1, f(a))$};
    \node[red!70!black, right] at (axis cs:2.5, 2.2) {\small pente $= \dfrac{\partial f}{\partial x_1}(a)$};
  \end{axis}
\end{tikzpicture}
\captionof{figure}{La dérivée partielle $\dfrac{\partial f}{\partial x_1}(a)$ est la pente de la courbe obtenue en coupant le graphe de~$f$ par le plan $x_2 = a_2$.}
\label{fig:derivee-partielle-tranche}
\end{center}

\subsection{Existence des dérivées partielles et continuité}

\begin{theoreme}
  \label{thm:partials-not-continuity}
  \index{dérivée partielle!ne suffit pas pour la continuité}
  L'existence de toutes les dérivées partielles en un point n'entraîne pas la continuité en ce point.
\end{theoreme}

\begin{exemple}[Contre-exemple fondamental]
  \label{ex:partials-not-continuity}
  Considérons la fonction $f \colon \R^2 \to \R$ définie par
  \[
    f(x,y) = \begin{cases}
      \dfrac{x^2 y}{x^4 + y^2} & \text{si } (x,y) \neq (0,0), \\[6pt]
      0 & \text{si } (x,y) = (0,0).
    \end{cases}
  \]
  \index{contre-exemple!dérivées partielles sans continuité}
\end{exemple}

\begin{proof}[Vérification]
  \emph{Étape 1 : Existence des dérivées partielles en $(0,0)$.}

  Pour $\dfrac{\partial f}{\partial x}(0,0)$ : on calcule
  \[
    \frac{f(h,0) - f(0,0)}{h} = \frac{0 - 0}{h} = 0,
  \]
  donc $\dfrac{\partial f}{\partial x}(0,0) = 0$. De même, $\dfrac{\partial f}{\partial y}(0,0) = 0$.

  \emph{Étape 2 : $f$ n'est pas continue en $(0,0)$.}

  Considérons le chemin $y = x^2$ : pour $x \neq 0$,
  \[
    f(x, x^2) = \frac{x^2 \cdot x^2}{x^4 + x^4} = \frac{x^4}{2x^4} = \frac{1}{2}.
  \]
  Ainsi $\lim_{x \to 0} f(x, x^2) = \dfrac{1}{2} \neq 0 = f(0,0)$.

  Cependant, le long de toute droite passant par l'origine (c'est-à-dire pour $y = mx$ avec $m$ fixé), on a
  \[
    f(x, mx) = \frac{x^2 \cdot mx}{x^4 + m^2 x^2} = \frac{mx}{x^2 + m^2} \xrightarrow[x \to 0]{} 0.
  \]
  La fonction admet donc une limite nulle le long de toute droite, mais n'est pas continue à l'origine. Ceci montre que la vérification le long des droites ne suffit pas à établir l'existence d'une limite double.
\end{proof}

% ============================================================
\section{Différentiabilité}
\label{sec:differentiabilite}
\index{différentiabilité}

Le bon concept de dérivabilité en plusieurs variables est celui de \emph{différentiabilité}, qui fait intervenir une approximation \emph{linéaire} globale (et non direction par direction comme les dérivées partielles).

\subsection{Définition}

\begin{tcolorbox}[
  enhanced,
  colback=red!5!white,
  colframe=red!70!black,
  fonttitle=\bfseries\large,
  title={Définition centrale : Différentiabilité},
  breakable,
  before skip=12pt,
  after skip=12pt,
  left=10pt, right=10pt, top=8pt, bottom=8pt
]
\begin{definition}[Application différentiable]\label{def:differentiable}
  Soient $\Omega \subset \R^n$ un ouvert, $f \colon \Omega \to \R^p$ et $a \in \Omega$. On dit que $f$ est \emph{différentiable}\index{différentiabilité} en~$a$ s'il existe une application linéaire $L \colon \R^n \to \R^p$ telle que
  \[
    f(a + h) = f(a) + L(h) + \norm{h} \, \eps(h),
  \]
  où $\eps(h) \to 0$ quand $h \to 0$ dans~$\R^n$. De manière équivalente :
  \[
    \lim_{h \to 0} \frac{\norm{f(a+h) - f(a) - L(h)}}{\norm{h}} = 0.
  \]
  L'application linéaire~$L$, lorsqu'elle existe, est unique ; on l'appelle la \emph{différentielle}\index{différentielle} de~$f$ en~$a$ et on la note $\dd f_a$ ou $Df(a)$.
\end{definition}
\end{tcolorbox}

\begin{remarque}
  \label{rem:diff-dim-1}
  En dimension $n = p = 1$, la différentiabilité de $f$ en~$a$ signifie l'existence d'un réel~$\lambda$ tel que $f(a+h) = f(a) + \lambda h + o(h)$. On retrouve la dérivabilité ordinaire, avec $\lambda = f'(a)$.
\end{remarque}

\subsection{Différentiabilité implique continuité}

\begin{theoreme}
  \label{thm:diff-implique-cont}
  \index{différentiabilité!implique continuité}
  Si $f \colon \Omega \to \R^p$ est différentiable en~$a$, alors $f$ est continue en~$a$.
\end{theoreme}

\begin{proof}
  Par hypothèse, $f(a+h) = f(a) + \dd f_a(h) + \norm{h}\,\eps(h)$ avec $\eps(h) \to 0$. L'application linéaire $\dd f_a \colon \R^n \to \R^p$ est continue (toute application linéaire entre espaces de dimension finie est continue), donc $\dd f_a(h) \to 0$ quand $h \to 0$. De plus, $\norm{h}\,\eps(h) \to 0$ car $\norm{h} \to 0$ et $\eps(h)$ est bornée au voisinage de~$0$. Il vient
  \[
    \lim_{h \to 0} f(a+h) = f(a) + 0 + 0 = f(a),
  \]
  ce qui prouve la continuité de~$f$ en~$a$.
\end{proof}

\subsection{Différentiabilité et dérivées partielles}

\begin{theoreme}
  \label{thm:diff-implique-partielles}
  \index{différentielle!et dérivées partielles}
  Soient $\Omega \subset \R^n$ un ouvert et $f \colon \Omega \to \R$ différentiable en $a \in \Omega$. Alors toutes les dérivées partielles de~$f$ existent en~$a$ et
  \[
    \dd f_a(h) = \sum_{i=1}^n \frac{\partial f}{\partial x_i}(a) \, h_i
    \qquad \text{pour tout } h = (h_1, \ldots, h_n) \in \R^n.
  \]
\end{theoreme}

\begin{proof}
  Soit $e_i = (0, \ldots, 0, 1, 0, \ldots, 0)$ le $i$-ème vecteur de la base canonique. Par hypothèse de différentiabilité :
  \[
    f(a + t \, e_i) = f(a) + t\,\dd f_a(e_i) + \abs{t}\,\eps(t\,e_i),
  \]
  d'où
  \[
    \frac{f(a + t\,e_i) - f(a)}{t} = \dd f_a(e_i) + \frac{\abs{t}}{t}\,\eps(t\,e_i).
  \]
  Lorsque $t \to 0$, le terme $\eps(t\,e_i) \to 0$, donc
  \[
    \frac{\partial f}{\partial x_i}(a)
    = \lim_{t \to 0} \frac{f(a + t\,e_i) - f(a)}{t}
    = \dd f_a(e_i).
  \]
  Comme $\dd f_a$ est linéaire et déterminée par ses valeurs sur la base canonique :
  \[
    \dd f_a(h) = \dd f_a\!\left(\sum_{i=1}^n h_i \, e_i\right) = \sum_{i=1}^n h_i \, \dd f_a(e_i) = \sum_{i=1}^n \frac{\partial f}{\partial x_i}(a) \, h_i. \qedhere
  \]
\end{proof}

\subsection{Matrice jacobienne}

\begin{definition}[Matrice jacobienne]\label{def:jacobienne}
  Soient $\Omega \subset \R^n$ un ouvert et $f = (f_1, \ldots, f_p) \colon \Omega \to \R^p$ différentiable en $a \in \Omega$. La \emph{matrice jacobienne}\index{matrice jacobienne} de $f$ en~$a$ est la matrice de la différentielle $\dd f_a$ dans les bases canoniques de $\R^n$ et $\R^p$ :
  \[
    J_f(a) = \begin{pmatrix}
      \dfrac{\partial f_1}{\partial x_1}(a) & \cdots & \dfrac{\partial f_1}{\partial x_n}(a) \\[8pt]
      \vdots & \ddots & \vdots \\[4pt]
      \dfrac{\partial f_p}{\partial x_1}(a) & \cdots & \dfrac{\partial f_p}{\partial x_n}(a)
    \end{pmatrix}
    \in \mathcal{M}_{p,n}(\R).
  \]
  Lorsque $n = p$, le \emph{jacobien}\index{jacobien} de~$f$ en~$a$ est le déterminant $\det J_f(a)$.
\end{definition}

% ============================================================
\section{Condition suffisante de différentiabilité : \texorpdfstring{$\mathcal{C}^1$}{C1}}
\label{sec:C1-differentiable}

\begin{theoreme}[$\mathcal{C}^1 \Rightarrow$ différentiable]
  \label{thm:C1-diff}
  \index{$\mathcal{C}^1$!implique différentiable}
  Soient $\Omega \subset \R^n$ un ouvert et $f \colon \Omega \to \R$. Si toutes les dérivées partielles $\dfrac{\partial f}{\partial x_i}$ existent sur~$\Omega$ et sont continues en un point $a \in \Omega$, alors $f$ est différentiable en~$a$.
\end{theoreme}

Nous donnons la démonstration complète dans le cas $n = 2$.

\begin{proof}[Démonstration pour $n = 2$]
  Soient $a = (a_1, a_2)$ et $h = (h_1, h_2)$ avec $a + h \in \Omega$. On écrit
  \[
    f(a_1 + h_1, a_2 + h_2) - f(a_1, a_2)
    = \underbrace{\bigl[f(a_1 + h_1, a_2 + h_2) - f(a_1, a_2 + h_2)\bigr]}_{(I)}
    + \underbrace{\bigl[f(a_1, a_2 + h_2) - f(a_1, a_2)\bigr]}_{(II)}.
  \]

  \emph{Traitement de $(I)$.} La fonction $\varphi \colon t \mapsto f(t, a_2 + h_2)$ est dérivable sur un voisinage de~$a_1$ (car $\frac{\partial f}{\partial x_1}$ existe sur~$\Omega$). Par le théorème des accroissements finis\index{théorème des accroissements finis}, il existe $\theta_1 \in\, ]0,1[$ tel que
  \[
    (I) = h_1 \, \frac{\partial f}{\partial x_1}(a_1 + \theta_1 h_1, \, a_2 + h_2).
  \]
  Par continuité de $\frac{\partial f}{\partial x_1}$ en~$a$, on a
  \[
    \frac{\partial f}{\partial x_1}(a_1 + \theta_1 h_1, \, a_2 + h_2)
    = \frac{\partial f}{\partial x_1}(a) + \eps_1(h),
  \]
  où $\eps_1(h) \to 0$ quand $h \to 0$. Donc $(I) = h_1 \, \dfrac{\partial f}{\partial x_1}(a) + h_1 \, \eps_1(h)$.

  \emph{Traitement de $(II)$.} De même, en appliquant le théorème des accroissements finis à $\psi \colon t \mapsto f(a_1, t)$, il existe $\theta_2 \in\, ]0,1[$ tel que
  \[
    (II) = h_2 \, \frac{\partial f}{\partial x_2}(a_1, a_2 + \theta_2 h_2)
    = h_2 \, \frac{\partial f}{\partial x_2}(a) + h_2 \, \eps_2(h),
  \]
  où $\eps_2(h) \to 0$ quand $h \to 0$.

  \emph{Conclusion.} En sommant :
  \[
    f(a + h) - f(a) = \frac{\partial f}{\partial x_1}(a)\, h_1 + \frac{\partial f}{\partial x_2}(a)\, h_2 + h_1\,\eps_1(h) + h_2\,\eps_2(h).
  \]
  Or $\abs{h_i} \leqslant \norm{h}$ pour $i = 1,2$, donc
  \[
    \abs{h_1\,\eps_1(h) + h_2\,\eps_2(h)} \leqslant \norm{h}\,\bigl(\abs{\eps_1(h)} + \abs{\eps_2(h)}\bigr) = \norm{h}\,\eps(h),
  \]
  où $\eps(h) = \abs{\eps_1(h)} + \abs{\eps_2(h)} \to 0$ quand $h \to 0$.

  Ceci prouve que $f$ est différentiable en~$a$, avec $\dd f_a(h) = \dfrac{\partial f}{\partial x_1}(a)\,h_1 + \dfrac{\partial f}{\partial x_2}(a)\,h_2$.
\end{proof}

% ============================================================
\section{Règle de la chaîne}
\label{sec:regle-chaine}
\index{règle de la chaîne}

\begin{theoreme}[Règle de la chaîne --- Chain rule]
  \label{thm:chain-rule}
  Soient $\Omega \subset \R^n$ et $\Omega' \subset \R^p$ des ouverts, $f \colon \Omega \to \Omega'$ différentiable en~$a$ et $g \colon \Omega' \to \R^q$ différentiable en $f(a)$. Alors $g \circ f$ est différentiable en~$a$ et
  \[
    \dd(g \circ f)_a = \dd g_{f(a)} \circ \dd f_a.
  \]
  En termes de matrices jacobiennes :
  \[
    J_{g \circ f}(a) = J_g(f(a)) \cdot J_f(a).
  \]
\end{theoreme}

\begin{proof}
  Posons $b = f(a)$. Par hypothèse :
  \begin{align*}
    f(a + h) &= f(a) + \dd f_a(h) + \norm{h}\,\eps_1(h), \quad \text{avec } \eps_1(h) \to 0, \\
    g(b + k) &= g(b) + \dd g_b(k) + \norm{k}\,\eps_2(k), \quad \text{avec } \eps_2(k) \to 0.
  \end{align*}
  Posons $k = f(a+h) - f(a) = \dd f_a(h) + \norm{h}\,\eps_1(h)$. Alors
  \[
    g(f(a+h)) = g(b + k) = g(b) + \dd g_b(k) + \norm{k}\,\eps_2(k).
  \]
  On a $\dd g_b(k) = \dd g_b(\dd f_a(h)) + \dd g_b(\norm{h}\,\eps_1(h))$, et par linéarité de $\dd g_b$ :
  \[
    \dd g_b(k) = (\dd g_b \circ \dd f_a)(h) + \norm{h}\,\dd g_b(\eps_1(h)).
  \]
  Comme $\dd g_b$ est linéaire continue et $\eps_1(h) \to 0$, on a $\dd g_b(\eps_1(h)) \to 0$.

  Il reste à montrer que $\norm{k}\,\eps_2(k) = o(\norm{h})$. Il existe une constante $C > 0$ telle que $\norm{\dd f_a(h)} \leqslant C\,\norm{h}$ (car $\dd f_a$ est linéaire continue), d'où
  \[
    \norm{k} \leqslant C\,\norm{h} + \norm{h}\,\norm{\eps_1(h)} = \norm{h}\,(C + \norm{\eps_1(h)}).
  \]
  Pour $h$ suffisamment petit, $C + \norm{\eps_1(h)} \leqslant C + 1$, donc $\norm{k} \leqslant (C+1)\,\norm{h}$. Par conséquent :
  \[
    \frac{\norm{k}\,\norm{\eps_2(k)}}{\norm{h}} \leqslant (C+1)\,\norm{\eps_2(k)} \to 0
  \]
  puisque $k \to 0$ quand $h \to 0$ (par continuité de~$f$), donc $\eps_2(k) \to 0$.

  En rassemblant tous les termes :
  \[
    g(f(a+h)) = g(f(a)) + (\dd g_b \circ \dd f_a)(h) + \norm{h}\,\eta(h),
  \]
  où $\eta(h) \to 0$, ce qui prouve la différentiabilité de $g \circ f$ en~$a$ avec $\dd(g \circ f)_a = \dd g_{f(a)} \circ \dd f_a$.
\end{proof}

\begin{corollaire}
  \label{cor:chain-rule-partielles}
  Si $f = (f_1, \ldots, f_p)$ et $g$ sont comme ci-dessus, avec $f$ et $g$ de classe $\mathcal{C}^1$, alors pour tout $i \in \{1, \ldots, n\}$ :
  \[
    \frac{\partial (g \circ f)}{\partial x_i}(a)
    = \sum_{j=1}^{p} \frac{\partial g}{\partial y_j}(f(a)) \cdot \frac{\partial f_j}{\partial x_i}(a).
  \]
\end{corollaire}

% ============================================================
\section{Gradient et interprétation géométrique}
\label{sec:gradient}
\index{gradient}

\begin{definition}[Gradient]\label{def:gradient}
  Soient $\Omega \subset \R^n$ un ouvert et $f \colon \Omega \to \R$ différentiable en $a \in \Omega$. Le \emph{gradient}\index{gradient} de~$f$ en~$a$ est le vecteur
  \[
    \nabla f(a) = \operatorname{grad} f(a) = \left(\frac{\partial f}{\partial x_1}(a), \ldots, \frac{\partial f}{\partial x_n}(a)\right) \in \R^n.
  \]
  La différentielle s'écrit alors $\dd f_a(h) = \langle \nabla f(a), h \rangle$, où $\langle \cdot, \cdot \rangle$ désigne le produit scalaire euclidien.
\end{definition}

\begin{proposition}[Direction de plus grande pente]
  \label{prop:gradient-direction}
  \index{gradient!direction de plus grande pente}
  Si $\nabla f(a) \neq 0$, alors le gradient $\nabla f(a)$ pointe dans la direction dans laquelle~$f$ croît le plus rapidement à partir de~$a$. Plus précisément, pour tout vecteur unitaire $u \in \R^n$ avec $\norm{u}_2 = 1$, on a
  \[
    \dd f_a(u) = \langle \nabla f(a), u \rangle \leqslant \norm{\nabla f(a)}_2,
  \]
  avec égalité si et seulement si $u = \dfrac{\nabla f(a)}{\norm{\nabla f(a)}_2}$.
\end{proposition}

\begin{proof}
  Par l'inégalité de Cauchy--Schwarz\index{inégalité de Cauchy--Schwarz}, $\langle \nabla f(a), u \rangle \leqslant \norm{\nabla f(a)}_2 \, \norm{u}_2 = \norm{\nabla f(a)}_2$. L'égalité a lieu si et seulement si $u$ et $\nabla f(a)$ sont colinéaires de même sens, c'est-à-dire $u = \nabla f(a) / \norm{\nabla f(a)}_2$.
\end{proof}

\begin{remarque}
  \label{rem:gradient-orthogonal-lignes-niveau}
  Le gradient est \emph{orthogonal} aux lignes de niveau de~$f$. En effet, si $\gamma \colon t \mapsto \gamma(t)$ est une courbe tracée sur la ligne de niveau $\{f = c\}$, alors $f(\gamma(t)) = c$ pour tout~$t$, d'où par la règle de la chaîne : $\langle \nabla f(\gamma(t)), \gamma'(t) \rangle = 0$.
\end{remarque}

\begin{center}
\begin{tikzpicture}[scale=1]
  \begin{axis}[
    view={0}{90},
    xlabel={$x$}, ylabel={$y$},
    xmin=-2.5, xmax=2.5,
    ymin=-2.5, ymax=2.5,
    width=9cm, height=9cm,
    axis lines=middle,
    xtick=\empty, ytick=\empty,
    clip=false
  ]
    % Contour lines for f(x,y) = x^2 + 2y^2
    \foreach \c in {0.5, 1, 2, 3, 4, 5, 6, 7} {
      \addplot[thick, blue!50!black, domain=0:360, samples=80]
        ({sqrt(\c)*cos(x)}, {sqrt(\c/2)*sin(x)});
    }
    % Gradient arrows at selected points
    \draw[-{Stealth[length=6pt]}, red!70!black, very thick] (axis cs:1,0.5) -- (axis cs:1.6,1);
    \draw[-{Stealth[length=6pt]}, red!70!black, very thick] (axis cs:-1,0.7) -- (axis cs:-1.6,1.4);
    \draw[-{Stealth[length=6pt]}, red!70!black, very thick] (axis cs:0.5,-1) -- (axis cs:0.8,-2);
    \draw[-{Stealth[length=6pt]}, red!70!black, very thick] (axis cs:-0.8,-0.5) -- (axis cs:-1.28,-1);
    \draw[-{Stealth[length=6pt]}, red!70!black, very thick] (axis cs:1.2,-0.6) -- (axis cs:1.92,-1.2);
    % Labels
    \node[blue!50!black] at (axis cs:2.1,0.5) {\small $f = c$};
    \node[red!70!black] at (axis cs:2.1,1.3) {\small $\nabla f$};
  \end{axis}
\end{tikzpicture}
\captionof{figure}{Lignes de niveau de $f(x,y) = x^2 + 2y^2$ et vecteurs gradient (en rouge), orthogonaux aux courbes de niveau.}
\label{fig:gradient-contour}
\end{center}

% ============================================================
\section{Dérivées partielles d'ordre supérieur et théorème de Schwarz}
\label{sec:derivees-superieures-schwarz}
\index{dérivées partielles!d'ordre supérieur}

\begin{definition}[Dérivées partielles d'ordre supérieur]\label{def:derivees-superieures}
  Si $\dfrac{\partial f}{\partial x_i}$ est elle-même dérivable par rapport à~$x_j$, on note
  \[
    \frac{\partial^2 f}{\partial x_j \partial x_i} = \frac{\partial}{\partial x_j}\left(\frac{\partial f}{\partial x_i}\right)
  \]
  la dérivée partielle seconde\index{dérivée partielle!seconde} de~$f$ par rapport à~$x_i$ puis~$x_j$. On définit de même les dérivées d'ordre supérieur par récurrence.
\end{definition}

\begin{theoreme}[Schwarz]
  \label{thm:schwarz}
  \index{théorème de Schwarz}
  Soient $\Omega \subset \R^n$ un ouvert et $f \colon \Omega \to \R$. Si les dérivées partielles $\dfrac{\partial^2 f}{\partial x_i \partial x_j}$ et $\dfrac{\partial^2 f}{\partial x_j \partial x_i}$ existent sur~$\Omega$ et sont continues en un point $a \in \Omega$, alors
  \[
    \frac{\partial^2 f}{\partial x_i \partial x_j}(a) = \frac{\partial^2 f}{\partial x_j \partial x_i}(a).
  \]
\end{theoreme}

\begin{proof}
  On se ramène au cas $n = 2$, $i = 1$, $j = 2$, en gelant les autres variables. Il s'agit de montrer que $\dfrac{\partial^2 f}{\partial y \partial x}(a,b) = \dfrac{\partial^2 f}{\partial x \partial y}(a,b)$.

  Pour $h, k$ petits, posons
  \[
    \Delta(h,k) = f(a+h, b+k) - f(a+h, b) - f(a, b+k) + f(a,b).
  \]
  Définissons $\varphi(t) = f(t, b+k) - f(t, b)$ ; alors $\Delta(h,k) = \varphi(a+h) - \varphi(a)$.

  Comme $\varphi$ est dérivable avec $\varphi'(t) = \dfrac{\partial f}{\partial x}(t, b+k) - \dfrac{\partial f}{\partial x}(t, b)$, le théorème des accroissements finis\index{théorème des accroissements finis} donne un $\theta_1 \in\, ]0,1[$ tel que
  \[
    \Delta(h,k) = h \left[\frac{\partial f}{\partial x}(a + \theta_1 h, b+k) - \frac{\partial f}{\partial x}(a + \theta_1 h, b)\right].
  \]
  En appliquant à nouveau le théorème des accroissements finis à $t \mapsto \dfrac{\partial f}{\partial x}(a + \theta_1 h, t)$, il existe $\theta_2 \in\, ]0,1[$ tel que
  \[
    \Delta(h,k) = hk \, \frac{\partial^2 f}{\partial y \partial x}(a + \theta_1 h, \, b + \theta_2 k).
  \]

  De manière symétrique, en posant $\psi(t) = f(a+h, t) - f(a, t)$ et en appliquant deux fois le théorème des accroissements finis, on obtient des $\theta_3, \theta_4 \in\, ]0,1[$ tels que
  \[
    \Delta(h,k) = hk \, \frac{\partial^2 f}{\partial x \partial y}(a + \theta_3 h, \, b + \theta_4 k).
  \]

  En divisant par $hk \neq 0$ et en faisant tendre $(h,k) \to (0,0)$ :
  \begin{align*}
    \frac{\partial^2 f}{\partial y \partial x}(a + \theta_1 h, \, b + \theta_2 k) &\xrightarrow[]{} \frac{\partial^2 f}{\partial y \partial x}(a,b), \\
    \frac{\partial^2 f}{\partial x \partial y}(a + \theta_3 h, \, b + \theta_4 k) &\xrightarrow[]{} \frac{\partial^2 f}{\partial x \partial y}(a,b),
  \end{align*}
  par continuité des dérivées secondes en $(a,b)$. L'unicité de la limite donne
  \[
    \frac{\partial^2 f}{\partial y \partial x}(a,b) = \frac{\partial^2 f}{\partial x \partial y}(a,b). \qedhere
  \]
\end{proof}

\begin{corollaire}
  \label{cor:Ck-symetrie}
  Si $f \in \mathcal{C}^k(\Omega, \R)$, alors les dérivées partielles d'ordre $\leqslant k$ sont indépendantes de l'ordre de dérivation.
\end{corollaire}

% ============================================================
\section{Formule de Taylor à plusieurs variables}
\label{sec:taylor-plusieurs-variables}
\index{formule de Taylor!plusieurs variables}

\subsection{Développement à l'ordre 1}

\begin{theoreme}[Taylor à l'ordre 1]
  \label{thm:taylor-ordre1}
  Soit $f \colon \Omega \to \R$ différentiable en $a \in \Omega$. Alors
  \[
    f(a + h) = f(a) + \sum_{i=1}^{n} \frac{\partial f}{\partial x_i}(a) \, h_i + o(\norm{h}).
  \]
\end{theoreme}

Ce n'est rien d'autre que la reformulation de la différentiabilité.

\subsection{Développement à l'ordre 2}

\begin{theoreme}[Taylor à l'ordre 2]
  \label{thm:taylor-ordre2}
  Soit $f \in \mathcal{C}^2(\Omega, \R)$ avec $a \in \Omega$. Alors
  \[
    f(a + h) = f(a) + \sum_{i=1}^{n} \frac{\partial f}{\partial x_i}(a) \, h_i
    + \frac{1}{2} \sum_{i=1}^{n} \sum_{j=1}^{n} \frac{\partial^2 f}{\partial x_i \partial x_j}(a) \, h_i \, h_j
    + o(\norm{h}^2).
  \]
\end{theoreme}

\begin{definition}[Matrice hessienne]\label{def:hessienne}
  Si $f \in \mathcal{C}^2(\Omega, \R)$, la \emph{matrice hessienne}\index{matrice hessienne} de~$f$ en~$a$ est la matrice symétrique (d'après le théorème de Schwarz, cf.\ théorème~\ref{thm:schwarz})
  \[
    H_f(a) = \left(\frac{\partial^2 f}{\partial x_i \partial x_j}(a)\right)_{1 \leqslant i,j \leqslant n} \in \mathcal{S}_n(\R).
  \]
  La formule de Taylor à l'ordre~2 se réécrit matriciellement :
  \[
    f(a+h) = f(a) + \langle \nabla f(a), h \rangle + \frac{1}{2} \langle H_f(a) \, h, h \rangle + o(\norm{h}^2).
  \]
\end{definition}

% ============================================================
\section{Extrema locaux}
\label{sec:extrema-locaux}
\index{extrema!locaux}

\subsection{Condition nécessaire du premier ordre}

\begin{theoreme}[Condition nécessaire : gradient nul]
  \label{thm:condition-necessaire-gradient}
  \index{point critique}
  Soient $\Omega \subset \R^n$ un ouvert et $f \colon \Omega \to \R$ différentiable. Si $f$ admet un extremum local en $a \in \Omega$, alors
  \[
    \nabla f(a) = 0.
  \]
  Un point vérifiant cette condition est appelé \emph{point critique}\index{point critique} ou \emph{point stationnaire} de~$f$.
\end{theoreme}

\begin{proof}
  Si $f$ admet un extremum local en~$a$, alors pour chaque $i$, la fonction d'une variable $t \mapsto f(a + t\,e_i)$ admet un extremum local en $t = 0$. Par le critère du premier ordre en une variable : $\dfrac{\partial f}{\partial x_i}(a) = 0$ pour tout~$i$, soit $\nabla f(a) = 0$.
\end{proof}

\subsection{Condition suffisante du second ordre}

\begin{theoreme}[Critère du second ordre]
  \label{thm:critere-second-ordre}
  \index{critère du second ordre}
  Soit $f \in \mathcal{C}^2(\Omega, \R)$ avec $\nabla f(a) = 0$ (point critique). Notons $H = H_f(a)$ la matrice hessienne.
  \begin{enumerate}[label=(\roman*)]
    \item Si $H$ est \emph{définie positive}\index{définie positive}, alors $f$ admet un \emph{minimum local strict} en~$a$.
    \item Si $H$ est \emph{définie négative}\index{définie négative}, alors $f$ admet un \emph{maximum local strict} en~$a$.
    \item Si $H$ est \emph{indéfinie} (c'est-à-dire possède des valeurs propres de signes opposés), alors $a$ est un \emph{point selle}\index{point selle}.
    \item Si $H$ est semi-définie (positive ou négative), on ne peut pas conclure sans informations supplémentaires.
  \end{enumerate}
\end{theoreme}

\begin{remarque}
  \label{rem:critere-dim2}
  Dans le cas $n = 2$, la matrice hessienne en un point critique $a$ est
  \[
    H = \begin{pmatrix}
      r & s \\ s & t
    \end{pmatrix},
    \quad \text{où } r = \frac{\partial^2 f}{\partial x^2}(a),\;
    s = \frac{\partial^2 f}{\partial x \partial y}(a),\;
    t = \frac{\partial^2 f}{\partial y^2}(a).
  \]
  Posons $\Delta = rt - s^2 = \det H$.
  \begin{itemize}
    \item Si $\Delta > 0$ et $r > 0$ : minimum local strict.
    \item Si $\Delta > 0$ et $r < 0$ : maximum local strict.
    \item Si $\Delta < 0$ : point selle.
    \item Si $\Delta = 0$ : cas douteux.
  \end{itemize}
\end{remarque}

\begin{exemple}
  \label{ex:extrema}
  \begin{enumerate}[label=(\alph*)]
    \item $f(x,y) = x^2 + y^2$. Point critique : $(0,0)$. Hessienne : $H = \begin{psmallmatrix} 2 & 0 \\ 0 & 2 \end{psmallmatrix}$, $\Delta = 4 > 0$, $r = 2 > 0$ : \emph{minimum local} (et global).

    \item $f(x,y) = -x^2 - y^2$. Point critique : $(0,0)$. $H = \begin{psmallmatrix} -2 & 0 \\ 0 & -2 \end{psmallmatrix}$, $\Delta = 4 > 0$, $r = -2 < 0$ : \emph{maximum local} (et global).

    \item $f(x,y) = x^2 - y^2$. Point critique : $(0,0)$. $H = \begin{psmallmatrix} 2 & 0 \\ 0 & -2 \end{psmallmatrix}$, $\Delta = -4 < 0$ : \emph{point selle}.
  \end{enumerate}
\end{exemple}

\begin{center}
\begin{tikzpicture}
  \begin{axis}[
    view={35}{30},
    xlabel={$x$}, ylabel={$y$}, zlabel={$z$},
    xmin=-1.5, xmax=1.5,
    ymin=-1.5, ymax=1.5,
    zmin=-1.5, zmax=1.5,
    width=7cm, height=6cm,
    title={\small Minimum : $x^2+y^2$},
    samples=15,
    colormap/cool,
    mesh/ordering=y varies,
  ]
    \addplot3[surf, shader=flat, opacity=0.8, domain=-1.2:1.2, domain y=-1.2:1.2]
      {x^2 + y^2};
  \end{axis}
\end{tikzpicture}
\hfill
\begin{tikzpicture}
  \begin{axis}[
    view={35}{30},
    xlabel={$x$}, ylabel={$y$}, zlabel={$z$},
    xmin=-1.5, xmax=1.5,
    ymin=-1.5, ymax=1.5,
    zmin=-1.5, zmax=1.5,
    width=7cm, height=6cm,
    title={\small Maximum : $-x^2-y^2$},
    samples=15,
    colormap/cool,
    mesh/ordering=y varies,
  ]
    \addplot3[surf, shader=flat, opacity=0.8, domain=-1.2:1.2, domain y=-1.2:1.2]
      {-x^2 - y^2};
  \end{axis}
\end{tikzpicture}
\hfill
\begin{tikzpicture}
  \begin{axis}[
    view={35}{30},
    xlabel={$x$}, ylabel={$y$}, zlabel={$z$},
    xmin=-1.5, xmax=1.5,
    ymin=-1.5, ymax=1.5,
    zmin=-1.5, zmax=1.5,
    width=7cm, height=6cm,
    title={\small Point selle : $x^2-y^2$},
    samples=15,
    colormap/cool,
    mesh/ordering=y varies,
  ]
    \addplot3[surf, shader=flat, opacity=0.8, domain=-1.2:1.2, domain y=-1.2:1.2]
      {x^2 - y^2};
  \end{axis}
\end{tikzpicture}
\captionof{figure}{De gauche à droite : minimum local, maximum local, point selle.}
\label{fig:extrema-types}
\end{center}

% ============================================================
\section{Exercices}
\label{sec:exercices-chap9}

\begin{exercice}[Normes et topologie]
  \label{exo:normes}
  Soit $x = (3, -4, 1) \in \R^3$. Calculer $\norm{x}_1$, $\norm{x}_2$ et $\norm{x}_\infty$. Vérifier que $\norm{x}_\infty \leqslant \norm{x}_2 \leqslant \norm{x}_1 \leqslant 3\,\norm{x}_\infty$.
\end{exercice}

\begin{exercice}[Ouverts et fermés]
  \label{exo:ouverts-fermes}
  Déterminer si les ensembles suivants sont ouverts, fermés, les deux ou ni l'un ni l'autre :
  \begin{enumerate}[label=(\alph*)]
    \item $A = \{(x,y) \in \R^2 : x^2 + y^2 < 1\}$ ;
    \item $B = \{(x,y) \in \R^2 : x^2 + y^2 \leqslant 1\}$ ;
    \item $C = \{(x,y) \in \R^2 : 0 < x^2 + y^2 \leqslant 1\}$ ;
    \item $D = \{(x,y) \in \R^2 : x \in \Q\}$.
  \end{enumerate}
\end{exercice}

\begin{exercice}[Limites en plusieurs variables]
  \label{exo:limites}
  Déterminer si les limites suivantes existent et, le cas échéant, les calculer :
  \begin{enumerate}[label=(\alph*)]
    \item $\displaystyle \lim_{(x,y) \to (0,0)} \frac{x^2 y}{x^2 + y^2}$ ;
    \item $\displaystyle \lim_{(x,y) \to (0,0)} \frac{x^2 - y^2}{x^2 + y^2}$ ;
    \item $\displaystyle \lim_{(x,y) \to (0,0)} \frac{x^3 + y^3}{x^2 + y^2}$.
  \end{enumerate}
  \emph{Indication :} Pour (a) et (c), passer en coordonnées polaires. Pour (b), essayer $y = 0$ puis $x = 0$.
\end{exercice}

\begin{exercice}[Continuité et dérivées partielles]
  \label{exo:cont-partielles}
  Soit $f \colon \R^2 \to \R$ définie par
  \[
    f(x,y) = \begin{cases} \dfrac{x^3 y}{x^4 + y^2} & \text{si } (x,y) \neq (0,0), \\ 0 & \text{si } (x,y) = (0,0). \end{cases}
  \]
  Montrer que les dérivées partielles de~$f$ existent en $(0,0)$, mais que $f$ n'est pas continue en $(0,0)$.
\end{exercice}

\begin{exercice}[Calcul de différentielle]
  \label{exo:calcul-differentielle}
  Soit $f(x,y) = x^2 y + e^{xy}$. Calculer $\dd f_{(1,1)}$ et écrire l'approximation affine de~$f$ au voisinage de $(1,1)$.
\end{exercice}

\begin{exercice}[Matrice jacobienne]
  \label{exo:jacobienne}
  Soit $F \colon \R^2 \to \R^2$ définie par $F(r, \theta) = (r \cos\theta, \, r \sin\theta)$ (coordonnées polaires). Calculer la matrice jacobienne $J_F(r, \theta)$ et son déterminant. Interpréter géométriquement.
\end{exercice}

\begin{exercice}[Règle de la chaîne]
  \label{exo:chain-rule}
  Soit $f \colon \R^2 \to \R$ de classe $\mathcal{C}^1$, et soient $x = r\cos\theta$, $y = r\sin\theta$. Exprimer $\dfrac{\partial f}{\partial r}$ et $\dfrac{\partial f}{\partial \theta}$ en fonction de $\dfrac{\partial f}{\partial x}$ et $\dfrac{\partial f}{\partial y}$.
\end{exercice}

\begin{exercice}[Gradient et lignes de niveau]
  \label{exo:gradient-lignes-niveau}
  Soit $f(x,y) = x^2 - 2xy + 3y^2$.
  \begin{enumerate}[label=(\alph*)]
    \item Calculer $\nabla f(x,y)$.
    \item Déterminer l'équation de la ligne de niveau passant par le point $(1,1)$.
    \item Vérifier que $\nabla f(1,1)$ est orthogonal à cette ligne de niveau en ce point.
  \end{enumerate}
\end{exercice}

\begin{exercice}[Théorème de Schwarz]
  \label{exo:schwarz}
  Soit $f(x,y) = \arctan\!\left(\dfrac{y}{x}\right)$ pour $x > 0$. Calculer $\dfrac{\partial^2 f}{\partial x \partial y}$ et $\dfrac{\partial^2 f}{\partial y \partial x}$ et vérifier l'égalité du théorème de Schwarz.
\end{exercice}

\begin{exercice}[Recherche d'extrema]
  \label{exo:extrema-1}
  Trouver et classifier les points critiques de $f(x,y) = x^3 + y^3 - 3xy$.
\end{exercice}

\begin{exercice}[Extrema sur un compact]
  \label{exo:extrema-compact}
  Déterminer le maximum et le minimum de $f(x,y) = 2x + 3y$ sur le disque fermé $\overline{B}(0,1) = \{(x,y) : x^2 + y^2 \leqslant 1\}$.
  \emph{Indication :} Le maximum et le minimum sont atteints sur le bord (pourquoi ?). Paramétrer le bord.
\end{exercice}

\begin{exercice}[Taylor à l'ordre 2]
  \label{exo:taylor-2}
  Soit $f(x,y) = e^x \sin y$. Écrire le développement de Taylor de~$f$ à l'ordre~2 au voisinage de $(0, 0)$.
\end{exercice}

\begin{exercice}[Fonction non différentiable avec partielles]
  \label{exo:non-diff-partielles}
  Soit $f \colon \R^2 \to \R$ définie par
  \[
    f(x,y) = \begin{cases} \dfrac{xy}{\sqrt{x^2+y^2}} & \text{si } (x,y) \neq (0,0), \\ 0 & \text{sinon.} \end{cases}
  \]
  \begin{enumerate}[label=(\alph*)]
    \item Montrer que $f$ est continue en $(0,0)$.
    \item Calculer les dérivées partielles de~$f$ en $(0,0)$.
    \item Montrer que $f$ n'est pas différentiable en $(0,0)$.
  \end{enumerate}
\end{exercice}

% ============================================================
\section*{Résumé du chapitre}
\addcontentsline{toc}{section}{Résumé du chapitre}

\begin{tcolorbox}[
  enhanced,
  colback=gray!5!white,
  colframe=gray!60!black,
  fonttitle=\bfseries\large,
  title={Résumé --- Chapitre \thechapter},
  breakable,
  before skip=12pt,
  after skip=12pt,
  left=10pt, right=10pt, top=8pt, bottom=8pt
]
\begin{enumerate}[label=\textbf{\arabic*.}, leftmargin=2em]
  \item \textbf{Topologie de $\R^n$.} Toutes les normes sur $\R^n$ sont équivalentes (théorème~\ref{thm:equivalence-normes}). Les notions d'ouverts, fermés, compacts, convergence de suites sont indépendantes de la norme choisie.

  \item \textbf{Limites et continuité.} La limite double est plus forte que les limites itérées : celles-ci peuvent exister sans que la limite double n'existe (exemple~\ref{ex:limites-iterees}). Une fonction continue sur un compact atteint ses bornes (théorème~\ref{thm:bornes-atteintes-compact}).

  \item \textbf{Dérivées partielles.} L'existence de toutes les dérivées partielles en un point n'implique ni la continuité ni la différentiabilité (exemple~\ref{ex:partials-not-continuity}).

  \item \textbf{Différentiabilité.} Une fonction est différentiable en~$a$ si elle admet une approximation affine de la forme $f(a+h) = f(a) + L(h) + o(\norm{h})$, où $L$ est linéaire. La différentiabilité implique la continuité (théorème~\ref{thm:diff-implique-cont}) et l'existence des dérivées partielles (théorème~\ref{thm:diff-implique-partielles}).

  \item \textbf{Condition $\mathcal{C}^1$.} Si les dérivées partielles existent et sont continues, alors $f$ est différentiable (théorème~\ref{thm:C1-diff}). C'est le critère le plus utilisé en pratique.

  \item \textbf{Règle de la chaîne.} $\dd(g \circ f)_a = \dd g_{f(a)} \circ \dd f_a$ (théorème~\ref{thm:chain-rule}). En termes de jacobiennes : $J_{g \circ f}(a) = J_g(f(a)) \cdot J_f(a)$.

  \item \textbf{Gradient.} $\nabla f(a)$ pointe dans la direction de plus grande croissance et est orthogonal aux lignes de niveau.

  \item \textbf{Théorème de Schwarz.} Sous hypothèse de continuité des dérivées secondes, l'ordre de dérivation est indifférent (théorème~\ref{thm:schwarz}).

  \item \textbf{Taylor.} La formule de Taylor à l'ordre~2 fait intervenir la matrice hessienne $H_f(a)$ (théorème~\ref{thm:taylor-ordre2}).

  \item \textbf{Extrema.} Condition nécessaire : $\nabla f(a) = 0$. La nature du point critique est déterminée par le signe de la hessienne : définie positive (minimum), définie négative (maximum), indéfinie (point selle) (théorème~\ref{thm:critere-second-ordre}).
\end{enumerate}
\end{tcolorbox}% ============================================================
%  Chapitre 10 : Intégrales Multiples et Changement de Variables
% ============================================================

\chapter{Intégrales Multiples et Changement de Variables}
\label{chap:integrales-multiples}

\index{intégrale multiple}
\index{changement de variables}

Le calcul intégral d'une variable, développé au chapitre~7, permet de
mesurer des longueurs, des aires sous une courbe, des travaux le long
d'un axe. Mais dès que l'on souhaite calculer le \emph{volume} d'un
solide, la \emph{masse} d'une plaque de densité variable, ou le
\emph{centre de masse} d'un objet plan, il faut intégrer des fonctions
de \emph{plusieurs} variables. Ce chapitre étend la théorie de l'intégrale
de Riemann au cadre multidimensionnel et introduit l'outil fondamental
qu'est le changement de variables, dont les coordonnées polaires,
cylindriques et sphériques sont les applications les plus classiques.

% ============================================================
%  10.1 Intégrale double sur un rectangle
% ============================================================

\section{Intégrale double sur un rectangle}
\label{sec:integrale-double-rectangle}

\index{intégrale double}
\index{somme de Riemann!en dimension 2}

Soit $R = [a,b]\times[c,d]$ un rectangle fermé de $\R^2$. Une
\emph{partition}\index{partition!d'un rectangle} de $R$ est la donnée
de partitions $a = x_0 < x_1 < \cdots < x_m = b$ et
$c = y_0 < y_1 < \cdots < y_n = d$ des deux intervalles. Les
sous-rectangles sont
$R_{ij} = [x_{i-1}, x_i]\times[y_{j-1}, y_j]$, d'aire
$\Delta A_{ij} = (x_i - x_{i-1})(y_j - y_{j-1})$.

\begin{definition}[Somme de Riemann double]\label{def:somme-riemann-2d}
  Soit $f\colon R \to \R$ une fonction bornée et $P$ une partition de $R$.
  Pour tout choix de points $(\xi_{ij}, \eta_{ij}) \in R_{ij}$, la
  \emph{somme de Riemann}\index{somme de Riemann!double} associée est
  \[
    S(f, P) = \sum_{i=1}^{m}\sum_{j=1}^{n} f(\xi_{ij}, \eta_{ij})\,
    \Delta A_{ij}.
  \]
\end{definition}

\begin{definition}{Intégrabilité au sens de Riemann sur un rectangle}%
  {integrabilite-rectangle}
  La fonction $f\colon R \to \R$ bornée est dite \emph{intégrable au
  sens de Riemann} sur $R$ si les sommes de Darboux supérieure et
  inférieure coïncident, c'est-à-dire si
  \[
    \inf_P U(f,P) = \sup_P L(f,P),
  \]
  où $U(f,P) = \sum_{i,j}\sup_{R_{ij}}f\;\Delta A_{ij}$ et
  $L(f,P) = \sum_{i,j}\inf_{R_{ij}}f\;\Delta A_{ij}$. La valeur
  commune est notée
  \[
    \iint_R f(x,y)\,\dd x\,\dd y
    \quad\text{ou}\quad
    \iint_R f\,\dd A.
  \]
\end{definition}

\begin{theoreme}[Intégrabilité des fonctions continues]
  \label{thm:continue-integrable-2d}
  \index{intégrabilité!des fonctions continues}
  Toute fonction continue sur un rectangle fermé $R = [a,b]\times[c,d]$
  est intégrable au sens de Riemann sur $R$.
\end{theoreme}

\begin{proof}
  Puisque $R$ est compact et $f$ est continue sur $R$, la fonction $f$
  est uniformément continue sur $R$ (théorème de Heine). Soit $\eps > 0$.
  Il existe $\delta > 0$ tel que $\norm{(x,y)-(x',y')} < \delta$
  implique $\abs{f(x,y)-f(x',y')} < \frac{\eps}{\abs{R}}$, où
  $\abs{R} = (b-a)(d-c)$. Pour toute partition $P$ de pas
  $\norm{P} < \delta/\sqrt{2}$, le diamètre de chaque sous-rectangle
  $R_{ij}$ est inférieur à $\delta$. Par conséquent
  \[
    \sup_{R_{ij}} f - \inf_{R_{ij}} f < \frac{\eps}{\abs{R}},
  \]
  et donc
  \[
    U(f,P) - L(f,P) = \sum_{i,j}\bigl(\sup_{R_{ij}}f -
    \inf_{R_{ij}}f\bigr)\Delta A_{ij}
    < \frac{\eps}{\abs{R}} \sum_{i,j} \Delta A_{ij} = \eps.
  \]
  Le critère de Riemann assure l'intégrabilité de $f$.
\end{proof}


% ============================================================
%  10.2 Théorème de Fubini
% ============================================================

\section{Théorème de Fubini}
\label{sec:fubini}

\index{Fubini (théorème de)}

Le théorème de Fubini ramène le calcul d'une intégrale double à deux
intégrales itérées simples.

\begin{theoreme}[Fubini]\label{thm:fubini}
  Soit $f\colon [a,b]\times[c,d] \to \R$ une fonction continue. Alors
  \begin{equation}\label{eq:fubini}
    \iint_R f(x,y)\,\dd x\,\dd y
    = \int_a^b \!\left(\int_c^d f(x,y)\,\dd y\right)\dd x
    = \int_c^d \!\left(\int_a^b f(x,y)\,\dd x\right)\dd y.
  \end{equation}
\end{theoreme}

\begin{proof}
  Posons $\varphi(x) = \int_c^d f(x,y)\,\dd y$. Montrons d'abord que
  $\varphi$ est continue. Soit $x_0 \in [a,b]$. Pour tout $\eps > 0$,
  l'uniforme continuité de $f$ sur $R$ fournit $\delta > 0$ tel que
  $\abs{x - x_0} < \delta$ implique $\abs{f(x,y) - f(x_0,y)} <
  \frac{\eps}{d-c}$ pour tout $y \in [c,d]$. Alors
  \[
    \abs{\varphi(x) - \varphi(x_0)}
    \leqslant \int_c^d \abs{f(x,y) - f(x_0,y)}\,\dd y
    < \frac{\eps}{d-c}\cdot(d-c) = \eps.
  \]

  Montrons maintenant l'égalité~\eqref{eq:fubini} (première forme).
  Considérons une partition $P$ de $R$ en sous-rectangles
  $R_{ij} = [x_{i-1},x_i]\times[y_{j-1},y_j]$. Puisque $f$ est
  continue et $R_{ij}$ est compact, $f$ atteint ses bornes sur
  $R_{ij}$. On a, pour chaque $i$ fixé,
  \[
    \sum_{j=1}^n \inf_{R_{ij}} f \cdot (y_j - y_{j-1})
    \leqslant \int_c^d f(x_i^*, y)\,\dd y
    \leqslant \sum_{j=1}^n \sup_{R_{ij}} f \cdot (y_j - y_{j-1})
  \]
  pour tout $x_i^* \in [x_{i-1}, x_i]$. En sommant sur $i$ et en
  multipliant par $(x_i - x_{i-1})$, on encadre
  $\int_a^b \varphi(x)\,\dd x$ entre $L(f,P)$ et $U(f,P)$. Puisque
  $f$ est intégrable sur $R$, on obtient
  \[
    \int_a^b\!\left(\int_c^d f(x,y)\,\dd y\right)\dd x
    = \iint_R f(x,y)\,\dd x\,\dd y.
  \]
  La seconde égalité s'obtient par symétrie des rôles de $x$ et $y$.
\end{proof}

\begin{exemple}\label{ex:fubini-calcul}
  Calculons $\displaystyle I = \iint_R xy^2\,\dd x\,\dd y$ avec
  $R = [0,1]\times[0,2]$.
  \[
    I = \int_0^1\!\left(\int_0^2 xy^2\,\dd y\right)\dd x
    = \int_0^1 x\left[\frac{y^3}{3}\right]_0^2 \dd x
    = \int_0^1 \frac{8x}{3}\,\dd x
    = \frac{8}{3}\cdot\frac{1}{2} = \frac{4}{3}.
  \]
\end{exemple}


% ============================================================
%  10.3 Intégrale sur un domaine borné
% ============================================================

\section{Intégrale double sur un domaine borné}
\label{sec:integrale-domaine-borne}

\index{domaine!de type~I}
\index{domaine!de type~II}

En pratique, on intègre rarement sur un rectangle ; les domaines sont
délimités par des courbes.

\begin{definition}[Domaines simples]\label{def:domaines-simples}
  Un domaine $D \subset \R^2$ est dit de \emph{type~I} s'il est de la
  forme
  \[
    D = \{(x,y) \in \R^2 : a \leqslant x \leqslant b,\;
    g_1(x) \leqslant y \leqslant g_2(x)\}
  \]
  où $g_1, g_2\colon [a,b]\to\R$ sont continues avec $g_1 \leqslant g_2$.
  Le domaine $D$ est de \emph{type~II} s'il est de la forme
  \[
    D = \{(x,y) \in \R^2 : c \leqslant y \leqslant d,\;
    h_1(y) \leqslant x \leqslant h_2(y)\}
  \]
  où $h_1, h_2\colon [c,d]\to\R$ sont continues avec $h_1 \leqslant h_2$.
\end{definition}

\begin{center}
\begin{tikzpicture}[scale=0.9]
  % Domaine de type I
  \begin{scope}[xshift=-4.5cm]
    \draw[->] (-0.3,0) -- (3.5,0) node[right] {$x$};
    \draw[->] (0,-0.3) -- (0,3.5) node[above] {$y$};
    \draw[dashed] (0.5,0) node[below] {$a$} -- (0.5,0.4);
    \draw[dashed] (3,0) node[below] {$b$} -- (3,2.5);
    \draw[blue,thick,domain=0.5:3,smooth] plot (\x, {0.2+0.1*(\x-0.5)^2})
      node[right] {$g_1(x)$};
    \draw[red,thick,domain=0.5:3,smooth] plot (\x, {2.8-0.15*(\x-1.5)^2})
      node[right] {$g_2(x)$};
    \fill[blue!10,opacity=0.5]
      plot[domain=0.5:3,smooth] (\x,{0.2+0.1*(\x-0.5)^2}) --
      plot[domain=3:0.5,smooth] (\x,{2.8-0.15*(\x-1.5)^2}) -- cycle;
    \node at (1.8,1.5) {$D$};
    \node at (1.8,-0.8) {Type~I};
  \end{scope}
  % Domaine de type II
  \begin{scope}[xshift=4cm]
    \draw[->] (-0.3,0) -- (3.5,0) node[right] {$x$};
    \draw[->] (0,-0.3) -- (0,3.5) node[above] {$y$};
    \draw[dashed] (0,0.5) node[left] {$c$} -- (0.3,0.5);
    \draw[dashed] (0,3) node[left] {$d$} -- (0.5,3);
    \draw[blue,thick,domain=0.5:3,smooth]
      plot ({0.2+0.1*(\x-0.5)^2}, \x) node[above] {$h_1(y)$};
    \draw[red,thick,domain=0.5:3,smooth]
      plot ({2.8-0.15*(\x-1.5)^2}, \x) node[above] {$h_2(y)$};
    \fill[red!10,opacity=0.5]
      plot[domain=0.5:3,smooth] ({0.2+0.1*(\x-0.5)^2},\x) --
      plot[domain=3:0.5,smooth] ({2.8-0.15*(\x-1.5)^2},\x) -- cycle;
    \node at (1.5,1.8) {$D$};
    \node at (1.5,-0.8) {Type~II};
  \end{scope}
\end{tikzpicture}
\end{center}

\begin{proposition}\label{prop:integrale-domaine-type-I}
  Si $D$ est un domaine de type~I et $f\colon D \to \R$ est continue,
  alors
  \begin{equation}\label{eq:type-I}
    \iint_D f(x,y)\,\dd x\,\dd y
    = \int_a^b \!\left(\int_{g_1(x)}^{g_2(x)} f(x,y)\,\dd y\right)\dd x.
  \end{equation}
  Symétriquement, si $D$ est de type~II,
  \begin{equation}\label{eq:type-II}
    \iint_D f(x,y)\,\dd x\,\dd y
    = \int_c^d \!\left(\int_{h_1(y)}^{h_2(y)} f(x,y)\,\dd x\right)\dd y.
  \end{equation}
\end{proposition}

\begin{exemple}[Intégrale sur un triangle]
  \label{ex:integrale-triangle}
  Soit $D$ le triangle de sommets $(0,0)$, $(1,0)$, $(0,1)$, c'est-à-dire
  $D = \{(x,y) : 0 \leqslant x \leqslant 1,\; 0 \leqslant y \leqslant 1-x\}$.
  Calculons $\displaystyle I = \iint_D (x+y)\,\dd x\,\dd y$.
  \[
    I = \int_0^1\!\left(\int_0^{1-x}(x+y)\,\dd y\right)\dd x
    = \int_0^1\left[xy + \frac{y^2}{2}\right]_0^{1-x}\dd x
    = \int_0^1\left(x(1-x) + \frac{(1-x)^2}{2}\right)\dd x.
  \]
  Développons :
  \[
    x(1-x) + \frac{(1-x)^2}{2} = x - x^2 + \frac{1-2x+x^2}{2}
    = \frac{1}{2} - \frac{x^2}{2}.
  \]
  Donc $I = \int_0^1\left(\frac{1}{2} - \frac{x^2}{2}\right)\dd x
  = \frac{1}{2} - \frac{1}{6} = \frac{1}{3}$.
\end{exemple}

\begin{exemple}[Intégrale sur un disque]
  \label{ex:integrale-disque}
  Soit $D = \{(x,y) : x^2 + y^2 \leqslant 1\}$. En tant que domaine
  de type~I, $D = \{(x,y) : -1 \leqslant x \leqslant 1,\;
  -\sqrt{1-x^2} \leqslant y \leqslant \sqrt{1-x^2}\}$. Alors
  \[
    \iint_D 1\,\dd x\,\dd y
    = \int_{-1}^{1}\!\int_{-\sqrt{1-x^2}}^{\sqrt{1-x^2}}\dd y\,\dd x
    = \int_{-1}^{1} 2\sqrt{1-x^2}\,\dd x = \pi,
  \]
  cette dernière intégrale se calculant par le changement de variable
  $x = \sin\theta$. Le résultat attendu $\pi r^2$ (ici $r=1$) est
  retrouvé. Nous verrons en section~\ref{sec:changement-variables} un
  calcul plus élégant en coordonnées polaires.
\end{exemple}


% ============================================================
%  10.4 Changement de variables
% ============================================================

\section{Changement de variables dans les intégrales multiples}
\label{sec:changement-variables}

\index{changement de variables!intégrales multiples}
\index{jacobien}

\begin{theoreme}[Changement de variables]
  \label{thm:changement-variables}
  Soit $U$ un ouvert de $\R^n$, et $\Phi\colon U \to \R^n$ un
  difféomorphisme de classe $\mathcal{C}^1$ d'image $V = \Phi(U)$. Pour
  toute fonction $f\colon V \to \R$ continue à support compact (ou
  intégrable),
  \begin{equation}\label{eq:changement-variables}
    \int_V f(\mathbf{y})\,\dd\mathbf{y}
    = \int_U f\bigl(\Phi(\mathbf{u})\bigr)\,
    \abs{\det J_\Phi(\mathbf{u})}\,\dd\mathbf{u},
  \end{equation}
  où $J_\Phi(\mathbf{u})$ est la matrice jacobienne de $\Phi$ en
  $\mathbf{u}$.
\end{theoreme}

\begin{remarque}
  Le facteur $\abs{\det J_\Phi}$ s'appelle le \emph{jacobien}\index{jacobien}
  du changement de variables. Il mesure le rapport local des volumes :
  un élément infinitésimal de volume $\dd\mathbf{u}$ est transformé en
  un élément de volume $\abs{\det J_\Phi}\,\dd\mathbf{u}$.
\end{remarque}


\subsection{Coordonnées polaires}
\label{subsec:coord-polaires}

\index{coordonnées polaires}

Le changement de variables le plus fondamental en dimension~2 est le
passage en coordonnées polaires :
\[
  \Phi(r,\theta) = (r\cos\theta,\; r\sin\theta),
  \qquad r > 0,\quad \theta \in\; ]0, 2\pi[.
\]
La matrice jacobienne est
\[
  J_\Phi(r,\theta) = \begin{pmatrix}
    \dfrac{\partial x}{\partial r} & \dfrac{\partial x}{\partial \theta} \\[8pt]
    \dfrac{\partial y}{\partial r} & \dfrac{\partial y}{\partial \theta}
  \end{pmatrix}
  = \begin{pmatrix}
    \cos\theta & -r\sin\theta \\
    \sin\theta & r\cos\theta
  \end{pmatrix}.
\]
Le jacobien vaut
\begin{equation}\label{eq:jacobien-polaire}
  \det J_\Phi(r,\theta) = r\cos^2\theta + r\sin^2\theta = r.
\end{equation}
On a donc $\dd x\,\dd y = r\,\dd r\,\dd\theta$.

\begin{center}
\begin{tikzpicture}[scale=1.1]
  % Grille polaire
  \draw[->] (-2.5,0) -- (2.8,0) node[right] {$x$};
  \draw[->] (0,-2.5) -- (0,2.8) node[above] {$y$};
  \foreach \r in {0.5,1,1.5,2} {
    \draw[blue!40,thin] (0,0) circle (\r);
  }
  \foreach \a in {0,30,...,330} {
    \draw[red!40,thin] (0,0) -- (\a:2.2);
  }
  \node[below right] at (1.5,0) {$r$};
  \draw[->,thick,green!50!black] (0.6,0) arc (0:45:0.6);
  \node[green!50!black] at (0.85,0.35) {$\theta$};
  \fill[blue!20,opacity=0.5] (0,0) -- (30:1.5) arc (30:60:1.5) --
    cycle;
  \fill[blue!20,opacity=0.5] (0,0) -- (30:1) arc (30:60:1) -- cycle;
  \draw[thick,blue] (30:1) arc (30:60:1);
  \draw[thick,blue] (30:1.5) arc (30:60:1.5);
  \draw[thick,blue] (30:1) -- (30:1.5);
  \draw[thick,blue] (60:1) -- (60:1.5);
  \node at (45:1.25) {\small$\dd A$};
\end{tikzpicture}
\end{center}

\begin{exemple}[Aire du disque]\label{ex:aire-disque-polaire}
  Soit $D = \{(x,y) : x^2 + y^2 \leqslant R^2\}$. En coordonnées
  polaires,
  \[
    \text{Aire}(D) = \iint_D \dd x\,\dd y
    = \int_0^{2\pi}\!\int_0^R r\,\dd r\,\dd\theta
    = \int_0^{2\pi}\frac{R^2}{2}\,\dd\theta = \pi R^2.
  \]
\end{exemple}

\begin{theoreme}[Intégrale de Gauss]
  \label{thm:integrale-gauss}
  \index{intégrale de Gauss}
  \index{Gauss (intégrale de)}
  On a
  \begin{equation}\label{eq:gauss}
    \int_{-\infty}^{+\infty} e^{-x^2}\,\dd x = \sqrt{\pi}.
  \end{equation}
\end{theoreme}

\begin{proof}
  Posons $I = \int_{-\infty}^{+\infty} e^{-x^2}\,\dd x$. Puisque
  l'intégrande est positif et pair, $I > 0$ ; montrons que $I^2 = \pi$.
  Par le théorème de Fubini\index{Fubini (théorème de)} (étendu aux
  intégrales impropres, ce qui est licite ici car l'intégrande est
  positif),
  \[
    I^2 = \left(\int_{-\infty}^{+\infty} e^{-x^2}\,\dd x\right)
    \left(\int_{-\infty}^{+\infty} e^{-y^2}\,\dd y\right)
    = \iint_{\R^2} e^{-(x^2+y^2)}\,\dd x\,\dd y.
  \]
  Effectuons le passage en coordonnées polaires
  ($x = r\cos\theta$, $y = r\sin\theta$, $\dd x\,\dd y = r\,\dd r\,\dd\theta$) :
  \[
    I^2 = \int_0^{2\pi}\!\int_0^{+\infty} e^{-r^2}\,r\,\dd r\,\dd\theta
    = 2\pi \int_0^{+\infty} r\,e^{-r^2}\,\dd r.
  \]
  On calcule l'intégrale radiale par le changement de variable
  $u = r^2$, $\dd u = 2r\,\dd r$ :
  \[
    \int_0^{+\infty} r\,e^{-r^2}\,\dd r
    = \frac{1}{2}\int_0^{+\infty} e^{-u}\,\dd u
    = \frac{1}{2}\bigl[-e^{-u}\bigr]_0^{+\infty}
    = \frac{1}{2}.
  \]
  Donc $I^2 = 2\pi \cdot \frac{1}{2} = \pi$, et comme $I > 0$, on
  conclut $I = \sqrt{\pi}$.
\end{proof}

\begin{remarque}\label{rem:gauss-probabilites}
  L'intégrale de Gauss est fondamentale en théorie des probabilités : la
  densité de la loi normale $\mathcal{N}(0,1)$ est
  $\frac{1}{\sqrt{2\pi}}e^{-x^2/2}$, dont l'intégrale vaut~$1$
  précisément grâce à~\eqref{eq:gauss}.
\end{remarque}


% ============================================================
%  10.5 Intégrales triples
% ============================================================

\section{Intégrales triples}
\label{sec:integrales-triples}

\index{intégrale triple}

Le théorème de Fubini\index{Fubini (théorème de)} s'étend naturellement
à la dimension~3 : si $f$ est continue sur un domaine $E \subset \R^3$
borné et « simple », l'intégrale triple
$\iiint_E f(x,y,z)\,\dd x\,\dd y\,\dd z$ se ramène à des intégrales
itérées.

\subsection{Coordonnées cylindriques}
\label{subsec:coord-cylindriques}

\index{coordonnées cylindriques}

On pose $(x,y,z) = (r\cos\theta,\; r\sin\theta,\; z)$ avec $r \geqslant 0$,
$\theta \in [0, 2\pi[$. Le jacobien est
\[
  \det J = \det\begin{pmatrix}
    \cos\theta & -r\sin\theta & 0 \\
    \sin\theta & r\cos\theta  & 0 \\
    0          & 0            & 1
  \end{pmatrix} = r.
\]
Donc $\dd x\,\dd y\,\dd z = r\,\dd r\,\dd\theta\,\dd z$.

\subsection{Coordonnées sphériques}
\label{subsec:coord-spheriques}

\index{coordonnées sphériques}

On pose
\[
  \begin{cases}
    x = \rho\sin\varphi\cos\theta, \\
    y = \rho\sin\varphi\sin\theta, \\
    z = \rho\cos\varphi,
  \end{cases}
  \quad \rho \geqslant 0,\quad \varphi \in [0,\pi],\quad
  \theta \in [0, 2\pi[.
\]

\begin{center}
\begin{tikzpicture}[scale=1.4,tdplot_main_coords/.style={},
  every node/.style={font=\small}]
  % Axes
  \draw[->] (0,0,0) -- (2.5,0,0) node[right] {$x$};
  \draw[->] (0,0,0) -- (0,2.5,0) node[above] {$y$};
  \draw[->] (0,0,0) -- (0,0,2.5) node[above] {$z$};
  % Point
  \coordinate (P) at (1.3,1.1,1.5);
  \draw[thick,->,blue] (0,0,0) -- (P) node[above right] {$(\rho,\varphi,\theta)$};
  % Projections
  \draw[dashed] (P) -- (1.3,1.1,0) node[below] {$(x,y,0)$};
  \draw[dashed] (1.3,1.1,0) -- (0,0,0);
  % Angles
  \draw[->,red,thick] (0,0,0.7) arc[start angle=90,end angle=55,radius=0.7];
  \node[red] at (0.25,0.1,0.85) {$\varphi$};
  \draw[->,green!50!black,thick] (0.5,0,0) arc[start angle=0,end angle=40,radius=0.5];
  \node[green!50!black] at (0.55,0.2,0) {$\theta$};
  % Rho label
  \node[blue,left] at (0.6,0.5,0.8) {$\rho$};
\end{tikzpicture}
\end{center}

La matrice jacobienne est
\[
  J = \begin{pmatrix}
    \sin\varphi\cos\theta & \rho\cos\varphi\cos\theta & -\rho\sin\varphi\sin\theta \\
    \sin\varphi\sin\theta & \rho\cos\varphi\sin\theta & \rho\sin\varphi\cos\theta \\
    \cos\varphi           & -\rho\sin\varphi           & 0
  \end{pmatrix}.
\]
Un calcul direct (développement par rapport à la troisième ligne) donne
\begin{equation}\label{eq:jacobien-spherique}
  \det J = \rho^2 \sin\varphi.
\end{equation}
Donc $\dd x\,\dd y\,\dd z = \rho^2 \sin\varphi\,\dd\rho\,\dd\varphi\,\dd\theta$.

\begin{proof}[Dérivation du jacobien sphérique]
  Développons selon la troisième ligne :
  \begin{align*}
    \det J &= \cos\varphi\bigl(\rho\cos\varphi\sin\theta \cdot
    (-\rho\sin\varphi\sin\theta) -
    \rho\sin\varphi\cos\theta\cdot\rho\cos\varphi\cos\theta\bigr) \cdot(-1)^{3+2}\\
    &\quad + (-\rho\sin\varphi)\bigl(\sin\varphi\cos\theta\cdot
    \rho\sin\varphi\cos\theta + \sin\varphi\sin\theta\cdot
    \rho\sin\varphi\sin\theta\bigr)\cdot(-1)^{3+2} \\
    &\quad \text{(en tenant compte de la convention de signe)}.
  \end{align*}
  Plus simplement, développons par la troisième ligne :
  \begin{align*}
    \det J &= \cos\varphi\,
    \det\begin{pmatrix}
      \sin\varphi\cos\theta & -\rho\sin\varphi\sin\theta \\
      \sin\varphi\sin\theta & \rho\sin\varphi\cos\theta
    \end{pmatrix} \\
    &\quad - (-\rho\sin\varphi)\,
    \det\begin{pmatrix}
      \sin\varphi\cos\theta & \rho\cos\varphi\cos\theta \\
      \sin\varphi\sin\theta & \rho\cos\varphi\sin\theta
    \end{pmatrix}.
  \end{align*}
  Le premier déterminant vaut
  $\rho\sin^2\!\varphi(\cos^2\!\theta + \sin^2\!\theta) = \rho\sin^2\!\varphi$.
  Le second vaut $0$ (les colonnes sont proportionnelles : la seconde
  est $\rho\cos\varphi$ fois la première\ldots{} non, calculons).
  En fait, le second déterminant vaut
  $\rho\cos\varphi\sin\varphi(\sin\theta\cos\theta -
  \cos\theta\sin\theta)  = 0$. Mais il faut aussi le troisième cofacteur.

  Reprenons proprement le développement selon la troisième ligne :
  \begin{align*}
    \det J &= \cos\varphi\cdot\rho\sin^2\!\varphi
    +\rho\sin\varphi\cdot\det\begin{pmatrix}
      \sin\varphi\cos\theta & \rho\cos\varphi\cos\theta \\
      \sin\varphi\sin\theta & \rho\cos\varphi\sin\theta
    \end{pmatrix} + 0.
  \end{align*}
  Non --- recalculons directement. Développons par la première colonne :
  \begin{align*}
    \det J &= \sin\varphi\cos\theta\bigl(\rho\cos\varphi\sin\theta\cdot
    0 - \rho\sin\varphi\cos\theta\cdot(-\rho\sin\varphi)\bigr) \\
    &\quad -\sin\varphi\sin\theta\bigl(\rho\cos\varphi\cos\theta\cdot
    0 - (-\rho\sin\varphi\sin\theta)\cdot(-\rho\sin\varphi)\bigr) \\
    &\quad + \cos\varphi\bigl(\rho\cos\varphi\cos\theta\cdot
    \rho\sin\varphi\cos\theta +
    \rho\sin\varphi\sin\theta\cdot\rho\cos\varphi\sin\theta\bigr).
  \end{align*}
  Simplifions chaque terme :
  \begin{align*}
    &= \sin\varphi\cos\theta\cdot\rho^2\sin^2\!\varphi\cos\theta
    - \sin\varphi\sin\theta\cdot(-\rho^2\sin^2\!\varphi\sin\theta) \\
    &\quad + \cos\varphi\cdot\rho^2\cos\varphi\sin\varphi
    (\cos^2\!\theta + \sin^2\!\theta) \\
    &= \rho^2\sin^3\!\varphi\cos^2\!\theta
    + \rho^2\sin^3\!\varphi\sin^2\!\theta
    + \rho^2\sin\varphi\cos^2\!\varphi \\
    &= \rho^2\sin^3\!\varphi + \rho^2\sin\varphi\cos^2\!\varphi
    = \rho^2\sin\varphi. \qedhere
  \end{align*}
\end{proof}


\begin{exemple}[Volume de la boule]\label{ex:volume-boule}
  \index{boule!volume}
  Le volume de la boule $B_R = \{(x,y,z) : x^2+y^2+z^2 \leqslant R^2\}$
  est
  \[
    V = \int_0^{2\pi}\!\int_0^\pi\!\int_0^R \rho^2\sin\varphi\,
    \dd\rho\,\dd\varphi\,\dd\theta
    = 2\pi\cdot\left[-\cos\varphi\right]_0^\pi\cdot\frac{R^3}{3}
    = 2\pi\cdot 2\cdot\frac{R^3}{3} = \frac{4}{3}\pi R^3.
  \]
\end{exemple}

\begin{exemple}[Moment d'inertie de la boule homogène]
  \label{ex:moment-inertie-boule}
  \index{moment d'inertie}
  Soit $B_R$ la boule de centre $O$ et de rayon $R$, de densité
  $\mu = 1$. Le moment d'inertie par rapport à l'axe $Oz$ est
  \[
    I_z = \iiint_{B_R}(x^2+y^2)\,\dd x\,\dd y\,\dd z.
  \]
  En coordonnées sphériques, $x^2+y^2 = \rho^2\sin^2\!\varphi$, d'où
  \begin{align*}
    I_z &= \int_0^{2\pi}\!\int_0^\pi\!\int_0^R
    \rho^2\sin^2\!\varphi\cdot\rho^2\sin\varphi\,
    \dd\rho\,\dd\varphi\,\dd\theta \\
    &= 2\pi\cdot\frac{R^5}{5}\cdot\int_0^\pi\sin^3\!\varphi\,\dd\varphi
    = 2\pi\cdot\frac{R^5}{5}\cdot\frac{4}{3}
    = \frac{8\pi R^5}{15}.
  \end{align*}
  En termes de la masse $M = \frac{4}{3}\pi R^3$, on obtient
  $I_z = \frac{2}{5}MR^2$.
\end{exemple}


% ============================================================
%  10.6 Calcul d'aires, volumes, centres de masse
% ============================================================

\section{Applications : aires, volumes, centres de masse}
\label{sec:applications-integrales-multiples}

\index{centre de masse}
\index{aire}

\begin{proposition}[Aire d'un domaine plan]\label{prop:aire-domaine}
  L'aire d'un domaine mesurable $D \subset \R^2$ est
  $\text{Aire}(D) = \iint_D \dd x\,\dd y$.
\end{proposition}

\begin{definition}[Centre de masse]\label{def:centre-masse}
  \index{centre de masse}
  Soit $D \subset \R^2$ un domaine de densité surfacique $\mu(x,y)$. La
  masse totale est $M = \iint_D \mu\,\dd x\,\dd y$, et le centre de
  masse $(\bar{x}, \bar{y})$ est défini par
  \[
    \bar{x} = \frac{1}{M}\iint_D x\,\mu(x,y)\,\dd x\,\dd y,
    \qquad
    \bar{y} = \frac{1}{M}\iint_D y\,\mu(x,y)\,\dd x\,\dd y.
  \]
\end{definition}

\begin{exemple}[Centre de masse d'un demi-disque]\label{ex:centre-masse}
  Soit $D = \{(x,y) : x^2+y^2 \leqslant R^2,\; y \geqslant 0\}$
  de densité uniforme $\mu = 1$. Par symétrie, $\bar{x} = 0$. La masse
  est $M = \frac{\pi R^2}{2}$. En coordonnées polaires,
  \[
    M\bar{y} = \int_0^\pi\!\int_0^R r\sin\theta\cdot r\,\dd r\,\dd\theta
    = \frac{R^3}{3}\int_0^\pi\sin\theta\,\dd\theta
    = \frac{R^3}{3}\cdot 2 = \frac{2R^3}{3}.
  \]
  Donc $\bar{y} = \frac{2R^3}{3}\cdot\frac{2}{\pi R^2} = \frac{4R}{3\pi}$.
\end{exemple}


% ============================================================
%  10.7 Exercices
% ============================================================

\section{Exercices}
\label{sec:exercices-chap10}

\begin{exercice}[$\star$]
  Calculer $\displaystyle\iint_R (x^2 - y)\,\dd x\,\dd y$ avec
  $R = [0,2]\times[-1,1]$.
\end{exercice}

\begin{exercice}[$\star$]
  Calculer $\displaystyle\iint_D xy\,\dd x\,\dd y$ où $D$ est le
  triangle de sommets $(0,0)$, $(1,0)$, $(1,1)$.
\end{exercice}

\begin{exercice}[$\star$]
  En utilisant le théorème de Fubini, montrer que
  $\displaystyle\int_0^1\!\int_0^x e^{x^2}\,\dd y\,\dd x
  = \frac{e-1}{2}$.
\end{exercice}

\begin{exercice}[$\star$]
  Calculer l'aire du domaine limité par les paraboles $y = x^2$ et
  $x = y^2$.
\end{exercice}

\begin{exercice}[$\star\star$]
  Calculer $\displaystyle\iint_D e^{x^2+y^2}\,\dd x\,\dd y$ où
  $D = \{(x,y) : x^2 + y^2 \leqslant 1\}$.
\end{exercice}

\begin{exercice}[$\star\star$]
  Calculer $\displaystyle\int_0^1\!\int_y^1 e^{x^2}\,\dd x\,\dd y$
  en intervertissant l'ordre d'intégration.
\end{exercice}

\begin{exercice}[$\star\star$]
  Calculer le volume du solide délimité par le paraboloïde
  $z = x^2 + y^2$ et le plan $z = 4$.
\end{exercice}

\begin{exercice}[$\star\star$]
  Soit $D$ le domaine défini par $1 \leqslant x^2 + y^2 \leqslant 4$
  et $y \geqslant 0$. Calculer $\displaystyle\iint_D \frac{1}
  {(x^2+y^2)^{3/2}}\,\dd x\,\dd y$.
\end{exercice}

\begin{exercice}[$\star\star$]
  Calculer le centre de masse du quart de disque
  $D = \{(x,y) : x^2+y^2 \leqslant R^2,\; x \geqslant 0,\;
  y \geqslant 0\}$ de densité uniforme.
\end{exercice}

\begin{exercice}[$\star\star\star$]
  Soit $f\colon\R\to\R$ continue. Montrer que pour tout $x > 0$,
  \[
    \int_0^x\!\int_0^{t_1}\!\int_0^{t_2} f(t_3)\,\dd t_3\,\dd t_2\,
    \dd t_1 = \frac{1}{2}\int_0^x (x-t)^2 f(t)\,\dd t.
  \]
  \emph{Indication : intervertir les ordres d'intégration.}
\end{exercice}

\begin{exercice}[$\star\star\star$]
  Calculer $\displaystyle\iiint_{B_R}\frac{1}
  {\sqrt{x^2+y^2+z^2}}\,\dd x\,\dd y\,\dd z$ où $B_R$ est la boule
  de rayon $R$ centrée à l'origine.
\end{exercice}

\begin{exercice}[$\star\star\star$]
  Soit $T$ le tétraèdre $\{(x,y,z) : x,y,z \geqslant 0,\;
  x+y+z \leqslant 1\}$. Calculer son moment d'inertie par rapport
  à l'axe $Oz$, en supposant une densité uniforme.
\end{exercice}


\bigskip
\begin{center}
\fbox{\parbox{0.85\textwidth}{%
\textbf{Résumé du chapitre~\thechapter.}
\begin{itemize}[nosep]
  \item L'intégrale double se définit via les sommes de Darboux sur un
    rectangle ; toute fonction continue est intégrable.
  \item Le théorème de Fubini ramène une intégrale double à des
    intégrales itérées simples.
  \item Sur un domaine de type~I (resp.~II), on intègre d'abord en $y$
    (resp.~$x$) entre les courbes frontières.
  \item Le changement de variables fait intervenir le jacobien
    $\abs{\det J_\Phi}$.
  \item Coordonnées polaires : $\dd x\,\dd y = r\,\dd r\,\dd\theta$.
  \item Coordonnées cylindriques :
    $\dd x\,\dd y\,\dd z = r\,\dd r\,\dd\theta\,\dd z$.
  \item Coordonnées sphériques :
    $\dd x\,\dd y\,\dd z = \rho^2\sin\varphi\,
    \dd\rho\,\dd\varphi\,\dd\theta$.
  \item L'intégrale de Gauss $\int_{-\infty}^{+\infty}e^{-x^2}\dd x
    = \sqrt{\pi}$ se démontre par passage en polaires.
\end{itemize}}}
\end{center}


% ============================================================
% ============================================================
%  Chapitre 11 : Introduction aux Courbes et Intégrales Curvilignes
% ============================================================
% ============================================================

\chapter{Introduction aux Courbes et Intégrales Curvilignes}
\label{chap:courbes-integrales-curvilignes}

\index{courbe paramétrée}
\index{intégrale curviligne}

Le mouvement d'un point matériel dans le plan ou l'espace, le travail
d'une force le long d'une trajectoire, la longueur d'un fil courbe :
autant de situations qui exigent de savoir \emph{intégrer le long
d'une courbe}. Ce chapitre développe le formalisme des courbes
paramétrées, définit rigoureusement la notion de longueur d'arc, puis
introduit les deux types d'intégrales curvilignes (scalaire et
vectorielle). Il culmine avec le théorème de Green, qui relie une
intégrale curviligne à une intégrale double.


% ============================================================
%  11.1 Courbes paramétrées
% ============================================================

\section{Courbes paramétrées}
\label{sec:courbes-parametrees}

\index{courbe!paramétrée}
\index{courbe!régulière}

\begin{definition}[Courbe paramétrée]\label{def:courbe-parametree}
  Un \emph{arc paramétré}\index{arc paramétré} de $\R^n$ est une
  application $\gamma\colon [a,b] \to \R^n$ de classe $\mathcal{C}^0$.
  L'image $\gamma([a,b])$ est le \emph{support} (ou \emph{trace}) de
  la courbe. Les points $\gamma(a)$ et $\gamma(b)$ sont les
  \emph{extrémités} ; la courbe est dite \emph{fermée} si
  $\gamma(a) = \gamma(b)$.
\end{definition}

\begin{definition}[Courbe régulière]\label{def:courbe-reguliere}
  L'arc paramétré $\gamma\colon [a,b]\to\R^n$ est dit \emph{de classe
  $\mathcal{C}^1$} si $\gamma$ est de classe $\mathcal{C}^1$ sur
  $[a,b]$. Il est \emph{régulier}\index{courbe!régulière} si de plus
  $\gamma'(t) \neq 0$ pour tout $t \in [a,b]$. Il est
  \emph{$\mathcal{C}^1$ par morceaux}\index{courbe!C1 par morceaux@$\mathcal{C}^1$
  par morceaux} s'il existe une subdivision $a = t_0 < t_1 < \cdots < t_k = b$
  telle que $\gamma|_{[t_{i-1},t_i]}$ est de classe $\mathcal{C}^1$ pour
  tout $i$.
\end{definition}

\begin{definition}[Vecteur tangent]\label{def:vecteur-tangent}
  \index{vecteur tangent}
  Si $\gamma$ est de classe $\mathcal{C}^1$ et $\gamma'(t_0) \neq 0$,
  le \emph{vecteur tangent} à la courbe au point $\gamma(t_0)$ est
  $\gamma'(t_0)$, et la \emph{droite tangente}\index{droite tangente} est
  l'ensemble
  \[
    \bigl\{\gamma(t_0) + s\,\gamma'(t_0) : s \in \R\bigr\}.
  \]
\end{definition}

\begin{center}
\begin{tikzpicture}[scale=1.0,
  decoration={markings,mark=at position 0.35 with {\arrow{>}},
  mark=at position 0.7 with {\arrow{>}}}]
  \draw[blue,thick,postaction={decorate}]
    plot[domain=0.3:5.5,samples=200] ({1.8*cos(\x r)+0.3*\x},{1.2*sin(\x r)+0.2*\x});
  % Point et vecteur tangent
  \pgfmathsetmacro{\tval}{2.5}
  \pgfmathsetmacro{\px}{1.8*cos(\tval r)+0.3*\tval}
  \pgfmathsetmacro{\py}{1.2*sin(\tval r)+0.2*\tval}
  \pgfmathsetmacro{\dx}{-1.8*sin(\tval r)+0.3}
  \pgfmathsetmacro{\dy}{1.2*cos(\tval r)+0.2}
  \fill[red] (\px,\py) circle (2pt) node[above left] {$\gamma(t_0)$};
  \draw[->,red,very thick] (\px,\py) -- ++(0.8*\dx,0.8*\dy)
    node[right] {$\gamma'(t_0)$};
  % Autre point
  \pgfmathsetmacro{\tval}{4.2}
  \pgfmathsetmacro{\px}{1.8*cos(\tval r)+0.3*\tval}
  \pgfmathsetmacro{\py}{1.2*sin(\tval r)+0.2*\tval}
  \pgfmathsetmacro{\dx}{-1.8*sin(\tval r)+0.3}
  \pgfmathsetmacro{\dy}{1.2*cos(\tval r)+0.2}
  \fill[red] (\px,\py) circle (2pt);
  \draw[->,red,very thick] (\px,\py) -- ++(0.8*\dx,0.8*\dy)
    node[above] {$\gamma'(t_1)$};
  % Labels
  \node[blue] at (-1.5,1.3) {$\gamma$};
\end{tikzpicture}
\end{center}


\begin{exemple}[Exemples fondamentaux]\label{ex:courbes-classiques}
  \begin{enumerate}[label=\textup{(\roman*)}]
    \item \textbf{Cercle}\index{cercle (paramétrisation)} :
      $\gamma(t) = (R\cos t,\; R\sin t)$, $t \in [0, 2\pi]$. On a
      $\gamma'(t) = (-R\sin t,\; R\cos t)$, de norme $R \neq 0$ :
      la courbe est régulière.
    \item \textbf{Ellipse}\index{ellipse} :
      $\gamma(t) = (a\cos t,\; b\sin t)$, $t \in [0, 2\pi]$,
      $a, b > 0$.
    \item \textbf{Cycloïde}\index{cycloïde} :
      $\gamma(t) = (t - \sin t,\; 1 - \cos t)$, $t \in [0, 2\pi]$.
      On a $\gamma'(t) = (1-\cos t,\; \sin t)$, qui s'annule en
      $t = 0$ et $t = 2\pi$ : la courbe n'est pas régulière mais est
      $\mathcal{C}^1$ par morceaux.
    \item \textbf{Hélice circulaire}\index{hélice} :
      $\gamma(t) = (\cos t,\; \sin t,\; t)$, $t \in [0, 4\pi]$, dans
      $\R^3$. On a $\gamma'(t) = (-\sin t,\; \cos t,\; 1)$, de norme
      $\sqrt{2}$ : la courbe est régulière.
  \end{enumerate}
\end{exemple}


% ============================================================
%  11.2 Longueur d'arc
% ============================================================

\section{Longueur d'arc}
\label{sec:longueur-arc}

\index{longueur d'arc}

\begin{definition}[Longueur d'une courbe]\label{def:longueur-courbe}
  Soit $\gamma\colon[a,b]\to\R^n$ un arc continu. Pour toute
  subdivision $\sigma = (a = t_0 < t_1 < \cdots < t_k = b)$, on pose
  \[
    \ell(\gamma,\sigma) = \sum_{i=1}^k \norm{\gamma(t_i) - \gamma(t_{i-1})}.
  \]
  La \emph{longueur} de $\gamma$ est
  $L(\gamma) = \sup_\sigma \ell(\gamma,\sigma) \in [0, +\infty]$.
  La courbe est dite \emph{rectifiable}\index{courbe!rectifiable} si
  $L(\gamma) < +\infty$.
\end{definition}

\begin{theoreme}[Formule de la longueur d'arc]
  \label{thm:longueur-arc}
  \index{longueur d'arc!formule}
  Si $\gamma\colon[a,b]\to\R^n$ est de classe $\mathcal{C}^1$, alors
  $\gamma$ est rectifiable et
  \begin{equation}\label{eq:longueur-arc}
    L(\gamma) = \int_a^b \norm{\gamma'(t)}\,\dd t.
  \end{equation}
\end{theoreme}

\begin{proof}
  \textsc{Majoration.} Soit $\sigma = (t_0,\ldots,t_k)$ une
  subdivision de $[a,b]$. Pour chaque $i$, par le théorème fondamental
  de l'analyse,
  \[
    \gamma(t_i) - \gamma(t_{i-1}) = \int_{t_{i-1}}^{t_i}
    \gamma'(t)\,\dd t,
  \]
  d'où, par l'inégalité triangulaire pour les intégrales,
  \[
    \norm{\gamma(t_i) - \gamma(t_{i-1})}
    \leqslant \int_{t_{i-1}}^{t_i} \norm{\gamma'(t)}\,\dd t.
  \]
  En sommant sur $i$,
  $\ell(\gamma,\sigma) \leqslant \int_a^b\norm{\gamma'(t)}\,\dd t$.
  En prenant le supremum sur $\sigma$,
  $L(\gamma) \leqslant \int_a^b\norm{\gamma'(t)}\,\dd t < +\infty$.

  \textsc{Minoration.} Puisque $\gamma'$ est continue sur $[a,b]$
  (compact), elle est uniformément continue. Soit $\eps > 0$. Il existe
  $\delta > 0$ tel que $\abs{s-t} < \delta$ implique
  $\norm{\gamma'(s)-\gamma'(t)} < \eps$. Prenons une subdivision
  $\sigma$ de pas $\norm{\sigma} < \delta$. Pour $t \in [t_{i-1},t_i]$,
  \[
    \norm{\gamma'(t)} \leqslant \norm{\gamma'(t_{i-1})} + \eps,
  \]
  et de même $\norm{\gamma'(t_{i-1})} \leqslant \norm{\gamma'(t)}+\eps$,
  d'où $\norm{\gamma'(t)} \geqslant \norm{\gamma'(t_{i-1})} - \eps$.
  Par conséquent,
  \begin{align*}
    \norm{\gamma(t_i)-\gamma(t_{i-1})}
    &= \norm*{\int_{t_{i-1}}^{t_i}\gamma'(t)\,\dd t}
    \geqslant \norm*{\gamma'(t_{i-1})(t_i-t_{i-1})}
    - \int_{t_{i-1}}^{t_i}\norm{\gamma'(t)-\gamma'(t_{i-1})}\,\dd t \\
    &\geqslant \norm{\gamma'(t_{i-1})}(t_i-t_{i-1})
    - \eps(t_i-t_{i-1}).
  \end{align*}
  En sommant,
  \[
    \ell(\gamma,\sigma) \geqslant
    \sum_{i=1}^k\norm{\gamma'(t_{i-1})}(t_i-t_{i-1})
    - \eps(b-a).
  \]
  Or la somme de droite est une somme de Riemann pour
  $\int_a^b\norm{\gamma'(t)}\,\dd t$. En raffinant $\sigma$,
  \[
    L(\gamma) \geqslant \int_a^b\norm{\gamma'(t)}\,\dd t - \eps(b-a).
  \]
  Puisque $\eps > 0$ est arbitraire,
  $L(\gamma) \geqslant \int_a^b\norm{\gamma'(t)}\,\dd t$.
\end{proof}

\begin{definition}[Abscisse curviligne]\label{def:abscisse-curviligne}
  \index{abscisse curviligne}
  Soit $\gamma\colon[a,b]\to\R^n$ un arc $\mathcal{C}^1$ régulier.
  L'\emph{abscisse curviligne} à partir du point $\gamma(a)$ est la
  fonction $s\colon[a,b]\to[0,L]$ définie par
  \[
    s(t) = \int_a^t \norm{\gamma'(u)}\,\dd u.
  \]
  Puisque $\gamma$ est régulier, $s'(t) = \norm{\gamma'(t)} > 0$, donc
  $s$ est un difféomorphisme croissant de $[a,b]$ sur $[0,L]$.
\end{definition}

\begin{proposition}[Paramétrisation par longueur d'arc]
  \label{prop:param-longueur-arc}
  \index{paramétrisation par longueur d'arc}
  Si $\gamma\colon[a,b]\to\R^n$ est un arc $\mathcal{C}^1$ régulier et
  $s$ est l'abscisse curviligne, alors $\tilde{\gamma} = \gamma\circ s^{-1}
  \colon [0,L]\to\R^n$ vérifie
  $\norm{\tilde{\gamma}'(\sigma)} = 1$ pour tout $\sigma \in [0,L]$.
\end{proposition}

\begin{proof}
  Par la règle de dérivation des fonctions composées,
  $\tilde{\gamma}'(\sigma) = \gamma'(s^{-1}(\sigma))\cdot
  (s^{-1})'(\sigma)$.
  Or $(s^{-1})'(\sigma) = \frac{1}{s'(t)} = \frac{1}{\norm{\gamma'(t)}}$
  en $t = s^{-1}(\sigma)$. Donc
  $\norm{\tilde{\gamma}'(\sigma)} = \frac{\norm{\gamma'(t)}}
  {\norm{\gamma'(t)}} = 1$.
\end{proof}

\begin{exemple}[Longueur du cercle]\label{ex:longueur-cercle}
  Pour le cercle $\gamma(t) = (R\cos t, R\sin t)$, $t \in [0,2\pi]$,
  on a $\norm{\gamma'(t)} = R$, d'où
  $L = \int_0^{2\pi}R\,\dd t = 2\pi R$.
\end{exemple}

\begin{exemple}[Longueur de la cycloïde]\label{ex:longueur-cycloide}
  Pour $\gamma(t) = (t-\sin t, 1-\cos t)$, $t \in [0,2\pi]$,
  $\norm{\gamma'(t)}^2 = (1-\cos t)^2 + \sin^2 t = 2(1-\cos t)
  = 4\sin^2(t/2)$, donc $\norm{\gamma'(t)} = 2\abs{\sin(t/2)}
  = 2\sin(t/2)$ sur $[0,2\pi]$, et
  \[
    L = \int_0^{2\pi}2\sin\frac{t}{2}\,\dd t
    = \left[-4\cos\frac{t}{2}\right]_0^{2\pi} = -4(-1) + 4(1) = 8.
  \]
\end{exemple}


% ============================================================
%  11.3 Intégrale curviligne de première espèce
% ============================================================

\section{Intégrales curvilignes de première espèce}
\label{sec:int-curv-premiere-espece}

\index{intégrale curviligne!de première espèce}
\index{intégrale curviligne!scalaire}

\begin{definition}[Intégrale curviligne scalaire]\label{def:int-curv-scalaire}
  Soit $\gamma\colon[a,b]\to\R^n$ un arc $\mathcal{C}^1$ (ou
  $\mathcal{C}^1$ par morceaux) et $f$ une fonction continue définie
  sur le support de $\gamma$. L'\emph{intégrale curviligne de première
  espèce} (ou \emph{intégrale scalaire}) de $f$ le long de $\gamma$ est
  \begin{equation}\label{eq:int-curv-scalaire}
    \int_\gamma f\,\dd s
    = \int_a^b f\bigl(\gamma(t)\bigr)\,\norm{\gamma'(t)}\,\dd t.
  \end{equation}
\end{definition}

\begin{remarque}\label{rem:invariance-param}
  Cette intégrale ne dépend pas de la paramétrisation choisie pour la
  courbe : si $\tilde{\gamma} = \gamma\circ\varphi$ est un
  reparamétrage admissible (avec $\varphi$ croissant et de classe
  $\mathcal{C}^1$), alors
  $\int_{\tilde{\gamma}}f\,\dd s = \int_\gamma f\,\dd s$.
\end{remarque}

\begin{exemple}[Masse d'un fil]\label{ex:masse-fil}
  \index{masse d'un fil}
  Un fil en forme de demi-cercle $\gamma(t) = (\cos t, \sin t)$,
  $t \in [0,\pi]$, a une densité linéique $\mu(x,y) = y$. Sa masse est
  \[
    M = \int_\gamma \mu\,\dd s
    = \int_0^\pi \sin t \cdot 1\,\dd t
    = [-\cos t]_0^\pi = 2.
  \]
\end{exemple}


% ============================================================
%  11.4 Intégrale curviligne de seconde espèce
% ============================================================

\section{Intégrales curvilignes de seconde espèce}
\label{sec:int-curv-seconde-espece}

\index{intégrale curviligne!de seconde espèce}
\index{travail d'une force}

\begin{definition}{Intégrale curviligne de seconde espèce (travail)}%
  {int-curv-travail}
  Soit $\gamma\colon[a,b]\to\R^n$ un arc $\mathcal{C}^1$ par morceaux
  et $\mathbf{F}\colon\Omega \to \R^n$ un champ de vecteurs continu
  défini sur un ouvert $\Omega$ contenant le support de $\gamma$.
  L'\emph{intégrale curviligne de seconde espèce} (ou \emph{circulation})
  de $\mathbf{F}$ le long de $\gamma$ est
  \begin{equation}\label{eq:int-curv-travail}
    \int_\gamma \mathbf{F}\cdot\dd\gamma
    = \int_a^b \mathbf{F}\bigl(\gamma(t)\bigr)\cdot\gamma'(t)\,\dd t.
  \end{equation}
\end{definition}

\begin{remarque}\label{rem:travail-physique}
  En mécanique, si $\mathbf{F}$ est un champ de forces et $\gamma$ une
  trajectoire, $\int_\gamma\mathbf{F}\cdot\dd\gamma$ représente le
  \emph{travail}\index{travail} de la force $\mathbf{F}$ le long du
  chemin $\gamma$. Contrairement à l'intégrale de première espèce,
  cette intégrale \emph{change de signe} si l'on renverse le sens de
  parcours.
\end{remarque}

\begin{center}
\begin{tikzpicture}[scale=1.0,
  decoration={markings,mark=at position 0.5 with {\arrow{>}}}]
  % Courbe
  \draw[blue,thick,postaction={decorate}]
    plot[domain=0:3.5,samples=200] (\x, {0.8*sin(70*\x)+0.3*\x});
  \node[blue,below] at (0,0) {$\gamma(a)$};
  \node[blue,above] at (3.5,{0.8*sin(70*3.5)+0.3*3.5}) {$\gamma(b)$};
  % Champ de vecteurs
  \foreach \x in {0.5,1,...,3} {
    \foreach \yy in {-0.5,0,0.5,1,1.5} {
      \pgfmathsetmacro{\fx}{0.2*\yy}
      \pgfmathsetmacro{\fy}{0.2*(1-0.2*\x)}
      \draw[->,gray!60,thin] (\x,\yy) -- ++(\fx,\fy);
    }
  }
  % Un vecteur sur la courbe
  \pgfmathsetmacro{\tval}{1.8}
  \pgfmathsetmacro{\px}{\tval}
  \pgfmathsetmacro{\py}{0.8*sin(70*\tval)+0.3*\tval}
  \fill[red] (\px,\py) circle (1.5pt);
  \draw[->,red,thick] (\px,\py) -- ++(0.35,0.3)
    node[above right] {$\mathbf{F}(\gamma(t))$};
  \draw[->,orange,thick] (\px,\py) -- ++(0.4,0.1)
    node[below right] {$\gamma'(t)$};
  \node at (5,0.8) {$\int_\gamma\mathbf{F}\cdot\dd\gamma$};
\end{tikzpicture}
\end{center}

\begin{exemple}\label{ex:travail-force}
  Soit $\mathbf{F}(x,y) = (y,\; x)$ et $\gamma(t) = (\cos t,\;\sin t)$,
  $t \in [0,\pi/2]$. Alors $\gamma'(t) = (-\sin t,\;\cos t)$ et
  \begin{align*}
    \int_\gamma\mathbf{F}\cdot\dd\gamma
    &= \int_0^{\pi/2}\bigl(\sin t\cdot(-\sin t) + \cos t\cdot\cos t\bigr)\dd t \\
    &= \int_0^{\pi/2}(\cos^2 t - \sin^2 t)\,\dd t
    = \int_0^{\pi/2}\cos(2t)\,\dd t
    = \left[\frac{\sin(2t)}{2}\right]_0^{\pi/2} = 0.
  \end{align*}
\end{exemple}

Lorsque $\mathbf{F} = (F_1, F_2)$ dans $\R^2$, on note aussi
$\int_\gamma \mathbf{F}\cdot\dd\gamma = \int_\gamma F_1\,\dd x
+ F_2\,\dd y$, qu'on appelle une \emph{forme différentielle}\index{forme différentielle}.


% ============================================================
%  11.5 Champs de gradient et potentiel
% ============================================================

\section{Champs de gradient et potentiel}
\label{sec:gradient-potentiel}

\index{champ de gradient}
\index{potentiel}

\begin{definition}[Champ de gradient]\label{def:champ-gradient}
  Un champ de vecteurs $\mathbf{F}\colon\Omega\to\R^n$ (avec $\Omega$
  ouvert de $\R^n$) est un \emph{champ de gradient}\index{champ de gradient}
  (ou \emph{champ conservatif}\index{champ conservatif}) s'il existe une
  fonction $V\colon\Omega\to\R$ de classe $\mathcal{C}^1$ telle que
  $\mathbf{F} = \nabla V$. La fonction $V$ est appelée un
  \emph{potentiel}\index{potentiel} de $\mathbf{F}$.
\end{definition}

\begin{theoreme}[Formule fondamentale pour les champs de gradient]
  \label{thm:gradient-potentiel}
  \index{formule fondamentale!champs de gradient}
  Soit $\mathbf{F} = \nabla V$ un champ de gradient sur un ouvert
  $\Omega$ et $\gamma\colon[a,b]\to\Omega$ un arc $\mathcal{C}^1$ par
  morceaux. Alors
  \begin{equation}\label{eq:gradient-potentiel}
    \int_\gamma \mathbf{F}\cdot\dd\gamma
    = V\bigl(\gamma(b)\bigr) - V\bigl(\gamma(a)\bigr).
  \end{equation}
  En particulier, l'intégrale ne dépend que des extrémités et, si
  $\gamma$ est fermée, elle est nulle.
\end{theoreme}

\begin{proof}
  Par la règle de dérivation des fonctions composées, la fonction
  $t \mapsto V(\gamma(t))$ est dérivable et
  \[
    \frac{\dd}{\dd t}V\bigl(\gamma(t)\bigr)
    = \nabla V\bigl(\gamma(t)\bigr)\cdot\gamma'(t)
    = \mathbf{F}\bigl(\gamma(t)\bigr)\cdot\gamma'(t).
  \]
  Le théorème fondamental de l'analyse (chapitre~7) donne alors
  \[
    \int_\gamma\mathbf{F}\cdot\dd\gamma
    = \int_a^b \frac{\dd}{\dd t}V\bigl(\gamma(t)\bigr)\,\dd t
    = V\bigl(\gamma(b)\bigr) - V\bigl(\gamma(a)\bigr). \qedhere
  \]
\end{proof}

\begin{corollaire}\label{cor:circulation-nulle-gradient}
  Si $\mathbf{F} = \nabla V$ et $\gamma$ est une courbe fermée dans
  $\Omega$, alors $\displaystyle\int_\gamma\mathbf{F}\cdot\dd\gamma = 0$.
\end{corollaire}

\begin{proposition}[Condition nécessaire d'exactitude]
  \label{prop:condition-necessaire-exacte}
  \index{condition d'exactitude}
  \index{Schwarz (condition de)}
  Si $\mathbf{F} = (F_1, F_2)$ est un champ de gradient de classe
  $\mathcal{C}^1$ sur un ouvert $\Omega \subset \R^2$, alors
  \begin{equation}\label{eq:condition-exacte}
    \frac{\partial F_1}{\partial y} = \frac{\partial F_2}{\partial x}
    \qquad\text{sur } \Omega.
  \end{equation}
\end{proposition}

\begin{proof}
  Si $\mathbf{F} = \nabla V = \bigl(\frac{\partial V}{\partial x},
  \frac{\partial V}{\partial y}\bigr)$, alors
  $\frac{\partial F_1}{\partial y} = \frac{\partial^2 V}{\partial y\,
  \partial x}$ et $\frac{\partial F_2}{\partial x} = \frac{\partial^2 V}
  {\partial x\,\partial y}$. Le théorème de Schwarz\index{Schwarz (théorème de)}
  (les dérivées partielles croisées de classe $\mathcal{C}^1$ sont
  égales) donne le résultat.
\end{proof}

\begin{remarque}\label{rem:condition-suffisante}
  La condition~\eqref{eq:condition-exacte} n'est pas suffisante en
  général ; elle l'est si $\Omega$ est \emph{simplement
  connexe}\index{simplement connexe} (par exemple convexe, ou étoilé).
  Nous en verrons un contre-exemple célèbre en section~\ref{sec:contre-exemple}.
\end{remarque}


% ============================================================
%  11.6 Théorème de Green
% ============================================================

\section{Théorème de Green}
\label{sec:theoreme-green}

\index{Green (théorème de)}

\begin{theoreme}[Green]\label{thm:green}
  Soit $D \subset \R^2$ un domaine borné dont la frontière $\partial D$
  est une courbe fermée $\mathcal{C}^1$ par morceaux, parcourue dans le
  sens trigonométrique (positif). Soient $P, Q\colon \overline{D}\to\R$
  de classe $\mathcal{C}^1$. Alors
  \begin{equation}\label{eq:green}
    \oint_{\partial D}(P\,\dd x + Q\,\dd y)
    = \iint_D \left(\frac{\partial Q}{\partial x}
    - \frac{\partial P}{\partial y}\right)\dd x\,\dd y.
  \end{equation}
\end{theoreme}

\begin{center}
\begin{tikzpicture}[scale=1.2,
  decoration={markings,mark=at position 0.15 with {\arrow{>}},
  mark=at position 0.4 with {\arrow{>}},
  mark=at position 0.65 with {\arrow{>}},
  mark=at position 0.9 with {\arrow{>}}}]
  \fill[green!10] plot[smooth cycle,tension=0.8]
    coordinates {(0,0) (1.5,-0.3) (3,0.2) (3.2,1.5) (2.5,2.5)
    (1,2.2) (-0.3,1.2)};
  \draw[blue,thick,postaction={decorate}]
    plot[smooth cycle,tension=0.8]
    coordinates {(0,0) (1.5,-0.3) (3,0.2) (3.2,1.5) (2.5,2.5)
    (1,2.2) (-0.3,1.2)};
  \node at (1.5,1.1) {\Large$D$};
  \node[blue] at (3.5,0.8) {$\partial D$};
  % Double flèche illustrant la relation
  \node at (5,1.1) {$\displaystyle\oint_{\partial D}\leftrightarrow
    \iint_D$};
\end{tikzpicture}
\end{center}

\begin{proof}[Idée de la preuve pour un domaine de type~I et~II]
  Supposons que $D$ soit à la fois de type~I et de type~II. Montrons
  séparément que
  \begin{equation}\label{eq:green-P}
    \oint_{\partial D}P\,\dd x = -\iint_D \frac{\partial P}{\partial y}
    \,\dd x\,\dd y
  \end{equation}
  et
  \begin{equation}\label{eq:green-Q}
    \oint_{\partial D}Q\,\dd y = \iint_D \frac{\partial Q}{\partial x}
    \,\dd x\,\dd y.
  \end{equation}
  La somme donnera~\eqref{eq:green}.

  \emph{Preuve de~\eqref{eq:green-P}.} Écrivons $D$ comme domaine de
  type~I : $D = \{(x,y) : a \leqslant x \leqslant b,\;
  g_1(x) \leqslant y \leqslant g_2(x)\}$. Alors
  \[
    \iint_D \frac{\partial P}{\partial y}\,\dd x\,\dd y
    = \int_a^b\!\int_{g_1(x)}^{g_2(x)} \frac{\partial P}{\partial y}
    (x,y)\,\dd y\,\dd x
    = \int_a^b \bigl[P(x,g_2(x)) - P(x,g_1(x))\bigr]\,\dd x.
  \]
  D'autre part, la frontière $\partial D$, parcourue dans le sens
  positif, se décompose en :
  \begin{itemize}[nosep]
    \item le chemin inférieur $C_1\colon x \mapsto (x, g_1(x))$, de $a$
      à $b$ (sens positif en $x$),
    \item le chemin supérieur $C_2\colon x \mapsto (x, g_2(x))$, de $b$
      à $a$ (sens négatif en $x$).
  \end{itemize}
  On a
  \begin{align*}
    \oint_{\partial D}P\,\dd x
    &= \int_a^b P(x,g_1(x))\,\dd x + \int_b^a P(x,g_2(x))\,\dd x \\
    &= \int_a^b P(x,g_1(x))\,\dd x - \int_a^b P(x,g_2(x))\,\dd x \\
    &= -\int_a^b \bigl[P(x,g_2(x))-P(x,g_1(x))\bigr]\,\dd x
    = -\iint_D\frac{\partial P}{\partial y}\,\dd x\,\dd y.
  \end{align*}
  La preuve de~\eqref{eq:green-Q} est analogue en écrivant $D$ comme
  domaine de type~II.
\end{proof}

\begin{corollaire}[Formule de l'aire]\label{cor:aire-green}
  \index{aire!formule de Green}
  Si $D$ est un domaine vérifiant les hypothèses du théorème de Green,
  alors
  \begin{equation}\label{eq:aire-green}
    \text{Aire}(D) = \frac{1}{2}\oint_{\partial D}(x\,\dd y - y\,\dd x).
  \end{equation}
\end{corollaire}

\begin{proof}
  On applique~\eqref{eq:green} avec $P = -y$ et $Q = x$. Alors
  $\frac{\partial Q}{\partial x} - \frac{\partial P}{\partial y}
  = 1 - (-1) = 2$, d'où
  $\oint_{\partial D}(-y\,\dd x + x\,\dd y) = \iint_D 2\,\dd x\,\dd y
  = 2\,\text{Aire}(D)$.
\end{proof}

\begin{exemple}[Aire de l'ellipse]\label{ex:aire-ellipse-green}
  Pour l'ellipse $\gamma(t) = (a\cos t,\; b\sin t)$, $t \in [0,2\pi]$,
  \begin{align*}
    \text{Aire} &= \frac{1}{2}\int_0^{2\pi}
    \bigl(a\cos t\cdot b\cos t - b\sin t\cdot(-a\sin t)\bigr)\,\dd t \\
    &= \frac{ab}{2}\int_0^{2\pi}(\cos^2 t + \sin^2 t)\,\dd t
    = \frac{ab}{2}\cdot 2\pi = \pi ab.
  \end{align*}
\end{exemple}


% ============================================================
%  11.7 Contre-exemple : forme fermée non exacte
% ============================================================

\section{Contre-exemple : forme fermée non exacte}
\label{sec:contre-exemple}

\index{forme fermée non exacte}

Considérons le champ de vecteurs
\begin{equation}\label{eq:champ-non-exact}
  \mathbf{F}(x,y) = \left(\frac{-y}{x^2+y^2},\;
  \frac{x}{x^2+y^2}\right),
\end{equation}
défini sur $\Omega = \R^2\setminus\{(0,0)\}$.

\begin{proposition}\label{prop:forme-fermee}
  Le champ $\mathbf{F}$ vérifie $\frac{\partial F_1}{\partial y}
  = \frac{\partial F_2}{\partial x}$ sur $\Omega$ (la forme
  $F_1\,\dd x + F_2\,\dd y$ est \emph{fermée}), mais $\mathbf{F}$
  n'est pas un champ de gradient sur $\Omega$.
\end{proposition}

\begin{proof}
  \emph{La forme est fermée.} Calculons :
  \[
    \frac{\partial F_1}{\partial y}
    = \frac{-(x^2+y^2)+y\cdot 2y}{(x^2+y^2)^2}
    = \frac{y^2-x^2}{(x^2+y^2)^2},
  \]
  \[
    \frac{\partial F_2}{\partial x}
    = \frac{(x^2+y^2)-x\cdot 2x}{(x^2+y^2)^2}
    = \frac{y^2-x^2}{(x^2+y^2)^2}.
  \]
  On a bien $\frac{\partial F_1}{\partial y} = \frac{\partial F_2}
  {\partial x}$.

  \emph{La forme n'est pas exacte.} Prenons $\gamma(t) = (\cos t,
  \sin t)$, $t \in [0, 2\pi]$, qui est une courbe fermée dans $\Omega$.
  Alors $\gamma'(t) = (-\sin t, \cos t)$ et
  \begin{align*}
    \int_\gamma\mathbf{F}\cdot\dd\gamma
    &= \int_0^{2\pi}\left(\frac{-\sin t}{1}\cdot(-\sin t)
    + \frac{\cos t}{1}\cdot\cos t\right)\dd t \\
    &= \int_0^{2\pi}(\sin^2 t + \cos^2 t)\,\dd t = 2\pi \neq 0.
  \end{align*}
  Si $\mathbf{F}$ était un gradient, la circulation sur toute courbe
  fermée serait nulle (corollaire~\ref{cor:circulation-nulle-gradient}).
  C'est une contradiction.
\end{proof}

\begin{remarque}\label{rem:argument}
  Sur le demi-plan $\{(x,y) : x > 0\}$ (ou tout ouvert simplement
  connexe contenu dans $\Omega$), $\mathbf{F}$ \emph{est} un gradient :
  $\mathbf{F} = \nabla\bigl(\arctan(y/x)\bigr)$. Le défaut d'exactitude
  globale provient de la topologie non triviale de
  $\Omega = \R^2\setminus\{0\}$, qui n'est pas simplement connexe.
  L'intégrale $\frac{1}{2\pi}\int_\gamma\mathbf{F}\cdot\dd\gamma$
  mesure l'\emph{indice de $\gamma$ par rapport à l'origine}%
  \index{indice d'une courbe}, un invariant topologique.
\end{remarque}


% ============================================================
%  11.8 Exercices
% ============================================================

\section{Exercices}
\label{sec:exercices-chap11}

\begin{exercice}[$\star$]
  Calculer la longueur de l'arc $\gamma(t) = (t^2, t^3)$ pour
  $t \in [0,1]$.
\end{exercice}

\begin{exercice}[$\star$]
  Calculer la longueur d'une arche de cycloïde
  $\gamma(t) = (R(t-\sin t),\; R(1-\cos t))$, $t \in [0,2\pi]$.
\end{exercice}

\begin{exercice}[$\star$]
  Calculer l'intégrale curviligne $\int_\gamma(x^2+y^2)\,\dd s$ où
  $\gamma$ est le quart de cercle unité $\gamma(t) = (\cos t,\sin t)$,
  $t \in [0,\pi/2]$.
\end{exercice}

\begin{exercice}[$\star$]
  Soit $\mathbf{F}(x,y) = (2xy,\; x^2)$. Vérifier que $\mathbf{F}$
  est un champ de gradient et trouver un potentiel $V$.
\end{exercice}

\begin{exercice}[$\star\star$]
  Calculer la circulation de $\mathbf{F}(x,y) = (x-y,\; x+y)$ le long
  du cercle $\gamma(t) = (\cos t,\sin t)$, $t \in [0,2\pi]$, de deux
  façons : directement et par le théorème de Green.
\end{exercice}

\begin{exercice}[$\star\star$]
  Soit $\gamma(t) = (e^t\cos t,\; e^t\sin t)$, $t \in [0,2\pi]$.
  Calculer $L(\gamma)$.
\end{exercice}

\begin{exercice}[$\star\star$]
  Calculer $\displaystyle\int_\gamma y^2\,\dd x + x\,\dd y$ où
  $\gamma$ est la parabole $y = x^2$ de $(0,0)$ à $(1,1)$.
\end{exercice}

\begin{exercice}[$\star\star$]
  En utilisant la formule de l'aire de Green, calculer l'aire du domaine
  délimité par l'astroïde\index{astroïde}
  $\gamma(t) = (\cos^3 t,\; \sin^3 t)$, $t \in [0,2\pi]$.
\end{exercice}

\begin{exercice}[$\star\star\star$]
  Soit $\Omega = \R^2\setminus\{(0,0)\}$ et
  $\omega = \frac{-y\,\dd x + x\,\dd y}{x^2+y^2}$. Calculer
  $\int_\gamma\omega$ pour $\gamma(t) = (2\cos t + 1,\; 2\sin t)$,
  $t \in [0,2\pi]$. Le résultat change-t-il si l'on remplace le centre
  du cercle par $(3,0)$ ?
\end{exercice}

\begin{exercice}[$\star\star\star$]
  Soit $D$ le domaine délimité par $y = 0$, $x = 1$ et $y = x$.
  Calculer $\oint_{\partial D}(e^x\sin y)\,\dd x + (e^x\cos y)\,\dd y$
  de deux façons : directement et par le théorème de Green.
\end{exercice}

\begin{exercice}[$\star\star\star$]
  Montrer que tout champ de vecteurs $\mathbf{F}$ de classe
  $\mathcal{C}^1$ sur un ouvert étoilé $\Omega \subset \R^2$ vérifiant
  $\frac{\partial F_1}{\partial y} = \frac{\partial F_2}{\partial x}$
  est un champ de gradient. \emph{Indication : construire explicitement
  le potentiel par intégration le long de segments.}
\end{exercice}

\begin{exercice}[$\star\star\star$]
  Soit $\gamma$ une courbe fermée simple $\mathcal{C}^1$ dans $\R^2$ et
  $\mathbf{n}$ le vecteur normal unitaire extérieur. Pour $\mathbf{F}$
  de classe $\mathcal{C}^1$, montrer la formule de flux-divergence
  en dimension~2 :
  \[
    \oint_\gamma \mathbf{F}\cdot\mathbf{n}\,\dd s
    = \iint_D \mathrm{div}\,\mathbf{F}\,\dd x\,\dd y,
  \]
  où $D$ est le domaine intérieur à $\gamma$ et
  $\mathrm{div}\,\mathbf{F} = \frac{\partial F_1}{\partial x}
  + \frac{\partial F_2}{\partial y}$.
  \emph{Indication : déduire cette formule du théorème de Green.}
\end{exercice}


\bigskip
\begin{center}
\fbox{\parbox{0.85\textwidth}{%
\textbf{Résumé du chapitre~\thechapter.}
\begin{itemize}[nosep]
  \item Un arc paramétré $\gamma\colon[a,b]\to\R^n$ est régulier si
    $\gamma'(t) \neq 0$ partout.
  \item Longueur d'arc : $L(\gamma) = \int_a^b\norm{\gamma'(t)}\,\dd t$.
  \item L'abscisse curviligne paramétrise la courbe à vitesse~$1$.
  \item Intégrale curviligne scalaire :
    $\int_\gamma f\,\dd s = \int_a^b f(\gamma(t))\norm{\gamma'(t)}\,\dd t$.
  \item Intégrale curviligne de travail :
    $\int_\gamma\mathbf{F}\cdot\dd\gamma = \int_a^b
    \mathbf{F}(\gamma(t))\cdot\gamma'(t)\,\dd t$.
  \item Si $\mathbf{F} = \nabla V$, alors
    $\int_\gamma\mathbf{F}\cdot\dd\gamma = V(\gamma(b))-V(\gamma(a))$.
  \item Théorème de Green :
    $\oint_{\partial D}(P\,\dd x + Q\,\dd y)
    = \iint_D(\frac{\partial Q}{\partial x}
    - \frac{\partial P}{\partial y})\,\dd x\,\dd y$.
  \item La condition $\frac{\partial F_1}{\partial y}
    = \frac{\partial F_2}{\partial x}$ est nécessaire mais pas
    suffisante pour l'existence d'un potentiel (topologie !).
\end{itemize}}}
\end{center}


% ============================================================
%  Remarques conclusives
% ============================================================

\section*{Remarques conclusives : perspectives}
\addcontentsline{toc}{section}{Remarques conclusives}

Ce cours a parcouru un chemin allant de l'intégrale de Riemann en une
variable jusqu'au théorème de Green, en passant par les séries de
fonctions, le calcul différentiel en plusieurs variables et les
intégrales multiples. Voici quelques pistes de prolongement naturel.

\medskip

\noindent\textbf{Vers la théorie de la mesure.}\index{théorie de la mesure}
L'intégrale de Riemann, si intuitive soit-elle, connaît des limitations
sérieuses (fonctions intégrables trop restrictives, passages à la limite
insuffisants). La \emph{théorie de la mesure} de Lebesgue fournit un
cadre beaucoup plus souple, où les théorèmes de convergence dominée et
monotone remplacent avantageusement les conditions de convergence
uniforme. C'est le fondement de l'analyse moderne et des probabilités.

\medskip

\noindent\textbf{Vers l'analyse complexe.}\index{analyse complexe}
Les intégrales curvilignes prennent une dimension nouvelle lorsque l'on
remplace $\R^2$ par $\C$. La théorie des fonctions holomorphes
(Cauchy, résidus, séries de Laurent) révèle des liens profonds entre
l'analyse, l'algèbre et la géométrie. Le théorème de Green y apparaît
comme un cas particulier de la formule intégrale de Cauchy.

\medskip

\noindent\textbf{Vers la géométrie différentielle.}\index{géométrie différentielle}
Les formes différentielles que nous avons rencontrées (du type
$P\,\dd x + Q\,\dd y$) se généralisent en \emph{formes
différentielles de degré $k$} sur des variétés, et le théorème de
Green se généralise en le \emph{théorème de Stokes}\index{Stokes (théorème de)}
\[
  \int_{\partial M}\omega = \int_M \dd\omega,
\]
qui unifie les théorèmes de Green, de la divergence et de Stokes classique.

\bigskip

\hfill\emph{Bonne continuation.}


% ============================================================
%  Bibliographie
% ============================================================

\begin{thebibliography}{99}

\bibitem{rudin}
  W.~Rudin,
  \emph{Principles of Mathematical Analysis},
  3\textsuperscript{e} édition,
  McGraw-Hill, 1976.
  \index{Rudin}

\bibitem{zuily-queffelec}
  C.~Zuily et H.~Queffélec,
  \emph{Analyse pour l'agrégation},
  Dunod, 2013.
  \index{Zuily}\index{Queffélec}

\bibitem{ramis-deschamps-odoux}
  E.~Ramis, C.~Deschamps et J.~Odoux,
  \emph{Cours de mathématiques spéciales},
  tomes~1--4,
  Masson, 1991.
  \index{Ramis}\index{Deschamps}\index{Odoux}

\bibitem{tao}
  T.~Tao,
  \emph{Analysis~II},
  3\textsuperscript{e} édition,
  Hindustan Book Agency / Springer, 2016.
  \index{Tao}

\bibitem{pompeiani-tarrago}
  I.~Pompeïani et A.~Tarrago,
  \emph{Analyse},
  Ellipses, 2008.
  \index{Pompeïani}\index{Tarrago}

\bibitem{spivak}
  M.~Spivak,
  \emph{Calculus on Manifolds},
  W.~A.~Benjamin, 1965.
  \index{Spivak}

\end{thebibliography}


\printindex
\end{document}
